\documentclass[11pt]{amsart}
\usepackage{amsmath,amssymb,amsthm}
\usepackage{booktabs,xcolor}
\usepackage[colorlinks=true,linkcolor=blue!45!black,citecolor=blue!45!black,urlcolor=blue!45!black]{hyperref}
\hypersetup{pdftitle={Supporting estimates for small powers in the Waring--Siegel problem},pdfauthor={Morley Davidson},bookmarksnumbered=true}
\setlength{\textwidth}{6.55in}
\setlength{\oddsidemargin}{-0.025in}
\setlength{\evensidemargin}{-0.025in}
\setlength{\textheight}{8.7in}
\newtheorem{theorem}{Theorem}[section]
\newtheorem{lemma}[theorem]{Lemma}
\newtheorem{proposition}[theorem]{Proposition}
\newtheorem{corollary}[theorem]{Corollary}
\theoremstyle{definition}
\newtheorem{definition}[theorem]{Definition}
\newtheorem{remark}[theorem]{Remark}
\numberwithin{equation}{section}
\setlength{\emergencystretch}{2em}
\setcounter{tocdepth}{1}
\title[Supporting estimates for small powers]{Supporting estimates for small powers\\
in the Waring--Siegel problem}
\author{Morley Davidson}
\date{}
\begin{document}
\begin{abstract}
For number fields of arbitrary degree and signature we give the local and analytic estimates needed at degrees four through ten. The estimates combine a localized critical mean value, efficient differencing with principal ideal factors, and auxiliary cubic and quadratic sums. An exact parameter calculation supplies the resulting bounds for smooth moments.
\end{abstract}
\maketitle
\raggedbottom
\tableofcontents
\section{Conventions and results}\label{sec:introduction}
In this supplement write $K=\mathbb K$, $n=[K:\mathbb Q]$ and
$J_{K,k}=\sum_{x\in\mathcal O_K}\mathbb Z x^k$.
The quantity $G_{\mathbb K}(k)$ denotes the least $s$ for which every
totally positive $\nu\in J_{K,k}$ of sufficiently large absolute norm is a
sum of at most $s$ $k$th powers of totally positive integral roots with
$|N x|\ll_{K,k,s}|N\nu|^{1/k}$. A zero is allowed when padding the sum.
At complex places the roots have no sign restriction.
\begin{theorem}\label{thm:main-transfer}
The smooth-moment estimates give the upper bounds for $G_{\mathbb K}(k)$
in the first row below, for every number field $\mathbb K$.
\end{theorem}
\begin{center}
\small
\begin{tabular}{lrrrrrrr}
\toprule
Degree $k$ & 4 & 5 & 6 & 7 & 8 & 9 & 10 \\
\midrule
Smooth moments (this supplement) & 16 & 17 & 24 & 31 & 38 & 46 & 55 \\
\href{https://raw.githubusercontent.com/morley988/warings-problem/main/manuscript/waring-siegel_preprint.pdf\#nameddest=theorem.8.1}{Main article (Theorem~8.1)} & 16 & 17 & 24 & 31 & 38 & 44 & 51 \\
\bottomrule
\end{tabular}
\end{center}
The fibre estimates in the main article sharpen the last two entries.
Both arguments use the smooth moments and the bounds for primitive
local solubility proved here.
\begin{theorem}\label{thm:quartic-refinements}
One has $G_{\mathbb K}(4)\le\max\{12,H_{\mathbb K}(4)\}$, where
$H_{\mathbb K}(4)$ is the maximal primitive width over the finite
completions. The bound is fifteen if the field has no completion
isomorphic to $\mathbb Q_2$ and twelve if each dyadic completion has
residue degree at least two.
\end{theorem}

\section{Preliminaries}
\label{sec:preliminaries}
Fix a number field $K$ of degree $n$, with ring of integers
$\mathcal O=\mathcal O_K$ and different $\mathfrak D$.
Put $V=K\otimes_{\mathbb Q}\mathbb R$ and
$\mathbb T_K=V/\mathfrak D^{-1}$, with probability Haar measure.
We use the trace character $e_K(z)=\exp(2\pi i\operatorname{Tr}z)$;
the notation $E(z)$ has the same meaning. At an archimedean place
let $d_v=1$ or $2$, according as the place is real or complex.
Normalize Lebesgue measure so that $\mathcal O$ has covolume one,
and use the trace-dual measure in coefficient space.

A root box $\mathcal B(Y)$ has fixed shape and archimedean radii
comparable to $Y^{1/n}$; in particular it contains $O_K(Y)$
algebraic integers. Thus $X=P^n$ is the volume parameter for
ordinary roots of archimedean size $P$. Fixed changes in shape,
translation, or scale will be allowed. Moment exponents are measured
in the root volume $X$; auxiliary box parameters also denote volumes.

For a fixed smooth compactly supported weight define
\[
 f_k(\alpha)=\sum_{x\in\mathcal O}w(x/P)e_K(\alpha x^k).
\]
For representation integrals the weights are nonnegative product
weights supported away from zero at every embedding and positive
at the real embeddings. Let $\mathcal A(Y,R)$ consist of the nonzero
roots in $\mathcal B(Y)$ whose prime ideal divisors have norm at most
$R$, and write $g_k(\alpha;Y,R)$ for its exponential sum.
Fixed support restrictions and bounded nonnegative weights are
permitted. The zero root can be appended without changing any
permissible exponent used below.

\begin{definition}
An exponent $\lambda_s$ is \emph{permissible} for degree $k$ if,
for every $\epsilon>0$, some $\eta>0$ gives
\[
 \int_{\mathbb T_K}|g_k(\alpha;Y,R)|^{2s}\,d\alpha
 \ll_{K,k,s,\epsilon}Y^{\lambda_s+\epsilon}
 \qquad(1\le R\le Y^\eta).
\]
The excess at full order $r=2s$ is
$\Delta_r=\lambda_s-(r-k)$. We also write
$U_r=\int|g_k|^r$ when the degree and scale are fixed.
\end{definition}

Linear interpolation between adjacent moment orders always means
H\"older interpolation. A finite collection of permissible estimates
can be used with one sufficiently small positive smoothness exponent.
The constants and this exponent may depend on $K$.

\subsection{Conductor ideals and heights}
For $\gamma\in K/\mathfrak D^{-1}$ put
\[
 \mathfrak a_\gamma=\{a\in\mathcal O:a\gamma\in\mathfrak D^{-1}\},
 \qquad Q_\gamma=N\mathfrak a_\gamma.
\]
For $T\ge1$ define
\[
 A_T(\beta)=\prod_{v\mid\infty}\max(1,T^{1/n}|\beta_v|)^{d_v},
 \qquad
 \mathcal H_T(\alpha)=\min_{\gamma\in K}
              Q_\gamma A_T(\alpha-\gamma).
\]
The minimum includes all lifts of the rational centers. A bound
for the displayed expression leaves finitely many conductor classes
and bounds each archimedean error, so the minimum is attained.
For a degree-$k$ ordinary sum set $B_P=A_{X^k}$, thus
$B_P(\beta)=\prod_v\max(1,P^k|\beta_v|)^{d_v}$.
The broad arcs $\mathfrak M_*$ and narrow arcs $\mathfrak N_c$
are respectively the unions of
\[
 Q_\gamma B_P(\alpha-\gamma)\le C_KX,
 \qquad Q_\gamma B_P(\alpha-\gamma)\le X^c,
\]
where $C_K$ is a sufficiently large fixed constant and $0<c<1$.
Write $\mathfrak m_*=\mathbb T_K\setminus\mathfrak M_*$.

For fixed $0\ne c\in\mathcal O_K$ one has
$\mathcal H_T(c\alpha)\asymp_{K,c}\mathcal H_T(\alpha)$.
Indeed, multiplying or dividing a rational center by $c$ changes
its conductor norm by a factor bounded in terms of $(c)$, and
the archimedean factors change by fixed factors. Taking minima
in the two directions proves the assertion.

\begin{lemma}[Product-height Dirichlet approximation and separation]
\label{lem:g789-height}
For $T\ge1$ one has $\mathcal H_T(\alpha)\ll_K T^{1/2}$.
Moreover there is $c_K>0$ such that if
\[
 \alpha=\gamma_i+\beta_i,\qquad
 (Q_{\gamma_1}A_T(\beta_1))
 (Q_{\gamma_2}A_T(\beta_2))<c_KT,
\]
then the two lifted centers are equal: $\gamma_1=\gamma_2$.
\end{lemma}
\begin{proof}
Apply Minkowski's theorem in $V^2$ to the lattice
$\{(b,a-\alpha b):b\in\mathcal O_K,\ a\in\mathfrak D^{-1}\}$,
whose covolume is fixed independently of $\alpha$, using the bounds
\[
 |b_v|\le C_KT^{1/(2n)},\qquad
 |a_v-\alpha_vb_v|\le C_KT^{-1/(2n)}.
\]
For sufficiently large $T$ the resulting nonzero pair has $b\ne0$,
because a nonzero element of $\mathfrak D^{-1}$ has absolute norm
at least $(N\mathfrak D)^{-1}$. Set $\gamma=a/b$.
Then $Q_\gamma\le|Nb|$, and
\[
 A_T(\alpha-\gamma)
 \le \prod_v(C_KT^{1/(2n)}/|b_v|)^{d_v}
 \ll_K T^{1/2}/|Nb|.
\]
Bounded $T$ is harmless. For separation, a nonzero difference of the
two centers has norm at least
$(Q_{\gamma_1}Q_{\gamma_2}N\mathfrak D)^{-1}$, whereas
\[
 |N(\beta_1-\beta_2)|\le2^nT^{-1}A_T(\beta_1)A_T(\beta_2).
\]
Choosing $c_K$ sufficiently small proves the assertion.
\end{proof}

\subsection{Divisors and quadratic norm equations}
\begin{lemma}\label{lem:divisor-norm-counts}

For fixed $C$ and every $\epsilon>0$, the following assertions hold uniformly
when all displayed nonzero algebraic integers have archimedean size at most $X^C$.

\begin{enumerate}
\item The number of divisors $d\in\mathcal O_K$ of a fixed nonzero $a$, with
all $|d_v|\le X^C$, is $O(X^\epsilon)$.  More generally the number of ideals
of norm at most $X^C$ supported on the prime ideals dividing $(a)$ is
$O(X^\epsilon)$.
\item If $abc\ne0$, the number of solutions of
\begin{equation}\label{eq:g6-binary}
 ax^2+by^2=c,
 \qquad x,y\in\mathcal O_K,\quad |x_v|,|y_v|\le X^C,
\end{equation}
is $O(X^\epsilon)$, uniformly in $a,b,c$.
\end{enumerate}

\end{lemma}
\begin{proof}
For the first assertion, count ideal divisors and then generators.  Ratios of
generators are units of logarithmic height $O(\log X)$; Dirichlet's unit lattice
therefore contributes only a power of $\log X$.  For the stronger supported-ideal
assertion, split the supporting primes at a sufficiently large fixed norm $T$.
The finitely many primes below $T$ contribute a power of $\log X$.
For $\sigma>0$ the large-prime Euler product is bounded by
\[
 \prod_{\substack{\mathfrak p\mid(a)\\N\mathfrak p>T}}
 (1-N\mathfrak p^{-\sigma})^{-1}
 \le \exp\!\left(O\!\left(\frac{\log X}{T^\sigma\log T}\right)\right).
\]
Rankin's inequality, first choosing $\sigma$ small and then $T$ large, gives
the assertion.

To prove the second assertion, put $L=K(\sqrt{-ab})$ and
$z=ax+y\sqrt{-ab}$, so that $N_{L/K}z=ac$.
When $[L:K]=2$, the ideal $(z)$ divides $(ac)\mathcal O_L$.
The number of such ideals is $O(X^\epsilon)$ uniformly in $L$:
$[L:\mathbb Q]\le2n$, and a rational prime has at most $2n$ prime ideals above it,
so the ordinary divisor bound applied to the rational norm suffices.
For each ideal, admissible generators differ by units of logarithmic height
$O(\log X)$.  This count is uniform in $L$: non-torsion algebraic units of
degree at most $2n$ have a positive uniform lower bound for their logarithmic
height, by Northcott's theorem and Kronecker's theorem.  Packing in the
logarithmic unit space gives $O((\log X)^{2n})$ possibilities; the number of
torsion units is bounded in terms of $n$.  Recovering $x,y$ from $z$ is injective.
If $-ab=d^2$ in $K$, then $d\in\mathcal O_K$ and
$(ax-dy)(ax+dy)=ac$, so the first assertion applies directly.
The hypothesis $c\ne0$ is essential; zero right sides are handled separately below.

\end{proof}

\section{Mean values and a transfer criterion}
\label{sec:analytic-principles}\label{sec:ordinary-bundle}
\begin{lemma}[Ordinary moments and pruning]\label{lem:ordinary-estimates}
For fixed $k\ge3$, the ordinary sums of Section~\ref{sec:preliminaries} satisfy

\begin{align}
 \int_{\mathfrak m_*}|f|^{k(k+1)}
    &\ll X^{k(k+1)-k-1+\epsilon},\label{eq:ordinary-critical-bundle}\\
 \int_{\mathfrak M_*}|f|^p&\ll X^{p-k}\qquad(p>2k),\label{eq:ordinary-major-bundle}\\
 \int_{\mathfrak M_*\setminus\mathfrak N_c}|f|^p
    &\ll X^{p-k-c(p/k-2)/2}\qquad(p>2k).
       \label{eq:ordinary-pruning-bundle}
\end{align}
\end{lemma}
We prove these estimates in the following four steps.

\subsection{Packing and close pairs}

If $\mathfrak b$ is a nonzero fractional ideal, then in any translated
archimedean box of radii $R_v$,
\begin{equation}\label{eq:ordinary-packing}
 \#(\mathfrak b\cap\text{box})\ll_K
       1+\frac{\prod_vR_v^{d_v}}{N\mathfrak b}.
\end{equation}
Indeed, every nonzero difference $z\in\mathfrak b$ has
$|Nz|\ge N\mathfrak b$.  After dividing the $v$th coordinate by $R_v$,
the maximum-distance separation is at least
$(N\mathfrak b/\prod R_v^{d_v})^{1/n}$.
Packing disjoint balls in a fixed bounded box proves
\eqref{eq:ordinary-packing}, also for highly unequal radii and translated boxes.

The following close-pair observation is particularly useful.  Suppose
$|y_v|\le C P$, $a\in\mathfrak D^{-1}$, and
\[
 |\theta_vy_v+\eta_v-a_v|\le Z_v\qquad(v\mid\infty).
\]
If $X\prod Z_v^{d_v}\le c_K$ for a sufficiently small fixed $c_K$, then either
there is at most one pair $(y,a)$ or all the pairs lie on one affine $K$-line
$a=a_0+\gamma(y-y_0)$, and their number is at most
\begin{equation}\label{eq:ordinary-line}
 O_K\left(1+Q_\gamma^{-1}
       \prod_v\min\{P,Z_v/|\theta_v-\gamma_v|\}^{d_v}\right).
\end{equation}
A quotient by zero means infinity.
To prove this, take differences $(u_i,b_i)$ from one pair.
The determinants $b_1u_2-b_2u_1\in\mathfrak D^{-1}$ have conjugates
$O(PZ_v)$, hence norm less than $N\mathfrak D^{-1}$, so vanish.
Distinct pairs cannot have the same $y$ by the same norm argument.
They therefore have one slope $\gamma=b/u\in K$.
Their $y$-coordinates lie in a coset of $\mathfrak a_\gamma$ and in the
intersection of the original box with boxes of radii
$O(Z_v/|\theta_v-\gamma_v|)$; apply \eqref{eq:ordinary-packing}.
Without the small-volume hypothesis the elementary bound is
\begin{equation}\label{eq:ordinary-line-trivial}
 \#\{(y,a)\}\ll_K X\left(1+\prod_v Z_v^{d_v}\right),
\end{equation}
by counting $a$ separately for each $y$.

\subsection{The shifted kernel and the critical minor-arc moment}

Choose a smooth compactly supported $v_0$, put $H=P^{k-1}$, and define
$D_H(t)=\sum_{h\in\mathcal O}v_0(h/H)e_K(-th)$.
Poisson summation gives, for every fixed large $L$,
\[
 |D_H(t)|\ll H^n\sum_{a\in\mathfrak D^{-1}}
             \prod_v(1+H|t_v-a_v|)^{-L}.
\]
Let $Y$ be a fixed coordinate box of algebraic integers with side comparable
to $P$, so $\#Y\asymp_KX$.  Uniformly in $\eta\in V$ and
$\theta\in\mathfrak m_*$,
\begin{equation}\label{eq:ordinary-shifted-kernel}
 \frac1{\#Y}\sum_{y\in Y}|D_H(k\theta y+\eta)|\ll X^{k-2}.
\end{equation}
Here is the full counting argument.
Decompose the displayed Poisson majorant into boxes
$Z_v=z_v/H$, $z_v\ge1$ dyadic, and write $B=\prod z_v^{d_v}$.
If $B\le c_KX^{k-2}$, \eqref{eq:ordinary-line} applies to the close-pair
count $N_B$.  Apart from a single pair it gives
\[
 \frac{N_B}{B}\ll 1+
 \frac{X}{Q_\gamma
   \prod_v\max\{z_v,P^k|k\theta_v-\gamma_v|\}^{d_v}}
 \ll_K1.
\]
The last inequality follows from $\theta\notin\mathfrak M_*$:
the conductor of $\gamma/k$ is at most $k^nQ_\gamma$, and
$B_P(\theta-\gamma/k)\le B_P(k\theta-\gamma)$.
If $B>c_KX^{k-2}$, \eqref{eq:ordinary-line-trivial} gives
$N_B/B\ll X/B+X^{2-k}\ll1$, since $k\ge3$.
Summing the convergent dyadic weights and using $H^n/X=X^{k-2}$ proves
\eqref{eq:ordinary-shifted-kernel}.

Define the full Vinogradov sum over a fixed enlarged coordinate box by
\[
 F(\boldsymbol\alpha)=\sum_{x\in\mathcal B(CX)}
    e_K(\alpha_1x+\cdots+\alpha_kx^k),\qquad
 J_{u,k}(CX)=\int_{\mathbb T_K^k}|F(\boldsymbol\alpha)|^{2u}
                 d\boldsymbol\alpha.
\]
For any positive integer $u$, the shifted-kernel argument gives
\begin{equation}\label{eq:ordinary-localized-general}
 \int_{\mathfrak m_*}|f(\theta)|^{2u}d\theta
       \ll X^{k(k-1)/2-1}J_{u,k}(CX).
\end{equation}
For completeness, insert coefficients of $x,\ldots,x^{k-2}$ by periodic
Fourier kernels which equal one on every possible difference of sums of these
powers.  Their suprema cost
$X^{1+\cdots+(k-2)}=X^{(k-1)(k-2)/2}$.
After integrating those coefficients, expansion imposes
$\sigma_j=\sum_{i\le u}(x_i^j-y_i^j)=0$ for $1\le j\le k-2$.
A common translation by $z\in Y$ preserves these equations and gives
\[
 \sigma_k(\mathbf x-z,\mathbf y-z)
      =\sigma_k(\mathbf x,\mathbf y)-kz\sigma_{k-1}(\mathbf x,\mathbf y).
\]
Fourier-complete each shifted weight using
$w((x-z)/P)=\int\widehat w(\xi)e_K(\xi(x-z)/P)\,d\xi$.
The Fourier $L^1$ norm is fixed, and all $x$ now lie in one enlarged box.
Introduce $\alpha_{k-1}$ by $D_H$ with $v_0$ equal to one on all possible
$P^{1-k}\sigma_{k-1}$.
Averaging over $z$ gives the kernel
$D_H(\alpha_{k-1}+kz\theta)$ times factors of modulus one.
Apply \eqref{eq:ordinary-shifted-kernel}.
The remaining product contains $2u$ full Vinogradov sums, each with only a
linear coefficient shifted by a Fourier frequency.
Extend the coefficient domain to $\mathbb T_K^k$ and use H\"older with
exponent $2u$.  Translation of $\alpha_1$ preserves Haar measure, so every
factor is $J_{u,k}(CX)$.
The total cost is
$X^{(k-1)(k-2)/2+k-2}=X^{k(k-1)/2-1}$, proving
\eqref{eq:ordinary-localized-general}.

Wooley's number-field mean-value theorem is
\cite[Theorem~15.1]{Wooley2019}, applied to $t,t^2,\ldots,t^k$,
whose Wronskian is nonzero.  In its volume normalization it gives
\[
 J_{k(k+1)/2,k}(CX)\ll X^{k(k+1)/2+\epsilon}.
\]
Substitution in \eqref{eq:ordinary-localized-general} proves
\eqref{eq:ordinary-critical-bundle}.  This completes the localization of the number-field moment theorem to the product arcs.

\subsection{The classical Hua seeds}

The elementary field Hua estimates needed for the smooth-moment iteration are
\begin{equation}\label{eq:ordinary-Hua}
 \int_{\mathbb T_K}|f|^{2^r}\ll X^{2^r-r+\epsilon}
       \qquad(1\le r\le k).
\end{equation}
The following differencing proof also provides the initial smooth moments.
By orthogonality a bounded weighted moment is at most a fixed multiple of
the unweighted moment on a containing box.
The case $r=1$ follows from $x^k=y^k$.
Assume the result for $r=J<k$.
Repeated Cauchy on the additive lattice gives
$|f|^{2^J}\ll X^{2^J-J-1}\sum_{x,\mathbf h}
e_K(\theta\Delta_{h_1}\cdots\Delta_{h_J}x^k)$, with the usual vertex-support
restrictions; the aggregate on the right is nonnegative.
Multiply by $|f|^{2^J}$ and integrate.
The equation so counted is
$\sum_{i\le2^{J-1}}(u_i^k-v_i^k)=\Delta_{h_1}\cdots\Delta_{h_J}x^k$.
If its common value is zero, the induction hypothesis supplies
$O(X^{2^J-J+\epsilon})$ choices of $\mathbf u,\mathbf v$, and there are
$O(X^J)$ choices of $x,\mathbf h$: a zero shift is immediate, while otherwise
the nonconstant polynomial in $x$ has boundedly many roots.
If its common value is nonzero, every $h_i$ divides it, so ideal divisors and
bounded-height units give $O(X^\epsilon)$ choices for all shifts, followed by
boundedly many $x$.  Thus the total count is $O(X^{2^J+\epsilon})$ in both
cases.  Multiplying by the prefactor proves the next Hua moment.

\subsection{Complete sums, Poisson summation, and pruning}

If $\psi$ is a primitive additive character modulo an integral ideal
$\mathfrak a$, $a$ is a unit modulo $\mathfrak a$, and $b$ is arbitrary, then
\begin{equation}\label{eq:ordinary-complete}
 \left|\sum_{x\bmod\mathfrak a}\psi(ax^k+bx)\right|
       \ll_{K,k,\epsilon}(N\mathfrak a)^{1-1/k+\epsilon}.
\end{equation}
We recall a uniform local proof because the linear twists are required by Poisson.
At a good prime of residue cardinality $q$, not dividing $2k(k-1)$,
write the modulus as $\pi^e$ and split into shells $v(x)=t$.
The deep shell $t\ge\lceil e/k\rceil$ has normalized size $O(q^{-e/k})$.
Outside it, put $r=v(b)$.  If $r\ne(k-1)t$, the residual polynomial on units
has unit derivative, so all shells with residual conductor at least two vanish.
Conductor one is bounded by the finite-field Weil theorem
\cite{Weil1948}.
There is at most one critical shell $r=(k-1)t$; its stationary points are
simple and its normalized size is
$O_k(q^{-t-(e-kt)/2})=O_k(q^{-e/k})$.
Thus the constant is uniform at every good prime.
At each of the finitely many bad primes, the local van der Corput theorem
\cite[Proposition~3.3]{Cluckers2011}, applied to the $k$th derivative $k!a$,
gives $O_{K_\mathfrak p,k}(q^{-e/k})$ uniformly in the integral lower coefficients.
Chinese remaindering and the ideal-divisor bound absorb the product of
good-prime constants into $(N\mathfrak a)^\epsilon$, proving
\eqref{eq:ordinary-complete}.

For $w=\prod_vw_v$ supported in archimedean annuli, let
\[
 I_v(t,u)=\int_{K_v}w_v(z)e_{K_v}(tz^k-uz)\,dz.
\]
Uniform stationary phase and integration by parts give, for every large fixed $L$,
\begin{equation}\label{eq:ordinary-stationary}
 |I_v(t,u)|\ll(1+|t|)^{-d_v/2}
                (1+|u|/(1+|t|))^{-L}.
\end{equation}
At a real place use the nonzero second derivative on the annulus.
At a complex place the real Hessian determinant of $\operatorname{Re}(tz^k)$
is a nonzero constant times $-|t|^2|z|^{2k-4}$.
After a finite partition, the stationary-phase theorem
\cite[Section~7.7]{Hormander2003} applies uniformly in the linear coefficient;
outside the gradient range repeated integration by parts gives the second factor.

For $\mathfrak a=\mathfrak a_\gamma$, $Q=N\mathfrak a$, ideal Poisson gives
\[
 f(\gamma+\beta)=\frac XQ
  \sum_{\xi\in\mathfrak a^{-1}\mathfrak D^{-1}}
    C_{\mathfrak a}(\gamma,\xi)
    \prod_v I_v(P^k\beta_v,P\xi_v),
 \quad
 C_{\mathfrak a}(\gamma,\xi)=
    \sum_{r\bmod\mathfrak a}e_K(\gamma r^k+\xi r).
\]
Primitivity and \eqref{eq:ordinary-complete} apply uniformly in every $\xi$.
Equation~\eqref{eq:ordinary-packing}, with radii
$(1+P^k|\beta_v|)/P$, bounds the sum of the decaying Fourier factors by
$O_K(1+QB_P(\beta)/X)$.
Hence, on the broad arcs,
\begin{equation}\label{eq:ordinary-broad-point}
 |f(\gamma+\beta)|\ll
       XQ^{-1/k+\epsilon}B_P(\beta)^{-1/2}.
\end{equation}
On $QB_P(\beta)\le X^c$ the same argument, using the norm lower bound for
each nonzero dual frequency, gives for every fixed $B>0$
\begin{equation}\label{eq:ordinary-core-formula}
 f(\gamma+\beta)=XG(\gamma)V(P^k\beta)+O(X^{-B}),
 \quad
 G(\gamma)=Q^{-1}C_{\mathfrak a}(\gamma,0),\quad
 V(t)=\prod_v I_v(t_v,0).
\end{equation}
Indeed the first nonzero normalized maximum frequency is at least
$c_K(X/(QB_P(\beta)))^{1/n}$; enough integrations by parts make its entire
tail an arbitrary power of $QB_P(\beta)/X$.

There are at most $Q$ center classes of conductor $\mathfrak a$.
Changing variables $t=P^k\beta$ in \eqref{eq:ordinary-broad-point} therefore
reduces the broad $p$th moment to
\[
 X^{p-k}\sum_{\mathfrak a}Q^{1-p/k+p\epsilon}
         \int_V A_1(t)^{-p/2}\,dt.
\]
Both factors converge for $p>2k$ and sufficiently small $\epsilon$, proving
\eqref{eq:ordinary-major-bundle}.
For the part outside $\mathfrak N_c$, every broad-arc parametrization has
$QA_1(t)>X^c$.  Put $a=p/k-2>0$, $\epsilon=a/(4p)$, and $\tau=a/2$.
Use
$\mathbf1_{\{QA_1(t)>X^c\}}\le X^{-c\tau}Q^\tau A_1(t)^\tau$.
The ideal exponent becomes $-1-a/4<-1$, and
$p/2-\tau=1+p(k-1)/(2k)>1$ makes every real and complex integral converge.
This proves \eqref{eq:ordinary-pruning-bundle}.
These upper bounds allow overlapping broad arcs.

\subsection{A common representation criterion}
\label{sec:product-endgame}

We first remove the small loss from a sufficiently high smooth moment.
This is the absorption argument underlying the mixed ordinary and
smooth method of \cite{Wooley2016}.

Throughout the representation argument, $g$ denotes the unweighted
smooth sum restricted to a fixed compact root annulus away from zero
at every embedding, positive at the real places. Permissible even
moments for a containing box also bound this subsum by orthogonality.
Choose the ordinary weight nonnegative and at least one on its support.

\begin{lemma}[Absorption of a critical moment]
\label{prop:bootstrap-here}
Put $w=k(k+1)$. Suppose that $u<S<w$ are even integers, $S>2k$, and
\[
 U_u=\int_{\mathbb T_K}|g|^u\ll X^{u-k+\Delta+\epsilon},
 \qquad D=(S-u)-(w-S)\Delta>0.
\]
Then $U_S\ll X^{S-k}$.
\end{lemma}
\begin{proof}
Choose a nonnegative ordinary weight which is at least one on the
support of $g$. Orthogonality gives
$U_S\le\int |f|^{S-u}|g|^u$.
On $\mathfrak m_*$ apply H\"older with weights
$(S-u)/w$, $(w-S)/w$, and $u/w$ to the moments of orders
$w$, $u$, and $S$, respectively. On $\mathfrak M_*$ use weights
$(S-u)/S$ and $u/S$, and the ordinary major-arc $S$th moment.
Writing $T=X^{-(S-k)}U_S$, Lemma~\ref{lem:ordinary-estimates} gives
\[
 T\ll X^{-D/w+\epsilon}T^{u/w}+T^{u/S}.
\]
Choose $\epsilon<D/(2w)$. Both powers of $T$ are smaller than one,
and the assertion follows.
\end{proof}

The following form permits either a single smooth moment or an
interpolation of several moments. It also separates the analytic
number of variables from a larger local requirement.

\begin{theorem}[Transfer from smooth mean values]
\label{thm:general-transfer}
Put $w=k(k+1)$. Let $a,b$ be positive integers and put $N_0=a+b>2k$.
Suppose that an even integer $S>b$ satisfies $S>2k$ and
\[
 U_S\ll X^{S-k}.
\]
Let $r_1,\ldots,r_t$ be even positive integers with permissible
excesses $\Delta_1,\ldots,\Delta_t$. Suppose that nonnegative
weights $\theta_1,\ldots,\theta_t$ satisfy
\begin{gather}
 \frac a{w}+\sum_i\theta_i=1,\qquad
 \sum_i r_i\theta_i=b,\qquad
 \delta=\frac a{w}-\sum_i\theta_i\Delta_i>0,
 \label{eq:general-minor-criterion}\\
 p=\frac{aS}{S-b}>2k.\label{eq:general-annulus-criterion}
\end{gather}
If $N\ge N_0$ and every element of every local power order has a
primitive representation by at most $N$ $k$th powers, then
$G_{\mathbb K}(k)\le N$.
\end{theorem}

We prove the theorem by replacing the smooth roots with finite
products of smaller ordinary roots. The same construction will also
be used with a joint ordinary--smooth moment for fifth powers.

\subsubsection{Finite products and the major arcs}
Fix an integer $\ell>1/\eta$ and nonnegative smooth weights $w_1,\ldots,w_\ell$
in small annuli of $V$, positive at all real places.  Define
\[
 h(\alpha)=\sum_{y_1,\ldots,y_\ell\in\mathcal O_K}
 \prod_{j=1}^{\ell}w_j(y_j/P^{1/\ell})
 e\!\left(\operatorname {Tr}\bigl(\alpha(y_1\cdots y_\ell)^k\bigr)\right).
\]
Each tuple supplies one $k$th-power root $y_1\cdots y_\ell$.
Every prime ideal dividing this root has norm at most
$C_KX^{1/\ell}<X^\eta$ for sufficiently large $X$.
The support can be placed inside the annulus used to define $g$.

For any prescribed $\epsilon>0$, the coefficient $b(x)$ of a nonzero root $x$
in $h$ satisfies $0\le b(x)\ll X^\epsilon$.
Indeed, the ideals $(y_1),\ldots,(y_\ell)$ factor $(x)$, so the number of ideal
tuples is bounded by its $\ell$-fold ideal divisor function, which is $O(X^\epsilon)$.
For each ideal tuple the ratios of two admissible generators are units with
all archimedean absolute values bounded above and below by fixed constants;
there are only $O_{K,\ell}(1)$ choices.  Consequently, for every fixed even $r$,
\begin{equation}\label{eq:product-moment-domination}
 \int|h|^r\ll X^\epsilon U_r.
\end{equation}
The same argument applies to joint even moments with ordinary factors.
The symbol $\epsilon$ can be made arbitrarily small by applying the divisor bound
with a correspondingly smaller exponent.  Real moments follow by interpolation.
The order of choices matters.  First choose the error exponent in the given
smooth-moment theorem small relative to the strict minor-arc saving, then fix a
permissible $\eta$ and prove the loss-free bootstrap moment.  Next choose
$\ell>1/\eta$ and $c<1/\ell$.  Finally choose the exponent in the independent
ideal-divisor bound small relative to $c$.  Thus the annular argument does not
require the potentially circular conditions $\eta\ll\epsilon\ll\eta$.

There is a direct major-arc formula for $h$.  Put
\begin{align*}
 H_\ell(\gamma)&=Q_\gamma^{-\ell}
 \sum_{\mathbf y\bmod\mathfrak a_\gamma}
 e\!\left(\operatorname {Tr}(\gamma(y_1\cdots y_\ell)^k)\right),\\
 V_\ell(t)&=\int_{V^\ell}\prod_j w_j(z_j)
 e\!\left(\operatorname {Tr}(t(z_1\cdots z_\ell)^k)\right)d\mathbf z.
\end{align*}
For any $0<c<1/\ell$ and any fixed $B>0$, Poisson summation in the $\ell$ factors gives
\begin{equation}\label{eq:product-core}
 h(\gamma+\beta)=XH_\ell(\gamma)V_\ell(P^k\beta)+O(X^{-B})
 \qquad(Q_\gamma B_P(\beta)\le X^c).
\end{equation}
Here $B$ is arbitrary because all factor weights are smooth.  To see the uniformity,
split each factor into residue classes modulo $\mathfrak a_\gamma$.
In the resulting Poisson formula, each nonzero dual coordinate belongs to
$\mathfrak a_\gamma^{-1}\mathfrak D^{-1}$ and has norm at least $c_K/Q_\gamma$.
The normalized dual radii are $(1+P^k|\beta_v|)/P^{1/\ell}$.
Their product has volume $B_P(\beta)/X^{1/\ell}$, so every nonzero mode has
normalized maximum radius at least
$c_K(X^{1/\ell}/(Q_\gamma B_P(\beta)))^{1/n}$.
Integration by parts once a dual coordinate exceeds the corresponding phase
gradient, followed by ideal-lattice packing, gives an arbitrary power of
$Q_\gamma B_P(\beta)/X^{1/\ell}$.
The trivial bound for the complete mixed sum is sufficient for these nonzero modes.

The pure complete factors satisfy
\begin{equation}\label{eq:product-complete}
 |H_\ell(\gamma)|\ll_{K,k,\ell,\epsilon}Q_\gamma^{-1/k+\epsilon}.
\end{equation}
For a prime-power conductor $\mathfrak p^E$, condition on
$T=\sum_{j<\ell}v_\mathfrak p(y_j)$.
The probability of a given $T$ is at most
$C_\ell(T+1)^{\ell-2}q^{-T}$, and the last-factor normalized $k$th-power sum
is at most $C_{K,k}q^{-\max(E-kT,0)/k}$.
Summing gives $C_{K,k,\ell}(E+1)^{\ell-1}q^{-E/k}$.
The finitely many primes dividing $k$ or the different are absorbed in the constant;
Chinese remaindering and the ideal divisor bound prove \eqref{eq:product-complete}.
The archimedean factors are rapidly decreasing: on the supports, differentiation
in any one factor has nonvanishing gradient of size comparable to $|t_v|$.

\paragraph{The representation main term and its normalization.}
For $a+b=N$, put $L=a+b\ell$ and let
\[
 \mathcal F(\mathbf x,\mathbf y)
   =\sum_{i=1}^a x_i^k+
      \sum_{j=1}^b(y_{j,1}\cdots y_{j,\ell})^k.
\]
Orthogonality makes
$R(\nu)=\int_{\mathbb T_K}f(\alpha)^ah(\alpha)^b e_K(-\alpha\nu)\,d\alpha$
the nonnegative weighted count of these representations.
The core arcs are disjoint for sufficiently large $X$.
Indeed, distinct centers of conductor norms $Q,Q'$ satisfy
$|N(\gamma-\gamma')|\gg_K(QQ')^{-1}$, whereas an overlap forces
$|N(\gamma-\gamma')|\ll_K X^{2c-k}(QQ')^{-1}$.
Here $2c<k$, since $c<1/\ell<1$ and $k\ge4$.

With the minor-arc and annular savings of Theorem~\ref{thm:general-transfer}, the two core Poisson formulae give
\[
 R(\nu)=X^{N-k}\mathfrak S(\nu)\mathfrak J(\nu/P^k)
            +o(X^{N-k}),
\]
uniformly for normalized targets in the compact set used below, where
\[
 \mathfrak S(\nu)=\sum_{\gamma\in K/\mathfrak D^{-1}}
      G(\gamma)^aH_\ell(\gamma)^b e_K(-\gamma\nu),
 \qquad
 \mathfrak J(t)=\int_V V_f(\beta)^aV_\ell(\beta)^b e_K(-\beta t)\,d\beta.
\]
Here $V_f(t)=\int_V w(z)e_K(tz^k)\,dz$ is the ordinary archimedean
factor, and all measures have the normalization fixed in
Section~\ref{sec:ordinary-bundle}.  The conductor majorant is summable
because $N>2k$; the archimedean factors decay rapidly.  These same majorants justify extending the truncated core
series and integral, with a uniform error tending to zero.
Chinese remaindering and finite-ring character orthogonality give
\[
 \mathfrak S(\nu)=\prod_{\mathfrak p}\sigma_{\mathfrak p}(\nu),\qquad
 \sigma_{\mathfrak p}(\nu)=\lim_{E\to\infty}
 q^{-E(L-1)}\#\{(\mathbf x,\mathbf y)\bmod\mathfrak p^E:
                       \mathcal F(\mathbf x,\mathbf y)\equiv\nu\},
 \quad q=N\mathfrak p.
\]
For completeness, a primitive local representation gives a target-uniform
positive lower bound.  Replace each product root $x$ by $(x,1,\ldots,1)$.
At a unit coordinate the derivative has valuation $h=v_{\mathfrak p}(k)$.
Fix the other $L-1$ variables modulo $\mathfrak p^{2h+1}$ at that exact solution.
Every further choice of those variables lifts in the distinguished coordinate
by Hensel's lemma.  Thus for every $E\ge2h+1$ the normalized count is at least
$q^{-(2h+1)(L-1)}$, independently of the target.
At all sufficiently large primes the complete-sum estimates give
$\sigma_{\mathfrak p}(\nu)=1+O(q^{1-N/k+\epsilon})$.
Choose $\epsilon<N/k-2$.  This summable good-prime tail, combined with the
positive lower bounds at the finitely many remaining primes, makes
$\mathfrak S(\nu)$ positive uniformly in $\nu$.

Finally take $P=|N\nu|^{1/(kn)}$ and balance $\nu$ using $k$th powers of
totally positive units, for example squares of units.  Their logarithmic
lattice has full rank, so the normalized targets lie in a compact set with
every embedding bounded away from zero.  Small root boxes near equal $k$th
roots, positive at the real places, give a uniformly positive singular
integral on a finite cover.  Undoing the totally positive unit preserves
the signs of all real roots and preserves their absolute norms.
Every resulting root has norm $O_K(P^n)=O_K(|N\nu|^{1/k})$, exactly as required
by the stated convention for $G_{\mathbb K}(k)$.

\begin{proof}[Completion of the proof of Theorem~\ref{thm:general-transfer}]
For $N=N_0$, the representation integral is $f^ah^b$.
By \eqref{eq:product-moment-domination} and
\eqref{eq:general-minor-criterion}, its minor arcs contribute
\[
 \int_{\mathfrak m_*}|f|^a|h|^b
 \ll X^{N_0-k-\delta+\epsilon}.
\]
On $\mathfrak M_*\setminus\mathfrak N_c$, use the ordinary
$p$th moment and the product $S$th moment, with weights
$1-b/S$ and $b/S$. The saving in the exponent of $X$ is
\[
 \frac c2\left(1-\frac bS\right)\left(\frac pk-2\right)>0.
\]
The independent divisor loss in the product coefficients is chosen
smaller than this saving. The core main term is positive by the
local and archimedean arguments above.

For $N>N_0$, insert $N-N_0$ additional ordinary roots. Their trivial
bound $O(X^{N-N_0})$ preserves both off-core savings, while the core
calculation has $N$ summands and uses the assumed primitive local
representations in that number of variables. This proves the theorem.
\end{proof}

\begin{corollary}[An even-variable criterion]
\label{cor:even-transfer}
Suppose that $u<S<w$ are even, $S>2k$, and
\[
 U_u\ll X^{u-k+\Delta+\epsilon},\qquad
 (S-u)-(w-S)\Delta>0.
\]
Under primitive local solubility in $N\ge S$ variables,
$G_{\mathbb K}(k)\le N$.
\end{corollary}
\begin{proof}
Lemma~\ref{prop:bootstrap-here} supplies $U_S\ll X^{S-k}$.
In Theorem~\ref{thm:general-transfer} take $a=S-u$, $b=u$,
and the two smooth moments of orders $u,S$, with weights
$(w-S)/w$ and $u/w$. The minor-arc saving is $D/w$ and the
ordinary annular order is $p=S$.
\end{proof}

\section{Primitive local representations}
\label{sec:local-solubility}\label{sec:higher-local-bounds}
For the integer ring $A$ of a finite extension of $\mathbb Q_p$,
let $J_k(A)$ be the additive group generated by its $k$th powers.
A representation is primitive if one root is a unit, including
when the target is zero. Let $h_k(A)$ denote the least number of
summands in such representations of every element of $J_k(A)$.

\begin{theorem}\label{thm:local-table}
For every such ring $A$ the following bounds hold:
\[
\begin{array}{c|rrrrrrr}
 k&4&5&6&7&8&9&10\\\hline
 h_k(A)\le&16&10&12&14&35&30&17.
\end{array}
\]
For fourth powers the bound is $5$ at odd primes, $10$ at dyadic
primes of residue degree at least two, and $15$ at every other
dyadic extension except $\mathbb Q_2$.
\end{theorem}

The case of a prime exponent follows from Joly's theorem
\cite[Theorem~7.34]{Joly1970}: the ordinary width is at most
$2k-1$. Apply it to $\nu-1$ and append $1^k$ to obtain the
primitive bounds $h_5\le10$ and $h_7\le14$. We prove the other
entries below.
\subsection{First-order carries and squarefree exponents}
\label{sec:valuation-one-carries}


Let $A$ be the integer ring of a finite extension of $\mathbb Q_p$,
with residue field $\mathbb F_q$, uniformizer $\pi$, and absolute
ramification index $e=v_\pi(p)$. Let $J_k(A)$ denote the additive
group generated by the $k$th powers in $A$. Write $h_k(A)$ for the
least number of positive power summands that represents every member
of $J_k(A)$ with at least one unit root. This requirement includes
the target zero.

\begin{theorem}\label{thm:carry-sixth-tenth}
For the integer ring of every finite extension of every $\mathbb Q_p$,
\[
 h_6(A)\le12,\qquad h_{10}(A)\le17.
\]
At places of residue characteristic five the second bound improves
to $h_{10}(A)\le16$.
\end{theorem}

We prove the assertion through an explicit valuation-one carry
argument. Its finite-field input is the usual $m$-term Waring bound
of Tornheim, also recorded in Joly, Proposition~7.16, and in the
proof of Ramanujam, Lemma~1 \cite{Joly1970,Ramanujam1963}.
The proof below keeps a unit coordinate throughout and treats all
ramification and residue degrees.

\subsubsection{Finite-field compression and primitive lifting}

Suppose $p\nmid m$, and put
\[
 H=(\mathbb F_q^\times)^m,
 \qquad E=\mathbb F_p(H).
\]
The additive span of $H$ is $E$: that span is a finite subring of
$\mathbb F_q$ containing one. Every member of $E$ is a sum of at
most $m$ members of $H\cup\{0\}$.
For completeness, the latter bound also follows from Kneser's
addition theorem. The additive stabilizer of a proper sumset
$r(H\cup\{0\})$ is stable under $H$, hence under $E$, and therefore
is zero. Until the whole field $E$ is reached, its cardinality is
at least $r|H|+1$. Taking $r=(|E|-1)/|H|\le m$ proves the claim.

A unit target in $J_m(A)$ consequently has a primitive representation
by at most $m$ $m$th powers: lift the finite-field representation
and use its unit coordinate in Hensel's lemma. Subtracting one first
gives a primitive representation of every target in $J_m(A)$ using
at most $m+1$ terms. If $m\ge3$ is odd, $m$ terms suffice, since
the zero residue already has the primitive two-term representation
$1+(-1)^m=0$.

We also record the exact lifting precisions needed below. For odd
$p$ and $k=pm$, a primitive power congruence modulo $p^2A$ lifts
without increasing its number of summands. At $p=2$, $k=2m$ with
$m$ odd, the same holds from modulus $8A$. Indeed
\begin{equation}\label{eq:carry-primitive-lifts}
 \frac{(1+pz)^{pm}-1}{p^2}\equiv mz\pmod\pi
 \quad(p>2),\qquad
 \frac{(1+4z)^{2m}-1}{8}\equiv mz\pmod\pi.
\end{equation}
Both normalized polynomials are integral and have a unit derivative.
Multiplying an existing unit root by $1+pz$ or $1+4z$ and applying
Hensel's lemma absorbs any error in the indicated ideal.

\subsubsection{The exact carry and its first-order cancellation}

Write $k=pm$ with $m\ge2$ and $p\nmid m$. For a finite list $t$ put
\[
 C_p(t)=\frac{(\sum_i t_i)^p-\sum_i t_i^p}{p}.
\]
This polynomial has integer coefficients, and
\begin{equation}\label{eq:carry-derivative}
 \frac{\partial C_p}{\partial t_i}
       =(\sum_jt_j)^{p-1}-t_i^{p-1}.
\end{equation}
When $p$ is odd, start with a signed expression
$\nu_0=\sum_i\epsilon_i x_i^{pm}$ and put $t_i=\epsilon_i x_i^m$.
Then $t_i^p=\epsilon_i x_i^{pm}$. In the dyadic case we use an
unsigned expression and put $t_i=x_i^m$; the availability of such
expressions is proved below.

Compress $S=\sum_i t_i\in J_m(A)$ exactly as
$S=\sum_{j=1}^r y_j^m$, with a unit $y_j$, and put $s_j=y_j^m$.
The bound is $r\le m+1$ in general and $r\le m$ when $m\ge3$ is odd.
Equality of the two sums gives the exact identity
\begin{equation}\label{eq:carry-exact}
 \nu_0-\sum_{j=1}^r y_j^k=pD,
 \qquad D=C_p(s)-C_p(t)\in A.
\end{equation}
Every residue involved belongs to $E$, so $\bar D\in E$.
Choose $h_1,\ldots,h_l\in H$, $l\le m$, with
$\sum_i h_i=\bar D$, allowing the empty sum, and set
\begin{equation}\label{eq:carry-constant}
 c=p\sum_{i=1}^l[h_i].
\end{equation}
Here brackets denote the ordinary Teichm\"uller lifts in $A$.
Frobenius permutes $H$, so each $[h_i]$ is a $k$th power of a
Teichm\"uller unit. Repeating it $p$ times expresses $c$ using at
most $pm=k$ positive powers, and $D-c/p\in\pi A$.

If $H\subseteq\mathbb F_p^\times$ and $e\ge2$, this improves to
\begin{equation}\label{eq:carry-second-order}
 D-c/p\in\pi^2A.
\end{equation}
To prove it, identify
$A/\pi^2A\simeq\mathbb F_q[\varepsilon]/(\varepsilon^2)$,
which has characteristic $p$. Denote its $\varepsilon$ coefficient
by a prime. A nonunit $t_i$ or $s_j$ is divisible by $\pi^m$ and
has zero first derivative. A unit has residue in
$\mathbb F_p^\times$, including signed $t_i$ when $p$ is odd, and
therefore its $(p-1)$st power is one. Using
\eqref{eq:carry-derivative} and the exact equality $\sum t_i=\sum s_j$
gives
\[
 D'=\sum_i t_i'-\sum_j s_j'=0.
\]
The constant coefficient of $D$ belongs to $\mathbb F_p$.
Since the lifts $[h_i]$ in \eqref{eq:carry-constant} lie in
$\mathbb Z_p$ and $p\in\pi^2A$, $c/p$ has the same complete class
modulo $\pi^2$. This proves \eqref{eq:carry-second-order}.
Notice that this is a statement about the explicit carry difference;
it does not assert $pA\subseteq J_k(A)$.

\subsubsection{Odd residue characteristic}

If $e=1$, compression and the constant correction leave an error
in $p\pi A=p^2A$. Equation~\eqref{eq:carry-primitive-lifts}
therefore gives the bound
\begin{equation}\label{eq:carry-odd-unramified-count}
 h_k(A)\le k+m+1.
\end{equation}

If $H\nsubseteq\mathbb F_p^\times$, choose Teichm\"uller
$m$th-power values $a,b$ with $a/b$ outside $\mathbb F_p$.
Reserve their two unit roots before compressing the remaining target.
Put $A_0=a^p$, $B_0=b^p$, and $\lambda=-a/b$; then
$A_0+B_0\lambda^p=0$. The polynomial
\[
 F(Z)=\frac{A_0(1+Z)^p+B_0(1+\lambda Z)^p-A_0-B_0}{p}
\]
is integral because the top and constant terms cancel. Its derivative
at zero is $B_0(\lambda-\lambda^p)$, a unit. After compression and
\eqref{eq:carry-constant}, solve $F(z)=D-c/p$ with $z\in\pi A$.
Replace the two underlying $m$th-power values by $a(1+z)$ and
$b(1+\lambda z)$; principal-unit $m$th roots exist because $p\nmid m$.
This absorbs the whole remaining error and gives
\begin{equation}\label{eq:carry-odd-pair-count}
 h_k(A)\le k+m+3.
\end{equation}

Now suppose $H\subseteq\mathbb F_p^\times$ and $e\ge2$.
Reserve four unit roots equal to one. After compression and
\eqref{eq:carry-constant}, the error is $p\pi^2R$ by
\eqref{eq:carry-second-order}. Put $z_4=-z_1-z_2-z_3$ and define
\[
 \mathcal F(z_1,z_2,z_3)
   =\frac{\sum_{i=1}^4(1+\pi z_i)^p-4}{p\pi^2}.
\]
The linear term vanishes. The intermediate binomial coefficients
are divisible by $p$, and $\sum_i z_i^p$ is divisible by $p$ as
an integral polynomial because $\sum_i z_i=0$.
Consequently $\mathcal F$ is integral and reduces to
\[
 -\frac12\sum_{i=1}^4z_i^2
 =-\frac12 z^{\mathsf T}(I_3+\mathbf1\mathbf1^{\mathsf T})z.
\]
The matrix has determinant four. A nondegenerate ternary quadratic
form over a finite field of odd characteristic represents every
value nonsingularly. Indeed, diagonalize and use the intersection
of two square sets to represent every nonzero value in two variables;
for zero, fix the third coordinate to one and represent the negative
of its contribution in the first two. Hensel lifting solves
$\mathcal F(z)=R$. The four values $1+\pi z_i$ remain $m$th powers,
so the four reserved coordinates absorb the error and stay units.
This proves
\begin{equation}\label{eq:carry-odd-four-count}
 h_k(A)\le k+m+5.
\end{equation}

\subsubsection{The two square-compression exceptions}
\label{subsec:carry-two-exceptions}

When $(p,m)=(3,2)$ or $(5,2)$, the last count improves by one by
using an existing compressed unit. The condition
$H\subseteq\mathbb F_p^\times$ forces $q=p$: indeed
$|H|=(q-1)/2\le p-1$, whereas $p^2>2p-1$.
Reserve only three ones and compress the remaining target by at
most three squares. Choose one unit square $b=y_1^2$ among these
compressed values and put $b_1=b$, $b_2=b_3=b_4=1$.
With $z_4=-z_1-z_2-z_3$, use
\begin{equation}\label{eq:carry-existing-unit}
 \mathcal F_b(z)
   =\frac{\sum_{i=1}^4(b_i+\pi z_i)^p-\sum_{i=1}^4b_i^p}
          {p\pi^2}.
\end{equation}
Its linear part is $((b^{p-1}-1)/\pi)z_1$, which is integral.
Modulo $\pi$, its quadratic part is
\[
 \frac{p-1}{2}
   \bigl(\bar b^{,p-2}z_1^2+z_2^2+z_3^2+z_4^2\bigr).
\]
All higher terms vanish modulo $\pi$, by the same top-power
divisibility as above. The underlying three-variable matrix is
\[
 \operatorname{diag}(\bar b^{,p-2},1,1)
       +\mathbf1\mathbf1^{\mathsf T},
 \qquad \det=3\bar b^{,p-2}+1.
\]
At $p=3$, a nonzero square has $\bar b=1$, so this determinant is
nonzero. At $p=5$, $\bar b\in\{1,-1\}$, giving determinants four
and minus two, again nonzero.

Completing the square turns the reduction of
\eqref{eq:carry-existing-unit} into a nondegenerate ternary quadratic
form plus a constant. It therefore represents every target
nonsingularly, including its critical value by a nonzero isotropic
vector. Hensel lifting absorbs the error $p\pi^2R$, and every
$b_i+\pi z_i$ remains a square. The term count is now
\begin{equation}\label{eq:carry-special-count}
 3\text{ compressed terms}+3\text{ reserved terms}
       +k\text{ constant terms}=k+6.
\end{equation}
Thus $h_6(A)\le12$ at $p=3$ and $h_{10}(A)\le16$ at $p=5$.
In the latter case the constant correction still uses at most ten
terms: every residue in $\mathbb F_5$ is a sum of at most two
members of $H=\{1,-1\}$, and both Teichm\"uller values are actual
tenth powers. 

\subsubsection{Dyadic valuation one}

Let $p=2$ and $k=2m$, where $m\ge3$ is odd.
There is a unit $r\in\mathbb Z_2$ with $r^{2m}=-7$.
For instance, Hensel lifting of $z^2+z+2=0$ gives a square root
$1+2z$ of $-7$, and the odd-power equation $r^m=1+2z$ has a
unit solution. Hence
\[
 -1=6\cdot1^k+r^k.
\]
Every signed expression, including one obtained by subtracting
reserved powers, can be made unsigned before applying
\eqref{eq:carry-exact}. Its original length is immaterial.
The exact $m$th-power compression uses at most $m$ terms.

If $H\ne\{1\}$, choose a Teichm\"uller $m$th-power value $a$
whose residue differs from one, and choose $\alpha$ with $\alpha^m=a$.
Reserve $(\alpha,1)$. Replacing these roots by
\[
 \alpha(1+z)^{1/m},\qquad(1-az)^{1/m},\qquad z\in\pi A,
\]
changes their $k$th-power sum by $2F_a(z)$, where
\begin{equation}\label{eq:carry-dyadic-pair}
 F_a(Z)=(a^2-a)Z+a^2Z^2.
\end{equation}
Its derivative at zero is a unit. The correction
\eqref{eq:carry-constant} uses at most $2m$ powers and leaves
$D-c/2\in\pi A$. Hensel lifting with $F_a$ absorbs it exactly.
The number of summands is at most
\begin{equation}\label{eq:carry-dyadic-nontrivial-count}
 m+2+2m=3m+2.
\end{equation}

Suppose $H=\{1\}$ and $e\ge2$. Reserve the pairs $(1,1)$ and
$(1+\pi,1)$; for the second put $a=(1+\pi)^m$.
Equation~\eqref{eq:carry-second-order} gives
$D\equiv n\pmod{\pi^2}$ with $n\in\{0,1\}$, so $c=2n$
uses at most two ones. Since $m$ is odd and $e\ge2$,
$a\equiv1+\pi$ and $a^2\equiv1\pmod{\pi^2}$. Therefore
\[
 G(Z)=\frac{F_a(\pi Z)}{\pi^2}
     =\frac{a^2-a}{\pi}Z+a^2Z^2\equiv Z+Z^2\pmod\pi.
\]
The remaining error is $2\pi^2R$. Choose $t_0$ lifting the square
root of $\bar R$, and solve $G(t)=R-t_0^2$ with $t\in\pi A$.
The first pair, with $z=\pi t_0$ in
\eqref{eq:carry-dyadic-pair} at $a=1$, contributes
$2\pi^2t_0^2$. The second, with $z=\pi t$, contributes
$2\pi^2G(t)$. Their sum is the required error; all four roots
remain units. The count is
\begin{equation}\label{eq:carry-dyadic-trivial-count}
 h_k(A)\le m+4+2=m+6.
\end{equation}

Finally suppose $H=\{1\}$ and $e=1$. The unit $m$th-power
values are exactly $1+2A$: their Teichm\"uller factors are killed,
and the $m$th-power map on principal units is bijective. Thus the
unit $k$th-power classes modulo $8A$ are
\[
 1+4V,\qquad V=\{z+z^2:z\in\mathbb F_q\},
\]
while nonunit $k$th powers vanish at that precision since $k\ge6$.
Their generated additive image is
$(\mathbb Z/8\mathbb Z)\cdot1+4V$. Every class $a+4v$ is a
sum of $r$ unit power classes, where $1\le r\le8$ and
$r\equiv a\pmod8$: use one class $1+4v$ and $r-1$ ones.
The dyadic lifting in \eqref{eq:carry-primitive-lifts} gives
\begin{equation}\label{eq:carry-dyadic-unramified-count}
 h_k(A)\le8.
\end{equation}

\subsubsection{Conclusion for the two exponents}

For $k=6$, the odd bad prime is $3$, where
\eqref{eq:carry-special-count} gives twelve; the other odd cases
give smaller counts. The three dyadic bounds are eleven, nine,
and eight. For $k=10$, the odd bad prime is $5$, where the bound
is sixteen, and the three dyadic bounds are seventeen, eleven,
and eight. At a prime not dividing $k$, finite-field compression
of $\nu-1$, appending $1^k$, and ordinary Hensel lifting give
$h_k(A)\le k+1$. All constructions keep a unit root, so they also
give primitive representations of zero. This proves
Theorem~\ref{thm:carry-sixth-tenth}.
\subsection{Witt coefficients with a positive summand count}
\label{sec:higher-witt-coefficients}

Let $p$ be a prime, let $q$ be a power of $p$, and put
$R=\mathbb F_q[[T]]$ and $B=W(\mathbb F_q)$, the unramified complete
discrete valuation ring with residue field $\mathbb F_q$.
We use $p$-typical Witt vectors $W_{a+1}(R)$, with coordinates
$(z_0,\ldots,z_a)$, Teichm\"uller vectors $[x]=(x,0,\ldots,0)$,
and ghost polynomials
\[
 w_j(z)=\sum_{i=0}^j p^i z_i^{p^{j-i}}.
\]
Their addition and multiplication are the universal integral
polynomials making the ghost maps homomorphisms. In particular
$[x][y]=[xy]$. All constant Witt vectors form the coefficient ring
$W_{a+1}(\mathbb F_q)=B/p^{a+1}B$.

Choose positive integers $N_0,\ldots,N_a$ satisfying
$p^{j-i}N_i\ge N_j$ for $i<j$, and put
\[
 I_N=\{z:z_j\in T^{N_j}R\},\qquad
 M_N=W_{a+1}(R)/I_N.
\]
Weighted homogeneity of the Witt operations shows that $I_N$ is an
ideal: a monomial of Witt weight $p^j$ in its coordinates has
$T$-order at least $N_j$. Multiplication is separately homogeneous
in each input, giving the same assertion when only one input is in
$I_N$. Triangularity in the coordinates gives
$|M_N|=q^{\sum_jN_j}$: choose the coordinates successively modulo
$T^{N_j}$, using the coefficient of $z_j$ equal to one in the
$j$th addition polynomial.

\begin{lemma}\label{lem:higher-witt-coefficients}
Let $G$ be the actual additive span of the vectors $[x]$, $x\in R$.
Every class of $G$ modulo $I_N$ is a positive sum of at most
\begin{equation}\label{eq:higher-witt-cost}
 C_p(N)=L_0+(p-1)\sum_{s=1}^a p^{s-1}L_s,
 \qquad
 L_s=\sum_{j=s}^a\left\lceil\frac{N_j}{p^s}\right\rceil
\end{equation}
actual Teichm\"uller vectors. The bound is independent of $q$.
\end{lemma}
\begin{proof}
Write $\ell(M)=\log_q|M|$ for finite $B$-modules. In characteristic
$p$, multiplication by $p^s$ has the coordinate description
\[
 p^sz=(\underbrace{0,\ldots,0}_{s},
             z_0^{p^s},\ldots,z_{a-s}^{p^s}).
\]
Thus its image modulo $I_N$ has cardinality $q^{L_s}$.
One can count the kernel after any sufficiently large common
$T$-truncation: its coordinate conditions are exactly
$z_i\in T^{\lceil N_{i+s}/p^s\rceil}R$ for $i\le a-s$.
Consequently $\ell(p^sM_N)=L_s$.

The image $M$ of $G$ in $M_N$ is a $B$-submodule, since multiplying
an actual generator by a constant Teichm\"uller vector preserves its
form. Adjoin actual generators $x_i$ successively to obtain a
filtration $0=M_0\subsetneq\cdots\subsetneq M_r=M$ with cyclic
quotients $M_i/M_{i-1}\simeq B/p^{\lambda_i}B$.
A coefficient modulo $p^\lambda$ has a Teichm\"uller digit expansion;
regarding its digit $p^j[c]$ as $p^j$ copies of $[c]$ uses at most
$Q_\lambda=1+p+\cdots+p^{\lambda-1}$ positive scalars. Working
backwards through the filtration therefore represents a target by
at most $\sum_iQ_{\lambda_i}$ actual generators.

For every $s$, the natural surjections
$p^sM_i/p^sM_{i-1}\twoheadrightarrow p^s(M_i/M_{i-1})$ give
\[
 \ell(p^sM)\ge\sum_i\max(\lambda_i-s,0).
\]
The identity
\[
 Q_\lambda=\lambda+
   (p-1)\sum_{s\ge1}p^{s-1}\max(\lambda-s,0)
\]
then bounds the positive cost by
$\ell(M)+(p-1)\sum_{s\ge1}p^{s-1}\ell(p^sM)$.
Each length is at most the corresponding length for $M_N$, since
$p^sM\subseteq p^sM_N$. This proves
\eqref{eq:higher-witt-cost}. In particular, the argument never
replaces actual power generators by arbitrary invariant-factor
coordinates.
\end{proof}

Here is the link to local integral powers. For a local integer ring
$A$ of residue field $\mathbb F_q$, choose its unramified coefficient
ring $B\subseteq A$ and the surjection $D=B[[T]]\to A$ sending
$T$ to $\pi$. On $W_{a+1}(R)$, choosing coordinate lifts to $D$
defines an additive map
\begin{equation}\label{eq:higher-last-ghost}
 \Phi:W_{a+1}(R)\longrightarrow D/p^{a+1}D,
 \qquad \Phi(z)=\sum_{i=0}^a p^i\widetilde z_i^{p^{a-i}}.
\end{equation}
It is independent of the lifts: changing a lift by a multiple of
$p$ changes its $p^{a-i}$th power by a multiple of $p^{a-i+1}$.
Additivity follows by lifting the universal Witt addition and using
the ghost identity. In particular $\Phi([x])=\widetilde x^{p^a}$.
Lift an eligible target by any finite signed power expression to
$D$; its associated signed Teichm\"uller sum reduces to an element
of $G$.

Suppose a centered chart covers $I_N$ in the form
\[
 w=\sum_{i=1}^r[x_i]-\sum_{i=1}^r[b_i],
\]
retaining a unit among the $x_i$. Subtract the fixed baseline
$\sum[b_i]$ from the target and use Lemma~\ref{lem:higher-witt-coefficients}
modulo $I_N$. The remaining ideal element is supplied by the chart,
and the baseline cancels. The result is a primitive representation
by at most $C_p(N)+r$ powers modulo $p^{a+1}D$, and hence modulo
$p^{a+1}A$. Fixed nonzero constant residues in the chart are
preserved by choosing lifts of the roots with those constants.

\subsection{Fourth powers}
\begin{lemma}[Primitive lifting]\label{lem:primitive-lifting}
At an odd prime, a fourth-power congruence modulo $\pi A$ with
a unit root lifts to an exact representation of the same target
using the same number of summands. At a dyadic prime, the same
assertion holds for a congruence modulo $16A$.
\end{lemma}
\begin{proof}
For odd $p$, the derivative $4x^3$ is a unit at a unit root.
For $p=2$, changing a unit root $x$ to $x+4z$ changes its fourth
power by $16F_x(z)$, where
\[
 F_x(z)=x^3z+6x^2z^2+16xz^3+16z^4.
\]
Modulo $\pi$, this is the invertible linear map $z\mapsto\bar x^3z$.
Simple-root Hensel lifting solves $F_x(z)=a$ for every $a\in A$.
Thus any error in $16A$ can be corrected by that root alone.
\end{proof}

\begin{lemma}[Odd residue characteristic]\label{lem:odd-local}
If $p\ne2$, then $h_4(A)\le5$.
\end{lemma}
\begin{proof}
The prime-to-residue-characteristic case of
Ramanujam~\cite[Lemma~1]{Ramanujam1963} says that the reduction of
$J_4(A)$ is the subfield generated by fourth powers and that each
of its elements is a sum of at most four fourth powers in the
residue field. If the target is a unit, such a representation has
a nonzero root and lifts. If it is a nonunit, apply the unit case
to $\nu-1$ and append $1^4$. This also treats $\nu=0$.
\end{proof}

We use only length-three $2$-typical Witt vectors. Their standard
integral polynomial construction is described in
Hazewinkel~\cite[Sections~5.2--5.22]{hazewinkel2009}; the specific
open-image assertion below is proved directly.
\subsubsection{The three-coordinate open-image lemma}

Let $C$ be a commutative ring, put $R=C[[T]]$, and let $W_3(R)$ denote
the $2$-typical Witt vectors of length three.  Its ghost polynomials are
\[
 w_0(X)=X_0,\qquad w_1(X)=X_0^2+2X_1,\qquad
 w_2(X)=X_0^4+2X_1^2+4X_2.
\]
Write $[y]=(y,0,0)$ for the Teichm\"uller vector.  Define
\[
 I=\{z\in W_3(R):z_0\in T^3R,\ z_1,z_2\in T^4R\}.
\]
\begin{lemma}\label{lem:wittopen}
The set $I$ is an ideal, and every $w\in I$ has the form
\begin{equation}\label{eq:quartic-specialized-open-image}
 w=[1+T^4z_0]+[T+T^3z_1]+[T^2+T^3z_2]
       -[1]-[T]-[T^2],\qquad z_0,z_1,z_2\in R.
\end{equation}
In particular $I$ lies in the additive subgroup generated by
Teichm\"uller vectors, and the root in the first summand of
\eqref{eq:quartic-specialized-open-image} is $1$ modulo $T$.
\end{lemma}

\begin{proof}
Give Witt coordinate $j$ weight $2^j$.  The integral Witt addition
and negation polynomials are homogeneous of the corresponding weight;
the multiplication polynomials are bihomogeneous in the two input
vectors.  The inequalities
\[
 2\cdot3>4,\qquad4\cdot3>4,\qquad2\cdot4>4
\]
show that the specified divisibilities are preserved under addition,
negation, and multiplication by an arbitrary Witt vector.  This proves
the ideal assertion.  It may equally be checked directly from the
three displayed ghost polynomials, working universally over $\mathbb Z$
before specializing to $C$.

For the open-image assertion, work first with integral universal Witt
polynomials and put
\[
 H(\eta_0,\eta_1,\eta_2)
   =\sum_{i=0}^2[T^i+\eta_i]-\sum_{i=0}^2[T^i].
\]
All coordinates of $H$ vanish at $\boldsymbol\eta=0$.  Its ghost
coordinates satisfy
\[
 w_j(H)=\sum_{i=0}^2\bigl((T^i+\eta_i)^{2^j}-T^{i2^j}\bigr).
\]
Differentiating at zero and using the vanishing of the lower
coordinates shows that the linear part of $H_j$ is
\[
 \sum_{i=0}^2 T^{i(2^j-1)}\eta_i.
\]
This calculation takes place over $\mathbb Z$ (or first over
$\mathbb Q$ with the resulting integral Witt polynomials); it does
not divide by $2$ inside $C$.

Substitute
\[
 (\eta_0,\eta_1,\eta_2)=(T^4z_0,T^3z_1,T^3z_2)
\]
and divide the three coordinates by $T^3,T^4,T^4$, respectively.
The resulting polynomial map $\Phi:R^3\longrightarrow R^3$ is
integral and has linear part
\[
 U=\begin{pmatrix}
 T&1&1\\
 1&1&T\\
 1&T^2&T^5
 \end{pmatrix},\qquad
 \det U=(T-1)(T^2-1)(T^3-1).
\]
For completeness, integrality and the nonlinear divisibilities follow
directly from the ghost recurrence.  The first coordinate is linear.
The second is
\[
 H_1=\sum_iT^i\eta_i-\sum_{i<j}\eta_i\eta_j,
\]
whose nonlinear terms have $T$-adic order at least six.  In the third
ghost equation, every nonlinear term in
$\sum_i((T^i+\eta_i)^4-T^{4i})$ has order at least eight, while
$H_0^4$ has order at least twelve and the nonlinear contribution of
$2H_1^2$ has order at least eight.  Dividing numerical coefficients by
$4$ does not change these $T$-adic orders, and the universal Witt
coordinate polynomial is integral.  Hence
\[
 \Phi(\mathbf z)=U\mathbf z+(0,T^2P_1(\mathbf z),T^4P_2(\mathbf z))
\]
for integral polynomials $P_1,P_2$.

The constant term of $\det U$ is $-1$, so $U$ is invertible over $R$.
For an arbitrary $w\in I$, divide its coordinates by $T^3,T^4,T^4$
to obtain $\mathbf b\in R^3$.  The equation $\Phi(\mathbf z)=\mathbf b$
has a solution by $T$-adic successive approximation: the map
\[
 \mathbf z\longmapsto U^{-1}\!\left(
 \mathbf b-(0,T^2P_1(\mathbf z),T^4P_2(\mathbf z))\right)
\]
increases the precision of two congruent inputs by at least one
power of $T$.  Completeness of $R$ gives a fixed point.  This proves
\eqref{eq:quartic-specialized-open-image} after specialization to $C$.
\end{proof}

\paragraph{The quartic ghost map.}
Let $B=W(\mathbb F_q)$ be an unramified dyadic integer ring, put
$D=B[[T]]$, and take $C=B/4B$.  The last ghost coordinate induces
an additive map
\[
 \Psi:W_3(D/4D)\longrightarrow D/16D,\qquad
 \Psi(z)=\widetilde z_0^4+2\widetilde z_1^2+4\widetilde z_2.
\]
Changes of lifts by $4D$ change the expression by a multiple of $16D$,
so it is well defined; the integral ghost identities prove additivity.
Moreover $\Psi([y])=\widetilde y^4$, and
\[
 4T^4D\pmod{16D}\subseteq\Psi(I),\qquad
 2T^8\pmod{16D}\in\Psi(I),
\]
by taking vectors $(0,0,T^4a)$ and $(0,T^4,0)$, respectively.
Applying $\Psi$ to \eqref{eq:quartic-specialized-open-image} gives
the three-positive-fourth-power correction with fixed baseline
$1+T^4+T^8$ used in the local endpoint proof.

If $A$ is the integer ring of any finite extension of $\mathbb Q_2$
with residue field $\mathbb F_q$, choose its maximal unramified
coefficient ring $B\subseteq A$ and a uniformizer $\pi$.
Substitution $T\mapsto\pi$ gives a surjective continuous map
$D\to A$: every element of $A$ has a convergent $\pi$-adic expansion
with coefficients in a system of residue representatives in $B$.
Composing the ghost map with this substitution gives
\[
 \Psi_A:W_3((B/4B)[[T]])\longrightarrow A/16A.
\]
Its image on the additive subgroup generated by the vectors $[y]$
is exactly the image of $J_4(A)$ modulo $16A$, because each root in
$A$ has a power-series preimage.  The same construction with
$B=\mathbb Z_2$ will be used for residue degree one.  Constant Witt
vectors evaluate through the corresponding ghost map into $B/16B$.

\subsubsection{The local fourth-power bound}
\label{sec:universal-quartic}

We now prove the primitive local bound needed in the global argument.
The finite additive-group calculation is established explicitly below.

\begin{theorem}\label{thm:universal-quartic-sixteen}
For every finite extension $L/\mathbb Q_p$,
\[
 h_4(\mathcal O_L)\le16.
\]
The constant is best possible in the primitive convention:
$h_4(\mathbb Z_2)=16$.
In dyadic residue degree at least two the bound is $10$; in residue
degree one it is $11$ when the ramification index is two and $15$
when that index is at least three.
\end{theorem}
\begin{proof}
At an odd prime, Lemma~\ref{lem:odd-local} gives the bound $5$.
The dyadic cases of residue degree at least two, ramification index
two and residue degree one, and ramification index at least three
and residue degree one are proved in
Propositions~\ref{prop:quartic-residue-degree}, \ref{prop:quartic-ramification-two}, \ref{prop:quartic-high-ramification} and~\ref{prop:quartic-cubic-ramification} below.
The only remaining local ring is $\mathbb Z_2$. Every class modulo
$16$ is a sum of $r$ copies of $1^4$ for some $1\le r\le16$.
Lemma~\ref{lem:primitive-lifting} gives the upper bound. Conversely,
every unit fourth power is $1$ modulo $16$, while every nonunit
fourth power is zero. In a primitive representation of zero the
positive number of unit summands is therefore divisible by $16$.
This proves equality over $\mathbb Z_2$ and the theorem.
\end{proof}

\subsubsection{Fourth powers when the dyadic residue degree is at least two}
\label{sec:quartic-residue-degree}

\begin{proposition}
\label{prop:quartic-residue-degree}
Let $A$ be the integer ring of a finite extension of $\mathbb Q_2$ with
residue field $\mathbb F_q$, $q=2^f$. If $f\ge2$, then
\[
 h_4(A)\le10.
\]
More precisely, nine summands suffice for eligible targets with nonzero
residue. The estimate is uniform in the ramification index.
\end{proposition}

\begin{proof}
Let $B=W(\mathbb F_q)$ and first consider a constant target
$\lambda\in B/16B$. Write $[c]_B$ for the Teichmuller lift of
$c\in\mathbb F_q$. Every $[c]_B$ is a fourth power, since the fourth-power
map is an automorphism of the group of roots of unity of odd order.

If $\bar\lambda\ne0$, choose distinct nonzero $c_1,c_2$ with
$c_1+c_2=\bar\lambda$. If $\bar\lambda=0$, choose distinct nonzero
$c_1,c_2$ and put $c_3=c_1+c_2$. These choices are possible because
$q\ge4$, and the selected residues sum to $\bar\lambda$.
After subtracting their Teichmuller lifts, represent the remaining
target modulo $8B$ by its next two binary digits, costing at most
$2+4=6$ additional Teichmuller fourth powers. Thus there is a
representation modulo $8B$ using at most eight powers when the residue
is nonzero, or nine otherwise. It includes two unit roots whose fourth
powers have distinct nonzero residues.

The last binary digit can be corrected without adding summands. For a
unit root $u$ and $z\in B$,
\[
 (u(1+2z))^4-u^4\equiv8u^4(z^2+z)\pmod{16B}.
\]
The image $V=\{z^2+z:z\in\mathbb F_q\}$ is the trace-zero hyperplane.
For distinct nonzero $c,d$, the hyperplanes $cV$ and $dV$ are distinct:
they are the kernels of the different linear functionals
$x\mapsto\operatorname{Tr}(x/c)$ and
$x\mapsto\operatorname{Tr}(x/d)$. Hence $cV+dV=\mathbb F_q$.
Varying the two chosen roots therefore corrects any error in the last
binary digit. Their residues remain distinct and nonzero. We have
proved the same eight- or nine-term bound modulo $16B$.

Now use the Witt model $W_3((B/4B)[[T]])$ for a target
$\nu\in J_4(A)$. Choose a preimage $z$ in the additive Teichmuller span
and let $z_{\mathrm c}$ be its evaluation at $T=0$. Its top ghost is some constant
$\lambda\in B/16B$. Construct a constant Teichmuller sum $s$ from the
preceding paragraph with this same top ghost. Replacing $z$ by
$z-z_{\mathrm c}+s$ changes no image in $A/16A$, since
$z_{\mathrm c}$ and $s$ have the same top ghost modulo $16B$.
The new vector evaluates exactly to $s$ at $T=0$.

Append one zero root to $s$. Among the roots are two units with distinct
nonzero residues $u,v$, together with the zero root. The derivative
matrix of the centered three-variable Witt sum has rows with exponents
$0,1,3$; its residue determinant is
\[
 \det\begin{pmatrix}1&1&1\\u&v&0\\u^3&v^3&0\end{pmatrix}
       =uv(v^2-u^2)\ne0.
\]
The centered Witt-coordinate map has the displayed invertible
linear part.  On substituting perturbations $Tz_i$ and dividing its
coordinates by $T$, all nonlinear terms acquire an extra factor $T$.
The same successive-approximation argument as in Lemma~\ref{lem:wittopen}
therefore realizes every Witt perturbation whose coordinates lie in
$T(B/4B)[[T]]$ using these three existing roots. It corrects the entire remaining
$T$-dependent part of the target, while preserving both unit roots.
All other coordinates are kept fixed, so no further summands are
introduced. Substitution $T=\pi$ gives a primitive congruence modulo
$16A$, and Lemma~\ref{lem:primitive-lifting} gives the exact representation in $A$.
The appended zero increases the counts to nine and ten, respectively.
\end{proof}


\subsubsection{Fourth powers at absolute ramification index two}
\label{sec:quartic-ramification-two}

\begin{proposition}
\label{prop:quartic-ramification-two}
If $A$ is a dyadic integer ring with residue field $\mathbb F_2$ and
absolute ramification index $2$, then $h_4(A)\le11$.
\end{proposition}
\begin{proof}
Normalize $v(\pi)=1$, so $v(2)=2$, and write $\pi^2=2u$ with
$u\in A^\times$. We first work modulo $\pi^7A$.
A fourth power depends only on its root modulo $\pi^2A$.
Indeed, for $x,h\in A$,
\[
 (x+\pi^2h)^4-x^4
 =4x^3\pi^2h+6x^2\pi^4h^2+4x\pi^6h^3+\pi^8h^4.
\]
The last two terms have valuation at least $8$. If $x\in\pi A$,
the first two have valuation at least $8$ as well. If $x$ is a unit,
the sum of the first two terms, divided by $\pi^6$, has residue
$\bar h+\bar h^2=0$ in $\mathbb F_2$; here
$4/\pi^4=u^{-2}$ and $6/\pi^2=3u^{-1}$ both have residue $1$.
This proves the assertion.

The four root classes are $0,1,\pi,1+\pi$. Put
\[
 t=\pi^4,\qquad w=(1+\pi)^4,\qquad r=w-1.
\]
In the additive group of $A/\pi^7A$ they satisfy
\[
 16=0,\qquad 4t=0,\qquad 2t=8,\qquad 2r=0.
\]
The first two relations follow from their valuations. For the third,
$2t-8=8(u^2-1)$ has valuation at least $8$, since both $u-1$ and
$u+1$ lie in $\pi A$. For the last relation, expand
\[
 r=4\pi+6\pi^2+4\pi^3+\pi^4
   =4\pi+4u(3+u)+4\pi^3.
\]
Every term has valuation at least $5$, so $2r\in\pi^7A$.

Consequently every eligible class has the form
\[
 n+jt+\varepsilon r,
 \qquad 0\le n\le7,\quad0\le j\le3,\quad
 \varepsilon\in\{0,1\}.
\]
If $\varepsilon=1$, replace $r$ by $w-1$, reduce the coefficient of
$1$ modulo $8$, and absorb the resulting multiple of $8$ into the
coefficient of $t$ using $8=2t$. This gives at most seven copies of
$1$, three copies of $t$, and one copy of $w$. It is primitive because
$w$ has the unit root $1+\pi$.
If $\varepsilon=0$ and $n>0$, the displayed form itself is primitive
and uses at most ten powers. If $\varepsilon=n=0$, use
\[
 jt=8+\bigl(j+2\bmod4\bigr)t
\]
modulo $\pi^7A$. This uses at most eleven powers and has unit roots.
Thus every eligible class has a primitive representation in at most
eleven fourth powers.

Finally, a primitive congruence modulo $\pi^7A$ lifts with the same
number of summands. For a unit root $x$, the polynomial
\[
 \frac{(x+\pi^3 Z)^4-x^4}{\pi^7}
\]
has coefficients in $A$, has unit linear coefficient
$4x^3/\pi^4$, and all its higher coefficients lie in $\pi A$.
Hensel's lemma therefore corrects any residual error by changing that
root alone. This proves the proposition.
\end{proof}

\subsubsection{The finite quotient in residue degree one}
\label{sec:quartic-conductor}

The following finite quotient handles all ramification indices at least
three in residue degree one. Its positive width is determined explicitly.

\subsubsection{A universal finite quotient}

Put $D=\mathbb Z_2[[T]]$, $R=D/4D$, and
$W=W_3(R)$.  Let $G$ be the additive subgroup of $W$ generated by all
Teichm\"uller vectors $[y]$, $y\in R$.  The last ghost component gives
an additive map
\[
 \Psi:W\longrightarrow D/16D,
 \qquad
 (z_0,z_1,z_2)\longmapsto
 \widetilde z_0^4+2\widetilde z_1^2+4\widetilde z_2\pmod {16D},
\]
independent of the chosen lifts.  Its image on $G$ is the image of
$J_4(D)$.  Lemma~\ref{lem:wittopen} gives
\[
 I=\{z\in W:z_0\in T^3R,\ z_1,z_2\in T^4R\}\subseteq G.
\]
Write $K=\Psi(I)$.  In particular,
\begin{equation}\label{eq:quartic-conductor-basic}
 4T^4D\pmod {16D}\subseteq K,
 \qquad 2T^8\pmod {16D}\in K.
\end{equation}

Define the additive coefficient projection
\[
 \Pi:D/16D\longrightarrow
 V:=\mathbb Z/8\mathbb Z\oplus(\mathbb Z/16\mathbb Z)^3
       \oplus(\mathbb Z/4\mathbb Z)^4\oplus\mathbb Z/2\mathbb Z
\]
by reading coefficients of $T^0,\ldots,T^8$, in that order, with the
indicated moduli.  Higher coefficients are discarded.  Polynomial
notation below denotes classes in this additive group; no ring
structure on $V$ is asserted.  Set
\[
 Q=\Pi(J_4(D)).
\]

\begin{lemma}\label{lem:quartic-universal-quotient}
The kernel of $\Pi$ on the image of $J_4(D)$ in $D/16D$ is
$K+8\mathbb Z_2$.  Every class of $Q$ is a sum of at most twelve
fourth-power images, with at least one root having odd constant
coefficient.
\end{lemma}

\begin{proof}
First, $K+8\mathbb Z_2$ is killed by $\Pi$.  Indeed, for $z\in I$
its three ghost terms have orders at least $12,8,4$, respectively.
The first is discarded; the second has an even coefficient at $T^8$;
and the third has coefficients divisible by $4$ in degrees $4$ through
$7$ and an even coefficient in degree $8$.  These are all zero in $V$,
as is the constant $8$.  Conversely, given
$\alpha\in J_4(D)$, choose $z\in G$ with $\Psi(z)=\alpha\pmod {16D}$.
By successively reducing its three Witt coordinates modulo $I$, choose
a representative $w$ of the same coset with its coordinate degrees
less than $3,4,4$, respectively.  The successive reduction is performed using the triangular form of Witt addition.  Therefore
\[
 \alpha\equiv\gamma:=\widetilde w_0^4+
                 2\widetilde w_1^2+4\widetilde w_2\pmod K,
 \qquad \deg\gamma\le8.
\]
If $\Pi(\alpha)=0$, the coefficients of $\gamma$ have the indicated
zero residues.  Thus
\[
 \gamma\in8\mathbb Z_2+4T^4D+2\mathbb Z_2T^8
       \pmod {16D},
\]
which lies in $K+8\mathbb Z_2$ by
\eqref{eq:quartic-conductor-basic}.  This proves the kernel assertion.

Every root can be reduced modulo $(4,T^4)$ for this calculation.
Indeed, if $\Delta\in T^4R$, the coordinates of
$[y+\Delta]-[y]$ are
\[
 \Delta,\qquad y\Delta,\qquad
 y^3\Delta+y^2\Delta^2+y\Delta^3,
\]
so this difference lies in $I$.  Changes of a lift by $4D$ do not
change its fourth power modulo $16D$.  There are consequently only
$4^4=256$ root classes to consider.

Here is an explicit description of their additive image.  In $V$, put
\[
 h_1=8T,\qquad h_2=8T^2,\qquad h_3=4T^3+2T^6,
 \qquad H=\langle h_1,h_2,h_3\rangle.
\]
The orders of these independent generators are $2,2,4$, so $|H|=16$.
For odd constant coefficient, changing the sign of the root permits
normalization to
\[
 F(a,b,c):=\Pi\bigl((1+aT+bT^2+cT^3)^4\bigr),
 \qquad a,b,c\in\mathbb Z/4\mathbb Z.
\]
Direct expansion gives
\begin{align*}
 F(a,b,c)={}&1+4aT+(4b+6a^2)T^2
       +(4c+12ab+4a^3)T^3\\
 &+(a^4+2b^2)T^4+2(c^2+a^2b^2)T^6+b^4T^8.
\end{align*}
All coefficients here are read in $V$.  In particular,
\begin{align*}
 F(a,b,c+1)-F(a,b,c)&=h_3,\\
 F(a,b+2,c)-F(a,b,c)&=h_2+2a h_3,\\
 F(a+2,b,c)-F(a,b,c)&=h_1+(a+1)h_2+2(a+b)h_3.
\end{align*}
It follows that $H\subseteq Q$ and that the unit-power image is
stable under translation by $H$.  Its four cosets modulo $H$ are
obtained by taking $a,b\in\{0,1\}$ and $c=0$.

Set, in $Q/H$,
\[
 u=1,\qquad x=4T+6T^2+4T^3+T^4,\qquad
 \alpha=T^4,\quad\beta=T^8,\quad\zeta=2T^6.
\]
These elements are independent, of orders $8,4,4,2,2$, respectively.
For example, a relation is first inspected in its constant and $T$
coefficients; the $T$ coefficient gives $4X\in8\mathbb Z$ and hence $X$ even.
The $T^2$ coefficient then gives $6X\in8\mathbb Z$ modulo $16$,
forcing $X=0$ modulo $4$.  The $T^4,T^8,T^3,T^6$ coefficients
settle the remaining generators.  The unit-power cosets are
\begin{equation}\label{eq:quartic-unit-cosets}
 u,\qquad u+x,\qquad u+2x+\beta,\qquad u+3x+\beta.
\end{equation}
For a root with even constant coefficient, expansion gives exactly
these nonunit cosets:
\begin{equation}\label{eq:quartic-nonunit-cosets}
 0,\qquad\alpha,\qquad\beta,\qquad\alpha+\beta+\zeta.
\end{equation}
For instance they are represented by $0,T,T^2,T+T^2$.
The possible extra contribution from constant coefficient $2$ is
$8a^2T^2+8a^3T^3$, which lies in $H$.
Equations \eqref{eq:quartic-unit-cosets} and
\eqref{eq:quartic-nonunit-cosets} generate all of $Q/H$, and show that
\[
 Q/H\simeq\mathbb Z/8\mathbb Z\oplus
 (\mathbb Z/4\mathbb Z)^2\oplus(\mathbb Z/2\mathbb Z)^2.
\]

To represent a target with coordinates $(C,X,A,B,Z)$ in this product,
choose $1\le r\le8$ with $r\equiv C\pmod8$.  A unit type indexed by
$j\in\{0,1,2,3\}$ contributes
$u+jx+\lfloor j/2\rfloor\beta$.  With three unit slots one can
prescribe both the $x$ and $\beta$ coordinates: for $\beta$ coordinate
zero, take triples $(0,0,0),(1,0,0),(1,1,0),(1,1,1)$ to obtain
$x$ coordinates $0,1,2,3$; for $\beta$ coordinate one, take
$(3,1,0),(3,1,1),(2,0,0),(3,0,0)$, respectively.
Extra unit slots may have index zero.

If $r\ge3$, append $Z\in\{0,1\}$ copies of the last type in
\eqref{eq:quartic-nonunit-cosets}, and then choose
$0\le s\le3$ copies of $\alpha$ so that $s+Z\equiv A\pmod4$.
The unit types can match the required $X$ and $B-Z$ coordinates.
The total is at most $r+s+Z\le12$.
If $r=1$ or $2$, use one unit type with index $X$, and index zero
for the remaining unit slot if present.  One additional $\beta$
corrects the remaining binary coordinate if necessary.  The total
is then at most $r+3+1+1\le7$.
Finally, because the unit-power image is $H$-stable, one unit
summand absorbs any discrepancy in $H$.  This proves the lemma.
\end{proof}

\begin{proposition}\label{prop:quartic-high-ramification}
Let $A$ be the valuation ring of a finite extension of $\mathbb Q_2$
with residue field $\mathbb F_2$ and ramification index $e\ge4$.
Then
\[
 h_4(A)\le15.
\]
\end{proposition}

\begin{proof}
Choose a uniformizer $\pi$.  Substitution $T\mapsto\pi$ defines a
surjective continuous map $D\to A$, hence sends the universal
power-generated group onto the corresponding group for $A$.
The image $K_A$ of $I$ under the last ghost map and this substitution
contains $4\pi^4A$ modulo $16A$.  Since $e\ge4$,
\[
 8=4\pi^4\,(2/\pi^4)\in4\pi^4A\subseteq K_A.
\]
Consequently the quotient in Lemma~\ref{lem:quartic-universal-quotient}
surjects onto $J_4(A)/K_A$, with all quotients here taken modulo $16A$.

Given $\nu\in J_4(A)$, apply that lemma to
\[
 \nu-(1+\pi^4+\pi^8)
\]
and obtain a sum $v$ of at most twelve fourth powers with
$\nu-(1+\pi^4+\pi^8)-v\in K_A$.
Lemma~\ref{lem:wittopen} supplies, for every element
of $I$, an identity
\[
 w=[y_0]+[y_1]+[y_2]-[1]-[T]-[T^2],
 \qquad y_0\equiv1\pmod T.
\]
Applying the last ghost component and substituting $T=\pi$ cancels
the baseline.  Thus
\[
 \nu\equiv v+y_0(\pi)^4+y_1(\pi)^4+y_2(\pi)^4\pmod {16A}
\]
is a representation by at most fifteen fourth powers containing a
unit root.  Lemma~\ref{lem:primitive-lifting} completes the proof.
\end{proof}


\begin{proposition}\label{prop:quartic-cubic-ramification}
If $A$ has residue field $\mathbb F_2$ and ramification index $e=3$,
then $h_4(A)\le15$.
\end{proposition}

\begin{proof}
For any uniformizer $\varpi$, the unit $\varpi^3/2$ has a cube root
in $A$, by simple-root Hensel lifting.  We may therefore choose a
uniformizer $\pi$ with $\pi^3=2$.
Repeat the universal coefficient calculation with constant coefficient
modulo $16$, instead of modulo $8$.  Denote the resulting group by
$Q_{16}$.  The proof of the kernel assertion now gives kernel $K$,
and the same subgroup $H$ and the same power-image lists give
\[
 Q_{16}/H\simeq\mathbb Z/16\mathbb Z\oplus
 (\mathbb Z/4\mathbb Z)^2\oplus(\mathbb Z/2\mathbb Z)^2,
\]
with generators $u,x,\alpha,\beta,\zeta$ as before; only the order
of $u$ changes.  On substituting $T=\pi$ one has
\[
 \zeta=2\pi^6=8=8u.
\]
Thus every specialized target may be written with $u$ coefficient
$0\le C\le7$, while retaining a binary $\zeta$ coefficient.
Choose $1\le r\le8$ with $r\equiv C\pmod8$.  When $C=0$ and
$r=8$, the extra $8u$ simply flips the required $\zeta$ coordinate.
The rest of the positive construction in
Lemma~\ref{lem:quartic-universal-quotient} is unchanged: at most
four nonunit types supply the $\alpha,\zeta$ coordinates for
$r\ge3$, while the cases $r=1,2$ cost at most seven terms in total.
Hence every class modulo $K_A$ has a representation with at most
twelve fourth powers.  The same three-term baseline cancellation
and Lemma~\ref{lem:primitive-lifting} give $h_4(A)\le15$.
\end{proof}

\subsection{The ninth-power bound}
\label{sec:higher-ninth-local}

Only residue characteristic three requires attention. Use $p=3$,
$a=2$, and coordinate orders $N=(3,5,5)$. In $W_3(R)$ take centers
$(b_0,b_1,b_2)=(0,1,T)$ and perturbations
\[
 (\delta_0,\delta_1,\delta_2)
       =(T^3u_0,T^5u_1,T^3u_2).
\]
For $H=\sum[b_i+\delta_i]-\sum[b_i]$, the universal ghost
recurrence gives $H_j\in T^{N_j}R$. After dividing these coordinates
by $T^3,T^5,T^5$, the reduction modulo $T$ is
\[
 (u_0+u_2,\ u_1+u_2,\ u_1).
\]
Indeed the linear term of $H_1$ is $\sum b_i^2\delta_i$, and that
of $H_2$ is $\sum b_i^8\delta_i$; the other terms have strictly
higher $T$-orders. The displayed linear map has determinant $-1$,
so formal Hensel inversion covers every element of $I_N$.
The perturbed second root is $1+T^5u_1$, hence a unit.

The image lengths are $L_0,L_1,L_2=13,4,1$, so the positive
coefficient cost is
\[
 C_3(3,5,5)=13+2(4+3\cdot1)=27.
\]
The three baseline terms give thirty ninth powers modulo $27A$.
A unit root absorbs the remaining error: the polynomial
\[
 \frac{(1+3z)^9-1}{27}\equiv z\pmod\pi
\]
is integral and has a unit derivative. Hensel's lemma gives the
exact primitive representation with the same number of summands.
At all other residue characteristics, finite-field compression of
$\nu-1$, appending one, and ordinary Hensel lifting use at most ten
terms. Thus $h_9(A)\le30$ without any restriction on residue degree
or ramification.

\subsection{A primitive eighth-power bound of thirty-five}

\begin{theorem}\label{thm:local-eighth-thirtyfive}
For the integer ring $A$ of every finite extension of $\mathbb Q_p$,
\[
 h_8(A)\le35.
\]
More precisely, at dyadic places the following bounds hold:
\[
 \begin{array}{c|c}
 \text{residue degree and absolute ramification}&\text{bound}\\\hline
 f\ge3&19\\
 f=2&30\\
 f=1,\ e=1&32\\
 f=1,\ e=2&32\\
 f=1,\ e=3,4&35\\
 f=1,\ e=5&22\\
 f=1,\ e\ge6&33.
 \end{array}
\]
The entries $e=5$ and $e\ge6$ use the exact finite certificates
described below. Their sizes are respectively $8192$ residue classes
and $64$ polynomial identities plus eight carry relations.
\end{theorem}

Write $\pi$ for a uniformizer. A primitive eighth-power congruence
modulo $16\pi A$ lifts without increasing its number of summands.
Indeed, changing a unit root $x$ to $x(1+2\pi z)$ changes its power by
$16\pi x^8F(z)$, where $F\in A[z]$ and $F(z)\equiv z\pmod\pi$.
All higher binomial terms are divisible by an additional $\pi$ after
division by $16\pi$. Simple-root Hensel lifting applies.
At an odd prime the usual derivative argument gives $h_8(A)\le9$.

\subsection{Residue degrees at least two}

Work first in $W_4(\mathbb F_q[[T]])$. For a centered sum of four
Teichm\"uller vectors with centers $b_i$, its Jacobian has entries
$b_i^{2^j-1}$, $0\le j,i\le3$, with the first row equal to one.
This follows directly by differentiating the universal ghost identities.

If $f\ge3$, take centers $0,c_1,c_2,c_3$, where the three nonzero
constants are linearly independent over $\mathbb F_2$. The zero column
is $(1,0,0,0)^t$; the other minor is an invertible Moore matrix, after
multiplication of its columns by the $c_i$. The centered sum therefore
covers the ideal whose four coordinates belong to $TR$: substitute
perturbations $Tz$, divide the four coordinates by $T$, and apply
formal Hensel inversion to the resulting map with invertible
linear reduction.
By Lemma~\ref{lem:higher-witt-coefficients}, the ambient coefficient cost is
$4+3+2\cdot2+4\cdot1=15$, so the four-variable baseline gives
nineteen powers modulo sixteen.

If $f=2$, choose $\omega^2+\omega+1=0$ and centers $(0,1,\omega,T)$.
The determinant is
\[
 (1+\omega)(T^7+T)=T\cdot\text{a unit}.
\]
Multiplication by the adjugate and the substitution
$\eta=T^2\operatorname{adj}(J)z$ normalize the centered map to
$uz+TP(z)$ after division by $T^3$, with $u$ a unit. Thus the
isotropic ideal of coordinate order three is covered. Its $2$-power
image lengths are $12,6,2,1$, and its coefficient cost is
\[
 12+6+2\cdot2+4\cdot1=26.
\]
Adding the four-variable baseline gives thirty powers modulo sixteen,
by \eqref{eq:higher-last-ghost} and the baseline cancellation above.

To justify the last integral lift, lift the represented roots from
$R=\mathbb F_q[[T]]$ to $D=B[[T]]$. The residual target is in $16D$
by \eqref{eq:higher-last-ghost}. Both baselines retain two units
with distinct eighth-power residues. The final digit modulo
thirty-two can therefore be corrected using these same roots, with
no extra summands. Explicitly,
\[
 (1+2z)^8\equiv1+16(z^2+z)\pmod{32}.
\]
If their lifted roots are $y_1,y_2$, the attainable correction divided
by sixteen is
\[
 \bar y_1^8(z_1^2+z_1)+\bar y_2^8(z_2^2+z_2)\quad\text{in }R.
\]
The Artin--Schreier map $z\mapsto z^2+z$ is bijective on $TR$ and
has image the absolute-trace-zero hyperplane on $\mathbb F_q$.
Multiplication by either unit $\bar y_i^8$ preserves $TR$.
Its two constant images are distinct trace hyperplanes because the
nonzero constants $\bar y_1(0)^8$ and $\bar y_2(0)^8$ are distinct;
the nondegenerate trace pairing identifies their annihilator lines
with their inverse scalars. The sum of these images is consequently
all of $R$. Choose $z_1,z_2$, lift them to $D$, and replace the two
roots by $y_i(1+2z_i)$. This matches the target modulo $32D$ while
preserving both unit roots. After specialization, $32A\subseteq16\pi A$,
so the primitive lifting argument already proved completes the exact
representations with nineteen and thirty terms.

\subsection{A seven-variable weighted chart}

Let $D=\mathbb Z_2[[T]]$ and set
\[
 I=\{(z_0,z_1,z_2,z_3)\in W_4(D):
 z_0\in T^2D,\ z_1\in T^4D,\ z_2,z_3\in T^6D\}.
\]
The inequalities $2^{j-i}N_i\ge N_j$ for
$N=(2,4,6,6)$ show, by weighted homogeneity of the universal Witt
operations, that $I$ is an ideal. The same definition applies after
reduction modulo two.

\begin{lemma}\label{lem:eighth-seven-chart}
Every $w\in I$ has a representation
\[
 w=\sum_{i=0}^6[x_i]-\sum_{i=0}^6[b_i],
 \qquad (b_0,\ldots,b_6)=(0,1,T,T^2,0,0,0),
\]
where $x_1\equiv1\pmod{T}$.
\end{lemma}
\begin{proof}
Use perturbations
\[
 \delta_1=T^6u_1,\qquad\delta_2=T^3u_2,\qquad
 \delta_i=T^2u_i\quad(i=0,3,4,5,6).
\]
Write $H=\sum[b_i+\delta_i]-\sum[b_i]$ and put
$v=(u_0,u_3,u_4,u_5,u_6)$. The universal ghost recurrence shows that
$H_j$ is divisible by $T^{N_j}$. Modulo $(2,T)$, the normalized map is
\[
 \left(\sigma_1(v),\ u_2+u_3+\sigma_2(v),\ u_1+u_2,\ u_1\right).
\]
For the second component, one may use the exact formula
$H_1=\sum b_i\delta_i-\sum_{i<j}\delta_i\delta_j$.
In the last two components, all nonlinear ghost terms have orders
strictly greater than six, leaving the displayed linear terms.

The last two equations fix $u_1,u_2$. For the first two, it is enough
to solve $\sigma_1(v)=s$ and $\sigma_2(v)+u_3=c$ nonsingularly
over $\mathbb F_2$. For $(s,c)=(0,0),(0,1),(1,0),(1,1)$, choose
respectively zero, two, one, and three ones among the four coordinates
other than $u_3$, keeping $u_3=0$. The two Jacobian rows are independent:
their difference from $1$ is the vector $v+e_3$, which is neither zero
nor the all-one vector for these choices. The full Jacobian has rank
four. First lift the chosen solution through the $2$-adic coefficient
ring, then through the $T$-adic parameter, freezing the unused variables.
This solves every normalized target in $D^4$ and proves the lemma.
The root $x_1=1+T^6u_1$ remains a unit.
\end{proof}

\subsection{The finite coefficient certificate}

Put $\bar D=\mathbb F_2[[T]]$ and let $G$ be the actual additive
Teichm\"uller span in $W_4(\bar D)/\bar I$.
The following eight actual generators suffice:
\[
 [1],\ [1+T^5],\ [1+T^4],\ [1+T^3],\ [1+T^2],\
 [1+T],\ [T],\ [T+T^2].
\]
In this order, their successive cyclic quotient orders are
\[
 16,2,2,2,2,4,4,2.
\]
Consequently every element of $G$ is a positive sum of at most
\[
 15+1+1+1+1+3+3+1=26
\]
of the displayed generators.

Here is the precise finite verification used for that assertion.
The combined ghost map sends a vector to $2^{3-j}w_j$, truncated
at degree $N_j$, modulo sixteen for each $j=0,1,2,3$.
It is additive and has kernel exactly $\bar I$, by induction on the
coordinates. Every root can be reduced modulo $T^6$, leaving exactly
sixty-four binary polynomials. The program in Appendix~\ref{app:eighth-ghost} checks their images
against the $8192$ normal forms with the coefficient ranges just
displayed. It also checks all eight carry relations, which make this
normal-form set an additive subgroup. If the generators above are
$g_0,\ldots,g_7$, these relations are
\[
 \begin{split}
 16g_0&=0,\quad 2g_1=2g_0,\quad2g_2=2g_0,\quad2g_3=2g_0,\\
 2g_4&=g_0+g_2,\quad4g_5=3g_0+g_2,\\
 4g_6&=12g_0+g_2+g_4+2g_5,\quad2g_7=2g_6.
 \end{split}
\]
The program also emits a coefficient witness for every binary root. All checks use exact integer arithmetic modulo sixteen.

\subsection{High ramification}

Suppose $f=1$ and $e\ge6$. Lift an eligible target through a finite
signed-power expression to $D$ and then to a Witt vector $g\in W_4(D)$
whose last ghost is that expression. Subtract the seven-term baseline,
reduce modulo two, and apply the preceding certificate modulo $\bar I$.
This uses at most twenty-six positive Teichm\"uller vectors. After
lifting them to $D$, the remaining Witt vector $\eta$ reduces into
$\bar I$. Choose $w\in I$ lifting this reduction. The last ghost of
$\eta-w$ is divisible by sixteen.

After specialization $T=\pi$, this remaining error belongs to
$16A\subseteq8\pi^6A$, since $e\ge6$. It can therefore be matched
exactly by the last ghost of a vector $V^3(T^6z)\in I$, after lifting
its coefficient from $A$ to $D$. Add that vector to $w$ and apply
Lemma~\ref{lem:eighth-seven-chart}. The fixed baseline cancels, leaving
at most $26+7=33$ positive eighth powers, one with a unit root.


\subsection{Small ramification}

The rational ring has $h_8(\mathbb Z_2)=32$: unit eighth powers are
one modulo thirty-two, nonunit eighth powers vanish there, and a
primitive zero needs thirty-two unit terms. Any residue class is
represented by between one and thirty-two ones, followed by lifting.

For $e=2,3,4$, work modulo $16\pi A=\pi^{4e+1}A$.
The eighth power of every unit is $1+d$ with
\[
 v_\pi(d)\ge7\quad(e=2,3),\qquad
 v_\pi(d)\ge9\quad(e=4).
\]
For $e=2,4$, the two lowest terms in the binomial expansion cancel
modulo $\pi$, since the residue field is $\mathbb F_2$; direct use
of $(1+\pi y)^8$ gives these bounds. The product of any two such
deviations vanishes at the working precision. Thus the deviations
form an additive subgroup $D_0$: multiplication of unit eighth powers
corresponds to addition of deviations.

When $e=2$, every nonunit eighth power is zero or $\pi^8$, and
$\pi^8\equiv16$ at this precision. The generated image is consequently
$\mathbb Z\cdot1+D_0$. It is represented by at most thirty-two unit
eighth powers, using one unit power for the prescribed deviation.

When $e=3$, choose $\pi^3=2$, by extracting the cube root of the
unit in any initial uniformizer relation. The nonunit images are
zero and $\pi^8$, and $4\pi^8=0$ at precision $\pi^{13}$.
Thus at most three copies of $\pi^8$, together with at most
thirty-two unit eighth powers, represent every eligible class.

When $e=4$, the nonunit images are zero, $\pi^8$, and $\pi^{16}$.
Writing $\pi^4=2u$ shows
\[
 \pi^{16}\equiv16,\qquad4\pi^8\equiv16\pmod{\pi^{17}}.
\]
The same bound of thirty-five follows. All these representations
have a unit root and lift.

Finally, if $e=5$, extracting a fifth root of a principal unit gives
$A=\mathbb Z_2[\pi]$ with $\pi^5=2$. Eighth powers modulo $16\pi$
depend only on roots modulo $2\pi$. The latter quotient has sixty-four
classes, represented by
\[
 a_0+a_1\pi+\cdots+a_4\pi^4,
 \qquad0\le a_0<4,\quad a_i\in\{0,1\}\ (i>0).
\]
The program in Appendix~\ref{app:eighth-tame} checks these classes in the exact ring
\[
 (\mathbb Z/32)\oplus
 (\mathbb Z/16)\pi\oplus\cdots\oplus(\mathbb Z/16)\pi^4,
 \qquad\pi^5=2.
\]
There are twenty distinct eighth-power images. Their generated group
has $8192$ elements and ordinary additive diameter twenty-one.
The breadth-first calculation verifies every parent representation
and closure under every generator. Applying it to the target minus
one and appending $1^8$ gives a primitive bound twenty-two.
This exhausts the ramification cases and proves
Theorem~\ref{thm:local-eighth-thirtyfive}.

\section{Efficient differencing with refined seed moments}
\label{sec:efficient-differencing}\label{sec:g6-transfer}\label{sec:quartic-principal-blocks}
The efficient-differencing method of Wooley and Vaughan--Wooley
\cite{WooleyThesis,Wooley1992,VW1995,VW2000} can be formulated
using principal blocks of prime ideals. At these fixed small exponents
the terminal estimates also use the number-field Vinogradov mean-value
theorem and coefficient-uniform curve bounds.
\subsection{Localizing the ordinary variables}

One may prove a permissible full-box moment by proving it for smooth roots in
arbitrarily small fixed boxes away from zero at every embedding.  Here are the
details that allow this reduction uniformly in the scale.
Partition nonzero roots into dyadic ranges for each $|x_v|$.  There are
$O((\log X)^{r_1+r_2})$ ranges, since $|Nx|\ge1$ supplies a polynomial lower
bound for every $|x_v|$.  Roots with $|Nx|\le X^{1/2}$ number
$O(X^{1/2+\epsilon})$, by counting ideals and bounded-height units; their
$2s$th moment is trivially $O(X^{s+\epsilon})$.
For each remaining annular box let $Y$ be the product of its archimedean radii,
with local dimensions counted.  Thus $X^{1/2}\ll Y\ll X$.
Multiplication by one unit balances these radii to fixed multiples of $Y^{1/n}$.
It preserves ideal smoothness and the mean value, since multiplication of
$\alpha$ by a unit's $k$th power is a torus automorphism.
Cover the resulting compact annulus by finitely many boxes of relative diameter
$\delta$, where $\delta>0$ is a sufficiently small constant depending only on
$K$ and the maximal differencing depth.  H\"older's inequality for the sum of
these pieces costs only a power of $\log X$.  Since $R=X^\eta\le Y^{2\eta}$,
a sufficiently small initial smoothness exponent works for all pieces.

For one such piece, replace two smooth roots by ordinary roots in that same
small box.  The other smooth roots can subsequently range over full boxes.
The two distinguished ordinary roots, and every pair produced from them by
efficient differencing, retain their actual small support.
In centered notation define
\[
 \Psi_j(z;\mathbf h;\mathbf m)=
 \Delta^*_j\bigl((z-\textstyle\sum_{i\le j}h_i m_i^k)^k;
                 2\mathbf h;\mathbf m\bigr),
 \qquad
 \Delta^*_1 f(z;h;m)=m^{-k}(f(z+hm^k)-f(z)).
\]
These are integral polynomials.  Their leading term is
$2^j(k)_j(h_1\cdots h_j)z^{k-j}$.
At each embedding, the actual support has $z$ within $O_j(\delta P)$ of a
fixed nonzero center of size comparable to $P$, and
\begin{equation}\label{eq:g6-small-increments}
 |h_{i,v}m_{i,v}^k|\le C_j\delta P\qquad(1\le i\le j).
\end{equation}
Fixed factors of two change the center and the $C_j$, harmlessly.
The central change of variables is $z=z_1+z_2$,
$m^kh=z_1-z_2$; divisibility by $2$ can be retained as a support restriction.
More explicitly, for exponent $k$ the centered polynomial identity multiplies
the differenced equation by $2^k$, doubles the earlier shift parameters, and
replaces every remaining quotient root $x$ by $2x$.
This identity applies at every repeated step.
The retained parity restrictions only reduce a solution count; subsequently
enlarging the bounding boxes changes their scales by fixed constants, and the
primes above $2$ belong to the smoothness range for sufficiently large $X$.
Consequently $\delta$ can be chosen so that, whenever all $h_i\ne0$,
\begin{equation}\label{eq:g6-nonzero}
 \Psi_j(z;\mathbf h;\mathbf m)\ne0,
 \qquad \partial_z\Psi_j(z;\mathbf h;\mathbf m)\ne0
 \quad(0\le j\le k-1).
\end{equation}
Indeed, after dividing by $h_1\cdots h_j$, the leading monomial dominates
all remaining monomials at each embedding.  The same is true for the derivative.
An ordinary differencing shift is a difference of two points in the small box;
it therefore retains \eqref{eq:g6-small-increments}.

In all subsequent counts the coupled restrictions
\eqref{eq:g6-small-increments} are retained.  The separate bounding boxes for
$h_i,m_i$ are used only to count parameter choices, never to replace the support
when \eqref{eq:g6-nonzero} is needed.  This is the step that removes the
positivity assumption implicit in the rational auxiliary counts.

\subsection{The fundamental efficient-differencing inequality}

Let $F_\Psi(\alpha)=\sum_{z,\mathbf u}\mathrm e_K(\alpha\Psi(z,\mathbf u))$,
where $z$ ranges in a box of volume $X$, the auxiliary parameter tuples number
at most $U$, every polynomial has bounded positive degree and coefficients of
polynomial height, and $\Psi'(z,\mathbf u)\ne0$ on the retained support.
Let $g_Q$ be the full smooth sum of volume $Q$, and put
$S_s=\int|F_\Psi|^2|g_Q|^{2s}$.
Choose $1\le M\le Q$.

Let $h_K$ be the order of the ideal class group.  Any sequence of $h_K$ ideal
classes has a nonempty zero-sum subsequence.  Partition the prime ideal factors
of a smooth ideal into such principal blocks, each with at most $h_K$ factors,
leaving fewer than $h_K$ factors.  If the part of $(x)$ coprime to
$(\Psi'(z,\mathbf u))$ has norm greater than $MR^{h_K}$, successive principal
blocks supply a divisor $(m)$ of $(x)$, coprime to $\Psi'$, with
\[
 M<N(m)\le MR^{h_K}.
\]
Choose a generator balanced by units.  There are $O(MR^{h_K+\epsilon})$
possible generators, and $x/m$ is a smooth algebraic integer in a fixed box
of volume $Q/M$.  Powers of $R$ will always be absorbed in $X^\epsilon$.

If no such divisor exists, the prime-to-$\Psi'$ part of $(x)$ has norm at most
$MR^{h_K}$.  For fixed $z,\mathbf u$ the number of these bad $x$ is at most
$MX^\epsilon$: count the prime-to-$\Psi'$ ideal and use the supported-ideal
bound for its complementary ideal, then count bounded-height generators.
This bound also covers roots of norm at most $MR^{h_K}$.

The contribution of a bad variable is at most
$S_s^{1/2}(QMX^\epsilon S_{s-1})^{1/2}$.
Indeed, if
$G_x=\sum_{z,\mathbf u:,x\text{ bad}}\mathrm e_K(\alpha\Psi(z,\mathbf u))$,
then
\[
 \left|\sum_x\mathrm e_K(\alpha x^k)G_x\right|^2
 \le O(Q)\sum_x|G_x|^2.
\]
After multiplication by $|g_Q|^{2s-2}$ and integration, a fixed pair of auxiliary
tuples admits at most $MX^\epsilon$ common bad $x$, giving the assertion.

For the good variables, extract a principal divisor separately from each of the
$2s$ smooth roots.  The dependence of this extraction on the derivative is
handled as follows.  Let $\mathcal M$ be the set of balanced extracted divisors,
so that $L=\#\mathcal M\ll MX^\epsilon$.  For $W\ne0$ put
\[
 F_W(\alpha)=\sum_{\substack{z,\mathbf u\text{ on the retained support}\\
              ((W),(\Psi'(z,\mathbf u)))=1}}
       \mathrm e_K(\alpha\Psi(z,\mathbf u)),
 \quad g=g_{Q/M},\quad
 V_m=\int|F_m(\alpha)|^2|g(m^k\alpha)|^{2s}\,d\alpha.
\]
For tuples $\mathbf m,\mathbf n\in\mathcal M^s$ put
$W=\prod_i m_i$ and $W'=\prod_i n_i$.
After enlarging the quotient-root boxes, orthogonality bounds the good count by
the sum over these tuples of
\[
 \int |F_W(\alpha)F_{W'}(\alpha)|
       \prod_{i=1}^s|g(m_i^k\alpha)g(n_i^k\alpha)|\,d\alpha.
\]
H\"older bounds each integral by the product of $2s$ factors of the form
$(\int|F_W|^2|g(m_i^k\alpha)|^{2s})^{1/(2s)}$ and their $W',n_i$ analogues.
The coefficient set of $F_W$ is contained in that of $F_{m_i}$, because
coprimality to $W$ implies coprimality to $m_i$.  Orthogonality, with its
nonnegative solution counts, therefore bounds that integral by $V_{m_i}$.
Consequently the total good count is at most
\[
 \left(\sum_{m\in\mathcal M}V_m^{1/(2s)}\right)^{2s}
 \le L^{2s-1}\sum_{m\in\mathcal M}V_m.
\]
Both derivative coprimalities are present in every $V_m$.

Here is the remaining Cauchy step explicitly.  Set $\mathfrak b=(m)^k$.
For $\rho\in\mathcal O_K/\mathfrak b$ and fixed $\mathbf u$, let
$\mathcal R(\rho,\mathbf u)$ be the nonsingular residue classes $r$ satisfying
$\Psi(r,\mathbf u)\equiv\rho\pmod{\mathfrak b}$.
Their number is at most $D^{\omega((m))}\ll X^\epsilon$, by Hensel lifting
and the Chinese remainder theorem.  Write
$F_{\mathbf u,r}=\sum_{z\equiv r\, (\mathfrak b)}
\mathrm e_K(\alpha\Psi(z,\mathbf u))$, retaining the original support.
Orthogonality kills pairs with different $\Psi$ residues, and hence
\[
 V_m=\sum_\rho\int
 \left|\sum_{\mathbf u}\sum_{r\in\mathcal R(\rho,\mathbf u)}
                   F_{\mathbf u,r}(\alpha)\right|^2
 |g(m^k\alpha)|^{2s}\,d\alpha.
\]
Cauchy's inequality inside each square costs at most $UX^\epsilon$ and leaves
the sum of $|F_{\mathbf u,r}|^2$.  Its underlying equations have the same
auxiliary tuple on both sides and $z\equiv w\pmod{(m)^k}$.
Write $z-w=m^kh$.  The case $h=0$ contributes at most
$XUM(Q/M)^{\lambda_s+\epsilon}$; the other cases are bounded by
$T_{\rm next}=\int|F_{\rm next}||g_{Q/M}|^{2s}$, with all actual support
restrictions retained.  Thus the good contribution is at most
\[
 X^\epsilon U M^{2s-1}
 \{XUM(Q/M)^{\lambda_s}+T_{\rm next}\}.
\]

The bad contribution can be absorbed into this same explicit diagonal bound
without a lower bound for the density of smooth algebraic integers.
Indeed, $S_0\ll XU^2$, by fixing $z,\mathbf u,\mathbf v$ and solving the
nonconstant equation for $w$.  H\"older gives
$S_{s-1}\le S_0^{1/s}S_s^{1-1/s}$.
If the bad contribution dominates, then
\[
 S_s\ll X^\epsilon XU^2(QM)^s
 \le X^\epsilon XU^2M^{2s}(Q/M)^{\lambda_s},
\]
since $Q\ge M$ and $\lambda_s\ge s$.
In every case, therefore,
\begin{equation}\label{eq:g6-fundamental}
 S_s\ll X^\epsilon U M^{2s-1}
 \{XUM(Q/M)^{\lambda_s}+T_{\rm next}\}.
\end{equation}
This is the required field version of the fundamental efficient-differencing
inequality.  Its proof is the congruence argument of Wooley's 1990 thesis,
Chapter~6, with principal ideal blocks and the preceding explicit treatment
of bad variables and the diagonal term.

\subsection{Auxiliary seed moments}
\begin{lemma}[Sparse progression endpoint over a number field]
\label{lem:field-sparse-endpoint}
Let $p\ge2$ and $J\ge1$ be fixed integers. For $1\le s\le J+1$,
\[
 \#\left\{\mathbf x,\mathbf y\in\mathcal B(M)^s:
       \sum_{i=1}^s x_i^{pr}=\sum_{i=1}^s y_i^{pr}\ (1\le r\le J)\right\}
 \ll_{K,p,J,\epsilon} M^{s+\epsilon}.
\]
\end{lemma}
\begin{proof}
For $s\le J$, Newton's identities determine the multiset of $p$th powers
on one side from the other. Each powered coordinate has at most $p$ roots,
so there are $O_{p,J}(M^s)$ solutions. The case $J=1,s=2$ is the ordinary
fourth-moment bound for the degree-$p$ sum, supplied by Hua's differencing
lemma over $K$.

Suppose $J\ge2$ and $s=J+1$. Interchanging the first coordinates of the two
sides turns the equations into equality of signed power sums
\[
 S_r(\mathbf x)=\sum_{i=1}^J x_i^{pr}-x_0^{pr}
             =S_r(\mathbf y)\qquad(1\le r\le J).
\]
These equalities imply
\[
 D(\mathbf x)=\prod_{i=1}^J(x_0^p-x_i^p)
             =\prod_{i=1}^J(y_0^p-y_i^p)=D(\mathbf y).
\]
Here is a direct algebraic verification. If $x_0^p=y_0^p$, apply Newton's
identities to the other $J$ powered coordinates. Otherwise the two monic
polynomials with roots $(y_0^p,x_1^p,\ldots,x_J^p)$ and
$(x_0^p,y_1^p,\ldots,y_J^p)$ differ by a constant, because their first $J$
power sums agree. Evaluating their difference at $x_0^p$ and $y_0^p$ gives
$(x_0^p-y_0^p)D(\mathbf x)=(x_0^p-y_0^p)D(\mathbf y)$.

Fix $\mathbf y$ with $D(\mathbf y)\ne0$. Every difference
$d_i=x_0^p-x_i^p$ is a bounded-height divisor of this fixed nonzero algebraic
integer. There are $O(M^\epsilon)$ choices for the whole difference tuple.
Moreover
\[
 (J-1)x_0^p=S_1(\mathbf y)+\sum_i d_i,
\]
so the powers of all coordinates, and hence the coordinates up to a fixed
multiplicity, are determined. This gives $O(M^{J+1+\epsilon})$ solutions.
If $D(\mathbf y)=0$, choose a coordinate with $y_i^p=y_0^p$; likewise one
coordinate of $\mathbf x$ has the same power as $x_0$.
Cancel these in the signed sums. The remaining $J-1$ powered coordinates
are determined up to permutation by Newton's identities. There are
$O(M^J)$ choices for $\mathbf y$ and $O(M)$ choices for the free center
$x_0$, again with bounded root multiplicities. This proves the lemma.
\end{proof}

For efficient differences choose $0\le\phi_i<1/k$ and let
\[
 M_i=X^{\phi_i},\quad H_i=X M_i^{-k},\quad
 \Phi_j=\sum_{i=1}^j\phi_i,\quad
 U_j=\prod_{i=1}^jH_iM_i.
\]
Fixed small constants and fixed smoothness powers are suppressed in these
scale definitions. The centered polynomial is denoted by $\Psi_j(z;\mathbf
h,\mathbf m)$, and $F_j$ is its sum over the small ordinary box and the
independent shift and smooth-coefficient boxes. The original ordinary
vertices remain in a common small box away from zero at every embedding.
Thus $\Psi_j/\prod h_i$ and its required derivatives are nonzero wherever
they are used. This is also true after the finitely many ordinary differences
below, which have small shifts and coefficient divisor $1$.

\begin{lemma}[Sparse moments from a nonvanishing Wronskian]
\label{lem:field-sparse-wronskian}
For fixed integers $p\ge2$ and $J\ge1$, put
\[
 \tau_J=\max\{J+1,J(J+1)/2\}.
\]
For every integer $1\le s\le\tau_J$, the number of solutions in a
box of volume $M$ to
\[
 \sum_{i=1}^s x_i^{pa}=\sum_{i=1}^s y_i^{pa}
 \qquad(1\le a\le J)
\]
is $O_{K,p,J,\epsilon}(M^{s+\epsilon})$.
\end{lemma}
\begin{proof}
The range $s\le J+1$ is Lemma~\ref{lem:field-sparse-endpoint}.
For $s\le J(J+1)/2$, apply the number-field mean-value theorem
of Wooley~\cite[Theorem~15.1]{Wooley2019} to the $J$ polynomials
$t^p,t^{2p},\ldots,t^{Jp}$. Their Wronskian is a nonzero monomial:
its coefficient is the product of their nonzero degrees and a
Vandermonde determinant in these distinct degrees.
The theorem applies to arbitrary complex coefficients, so restricting
to the prescribed box and using orthogonality gives the assertion.
\end{proof}

\begin{lemma}[Centered curve intersections]
\label{lem:centered-curve-intersections}
Fix $d\ge3$ and a nonzero leading coefficient $c\in K$.
Let $P,Q\in K[t]$ have degree $d$, leading coefficient $c$, and
vanishing coefficients of $t^{d-1}$. Suppose their coefficients
have fixed denominators and archimedean height bounded by a fixed
power of $X$. Let $\mathcal B$ be a sufficiently small box of
volume comparable to $X$ about a fixed nonzero archimedean center
at scale $X^{1/[K:\mathbb Q]}$. Then
\[
 \#\{(z,w)\in(\mathcal B\cap\mathcal O_K)^2:P(z)=Q(w)\}
 \ll X\,\mathbf1_{P=Q}+X^{1/3+\epsilon},
\]
uniformly in the lower coefficients of $P,Q$.
\end{lemma}
\begin{proof}
Factor $P(z)-Q(w)$ into irreducible factors over $K$.
Its highest homogeneous part is $c(z^d-w^d)$, which is square-free.
Every linear factor has the form $z-\zeta w-b$, with $\zeta^d=1$.
Substitution and comparison of the coefficient of $w^{d-1}$ give
$b=0$. Choosing the box smaller than a fixed fraction of the
distance between its center and its nontrivial root-of-unity
multiples excludes $z=\zeta w$ for $\zeta\ne1$.
The line $z=w$ occurs precisely when $P=Q$, and contributes $O(X)$.

The highest homogeneous part of a quadratic factor is the product
of two distinct factors of $z^d-w^d$. After normalizing one
coefficient it belongs to a fixed finite collection of nondegenerate
binary quadratic forms. Its lower coefficients have polynomial
height and fixed denominators. To see the denominator assertion,
clear the fixed leading coefficient by a fixed change in $z$;
the resulting polynomial is monic in $z$ over $\mathcal O_K[w]$.
Its monic factors belong to this integrally closed ring.
Standard fixed-degree factor bounds give the height assertion.
Translation to the center of the quadratic form therefore changes
the lattice by a fixed denominator only. The resulting equation is
a norm equation in one of finitely many fixed quadratic extensions
of $K$, or a product equation over $K$.
For a nonzero right side, the ideal-divisor estimate and the
bounded-height unit count give $O(X^\epsilon)$ solutions.
For a zero right side, an irreducible quadratic factor has at most
one $K$-point; split linear factors have already been treated.

Every remaining irreducible component has degree at least three.
The affine-curve estimate of Binyamini--Cluckers--Kato
\cite[Corollary~12]{BCK2025} bounds its integral points by
$O_K(X^{1/3}(\log X)^\kappa)$, uniformly in its coefficients.
Here the height parameter in that result is $O_K(X)$: the
embedding bound is $O_K(X^{1/[K:\mathbb Q]})$, as in their
definition preceding equation~(28).
There are at most $d$ components. Summing their contributions
proves the lemma.
\end{proof}

\begin{lemma}[Auxiliary moments for efficient differencing]
\label{lem:field-general-seed-aux}
Let $1\le\ell\le3$ and $1\le j\le k-1-\ell$, and put
\[
 J=\left\lfloor\frac{k-j-\ell+1}{2}\right\rfloor,
 \qquad \tau_J=\max\{J+1,J(J+1)/2\}.
\]
Assume the chronological divisor inequalities
\begin{equation}\label{eq:field-chronological-prefix}
 \Phi_i+(k-2)\phi_i\le1
 \qquad(\tau_J<i\le j).
\end{equation}
This condition is empty when $j\le\tau_J$.
Assume also that $\Phi_j\le2/3$ whenever at least one of the
residual degrees $k-j-r$, $0\le r<\ell$, is even and at least six.
Then
\[
 \int_{\mathbb T_K}|F_j(\alpha)|^{2^\ell}\,d\alpha
 \ll X^{2^\ell-\ell+\epsilon}U_j^{2^\ell-1}.
\]
\end{lemma}
\begin{proof}
Consider first a second moment after $r$ appended ordinary
differences, where $0\le r<\ell$. The retained small supports
make the divided phase nonzero. Grouping its nonzero value by
the divisor tuple and applying Cauchy's inequality costs
$X^\epsilon$ and reduces the count to a common tuple of efficient
and ordinary shifts on the two sides.
Polynomial coincidence implies
\begin{equation}\label{eq:seed-coefficient-coincidence}
 \sum_{i=1}^j h_i^{2a}(m_i^{2ka}-n_i^{2ka})=0
 \qquad(1\le a\le J).
\end{equation}
The implication follows by taking the logarithm of the centered
generating function
\[
 e^{zt}\prod_i\frac{\sinh(h_im_i^kt)}{h_im_i^kt}.
\]
Euler's product gives
\[
 [t^{2a}]\log(\sinh t/t)=
       \frac{(-1)^{a+1}\zeta(2a)}{a\pi^{2a}}\ne0.
\]
The ordinary coefficient blocks are equal on the two sides.

When $j\le\tau_J$, fix the shifts and apply H\"older to the $j$
coefficient sums on $\mathbb T_K^J$.
Multiplication of coordinate $a$ by $h_i^{2a}$ preserves probability
Haar measure. Lemma~\ref{lem:field-sparse-wronskian}, with $p=2k$,
therefore gives
\[
 \prod_{i=1}^j\left(M_i^{j+\epsilon}\right)^{1/j}
 \ll X^\epsilon\prod_{i=1}^j M_i
\]
coefficient pairs, including unequal coefficient volumes.
Summing the shifts gives $O(X^\epsilon U_j)$.

For $j>\tau_J$, eliminate the last coefficient block using the
first equation of \eqref{eq:seed-coefficient-coincidence}.
If $m_i^{2k}=n_i^{2k}$, at most $2k$ choices of $n_i$ occur for
each $m_i$, and all $J$ equations reduce to those of the preceding
blocks. Otherwise
$h_i^2(m_i^{2k}-n_i^{2k})=b\ne0$ has $O(X^\epsilon)$ solutions
after the other blocks are fixed: first count divisors $h_i$ of $b$,
then divisors $m_i-n_i$ of the remaining fixed integer, and finally
roots of a nonzero fixed-degree polynomial.
Consequently the coefficient count satisfies
\[
 N_i\ll X^\epsilon\left(H_iM_iN_{i-1}
                    +\prod_{a<i}H_aM_a^2\right).
\]
The inequality in \eqref{eq:field-chronological-prefix} is exactly
$\prod_{a<i}M_a\le H_iM_i$.
Stop the elimination at $i=\tau_J$ and apply the sparse estimate
there. Induction gives $N_j\ll X^\epsilon U_j$.

For odd residual degree, the centered phase has a factor $z$.
Divisor grouping in its nonzero value makes the two centers equal
at cost $X^\epsilon$; a nonzero difference of the two polynomials
then has boundedly many roots. This gives
$O(X^{r+1+\epsilon}U_j)$, since $\prod_iM_i\le X$.
At residual degree two, use a difference of squares. At degree
four, complete the squares in $z^2,w^2$; a nonzero level is a
divisor equation, and a zero level gives either polynomial
coincidence or a binary quadratic norm equation. The nontrivial
root-of-unity lines miss the common small supports. The same
bound follows in both cases.

For the remaining even residual degrees, divide the common-shift
polynomial by its leading coefficient. The two resulting centered
polynomials have the same fixed leading coefficient, vanishing
next coefficient, and fixed denominators. Their lower coefficients
have polynomial height. Lemma~\ref{lem:centered-curve-intersections}
gives
\[
 R_{j,r}\ll X^{r+1+\epsilon}U_j+
 X^{r+1/3+\epsilon}\Bigl(\prod_iH_i\Bigr)
                         \Bigl(\prod_iM_i\Bigr)^2.
\]
The second term is absorbed when $\Phi_j\le2/3$.
The padded ordinary-moment induction below now gives the stated
$2^\ell$th moment.
\end{proof}

\paragraph{The padded ordinary-moment induction.}
Here is an explicit induction from the appended second moments to the
higher moments used by A and B.  Let $F_r$ denote the family with $r$
ordinary differences already appended, and set $U_r=X^rU_j$.
Cauchy in the parameter tuple and ordinary differencing give
\[
                     |F_r|^2\ll XU_r^2+U_r|F_{r+1}|.
\]
The zero ordinary shift is the first term; all nonzero shifts retain the
small-support restrictions.  Put
$B_{r,a}=\int|F_r|^{2^a}$.  For $a\ge2$, multiply the pointwise inequality
by $|F_r|^{2^a-2}$.  Interpolation in the first term and H\"older in the
second, followed by absorption, yield
\[
 B_{r,a}\ll (XU_r^2)^{2^{a-2}}B_{r,a-1}
                  +U_r^{2^{a-1}}B_{r+1,a-1}.
\]
Starting with $B_{r,1}\ll X^{1+\epsilon}U_r$, induction gives
\[
                  B_{r,a}\ll X^{2^a-a+\epsilon}U_r^{2^a-1}.
\]
Indeed, the two powers of $X$ obtained on substitution are
$3\cdot2^{a-2}-a+1$ and $2^a-a$, and the former is no larger than the
latter.  The restriction $j+\ell\le k-1$ means that every final residual
degree is at least two; the extra multiplicity from a residual linear
phase never occurs.  The proof retains the parameter tuple throughout the
ordinary differencing and consequently introduces no new independent
copies of the efficient-divisor variables.

\paragraph{Preserving the supports.}
All arguments above may be performed on the original supports.  If
Cartesian boxes are wanted, split the balanced extracted generators into
dyadic norm annuli and then small archimedean boxes.  On one piece their
conjugates have comparable fixed scales.  The shift ranges can be enlarged
by fixed factors while all ordinary vertices remain in a slightly larger
common tiny box.  The enlargement is made in nonnegative orthogonality
counts before an absolute-value bound is used.  The number of pieces is
$X^\epsilon$ and changes in scale are $X^{O(\eta)}$.
Choose the initial relative diameter small enough for the fixed total
number of differences.  The leading monomial of every required residual
polynomial, and of its relevant derivative, then dominates at every
embedding.  This justifies all exclusions of zero phase values above.

\subsection{The exponent transformations}
Write $b_s$ for a permissible full exponent at moment $2s$, and
$b(t)$ for linear interpolation between adjacent integer indices.
Each seed row specifies $s$, a depth $j$, an order $\ell$, an A or B
terminal, and rational parameters $\phi_1,\ldots,\phi_j$. Put
\[
 q_i=1-\Phi_i,\qquad
 \kappa_\ell=2^\ell-\ell+(2^\ell-1)(j-(k-1)\Phi_j),
 \qquad u=s-j.
\]
The A terminal is
\begin{equation}\label{eq:seed-A-terminal}
 C_j=2^{-\ell}\kappa_\ell+
 (1-2^{-\ell})q_j b\bigl(u/(1-2^{-\ell})\bigr).
\end{equation}
For the B terminal choose an additional integer $t$ and put
$c=u/t$, $a=2(1-c)-2^{1-\ell}$, and
$b'=2^{1-\ell}-(1-c)$. If these weights are nonnegative, then
\begin{equation}\label{eq:seed-B-terminal}
 C_j=a\kappa_{\ell-1}+b'\kappa_\ell+cq_jb(t).
\end{equation}
These follow by H\"older from Lemma~\ref{lem:field-general-seed-aux}.
Proceeding backwards for $i=j-1,\ldots,1$, set
\begin{equation}\label{eq:seed-backwards}
 \begin{split}
 E_i&=\tfrac12\{i-(k-1)\Phi_i+(2(s-i-1)-1)\phi_{i+1}
                              +q_i b(s-i+1)\},\\
 D_i&=i+1-(k-1)\Phi_i+\phi_{i+1}+q_{i+1}b(s-i-1),\\
 C_i&=E_i+\tfrac12\max(D_i,C_{i+1}).
 \end{split}
\end{equation}
The first-stage output is
\begin{equation}\label{eq:seed-initial-output}
 T_s(b)=(2s-3)\phi_1+
       \max\{1+\phi_1+q_1b(s-1),C_1\}.
\end{equation}
The fundamental principal-block inequality proves these backwards
steps. Different terminal estimates may be substituted for $C_j$;
the remainder of the transformation is unchanged.

\begin{lemma}[Finite supersolution principle]\label{lem:supersolution}
Let a finite collection of moment transformations have nonnegative
coefficients in the input exponents, and let the sum of the
coefficients in every branch be at most $\rho<1$. Finite minima
and maxima of these branches are allowed. Suppose that a rational
vector $b$ satisfies $T_s(b)<b_s$ for every nontrivial coordinate,
and agrees with permissible elementary exponents at the remaining
coordinates. Then every $b_s$ is permissible. If outputs are clamped
below by the diagonal exponent $s$, assume also that $b_s\ge s$.
\end{lemma}
\begin{proof}
Start from a vector of elementary permissible exponents dominating
$b$. Monotonicity and the coefficient-sum bound show that, after
one transformation, its excess above $b$ in sup norm is at most
$\rho$ times the previous excess. The same assertion holds after
clamping, since the constant branch has coefficient sum zero.
After finitely many steps the excess is arbitrarily small.
For any prescribed error exponent, choose that number of steps
first and then choose the common smoothness parameter sufficiently
small for all the finitely many estimates. This gives the required
$X^\epsilon$ loss.
\end{proof}

\section{Coefficient sampling}
\label{sec:coefficient-sampling-common}
\begin{lemma}[A bilinear sampling bound]
\label{lem:g789-bilinear}
If $J,N\ge1$, $c(a)$ is supported on $\mathcal B(N)$, and
$\theta\in V$, then
\[
 \sum_{z\in\mathcal B(J)}
       \left|\sum_a c(a)E(\theta za)\right|^2
 \ll_K\left(J+N+\frac{JN}{\mathcal H_{JN}(\theta)}\right)
                         \sum_a|c(a)|^2.
\]
The same assertion holds for arbitrary subsets of the displayed
boxes and arbitrary fixed translates of those boxes.
\end{lemma}
\begin{proof}
First suppose $J\le N$. The multidimensional elementary large-sieve
inequality bounds the left side by $O_K(NM)\sum|c(a)|^2$, where $M$
is the maximum number of sample frequencies $\theta z$ in a translate
of an archimedean box of side $O_K(N^{-1/n})$ modulo
$\mathfrak D^{-1}$. One direct proof chooses a positive Gaussian
majorant for the coefficient box, applies Poisson summation to its
Gram matrix, and bounds its operator norm by the maximum row sum.
Covering successive frequency annuli by boxes of side $N^{-1/n}$
and using Gaussian decay bounds that row sum by $O_K(NM)$.

Partition the $z$ box into a fixed number of smaller boxes if
necessary. The affine field-line lemma in Section~\ref{sec:ordinary-bundle}
applies to the close pairs $(z,a)$, because their volume product is
$O_K(J/N)$, made sufficiently small by this fixed subdivision.
It gives either at most one point or a slope $\gamma\in K$ and
\[
 M\ll_K1+Q_\gamma^{-1}
       \prod_v\min\{J^{1/n},N^{-1/n}/|\theta_v-\gamma_v|\}^{d_v}
 =O_K\left(1+\frac{J}{Q_\gamma A_{JN}(\theta-\gamma)}\right).
\]
This is at most $O_K(1+J/\mathcal H_{JN}(\theta))$.
The large-sieve bound proves the result when $J\le N$.
When $J>N$, duality of the finite matrix
$(E(\theta za))_{z,a}$ interchanges the two boxes and proves the
bound with $J$ in place of $N$. Fixed translates merely insert
unit-modulus row and column factors; subsets can be handled by
zero extension or deletion of rows.
\end{proof}

\begin{lemma}[Counting small norms in an ideal]
\label{lem:g789-smallnorm}
Let $\mathfrak b$ be an integral ideal, $b=N\mathfrak b$, and
$H\ge1$. For $U\ge b$ one has
\[
 \#\{h\in\mathfrak b\cap\mathcal B(H):0<|Nh|\le U\}
       \ll_{K,\varepsilon}H^\varepsilon U/b.
\]
There are no such nonzero points when $U<b$.
\end{lemma}
\begin{proof}
For a counted element, $(h)\mathfrak b^{-1}$ is an integral ideal
of norm at most $U/b$. There are $O_K(U/b)$ such ideals.
For any fixed principal ideal $(h)$, its generators in the box
number $O_K((1+\log H)^{r_1+r_2-1})$: upper bounds for all conjugates,
together with $|Nh|\ge1$, confine their logarithmic vectors to a box
of side $O_K(1+\log H)$, and the unit logarithms form a fixed lattice.
This proves the asserted estimate, including uniformity in
$\mathfrak b$.
\end{proof}

\begin{proposition}[Refined cubic coefficient sampling over $K$]
\label{prop:g789-sampling}
Let $H,J,N\ge1$, put $V=H^3JN$, and choose
\[
 1\le Y\le c_K\min\{J,N,(JN/H^3)^{1/2}\}.
\]
Choose a lifted approximation $\alpha=\gamma+\beta$ with
$L=Q_\gamma A_V(\beta)\ll_K V^{1/2}$, as in
Lemma~\ref{lem:g789-height}. If $c(a)$ is supported on $\mathcal B(N)$,
then
\[
 \sum_{\substack{h\in\mathcal B(H)\\h\ne0}}
 \sum_{z\in\mathcal B(J)}
       \left|\sum_a c(a)E(\alpha h^3za)\right|^2
 \ll_{K,\varepsilon}
 V^\varepsilon JHN\left(Y^{-1}+\frac{H^2}{L}\right)
                          \sum_a|c(a)|^2.
\]
In particular $L$ can be replaced by the smaller intrinsic height
$\mathcal H_V(\alpha)$. In
applications the variables are products of nonzero shifts.
\end{proposition}
\begin{proof}
Write $W=JN$ and apply Lemma~\ref{lem:g789-bilinear} with
$\theta=\alpha h^3$. Since $Y\le\min(J,N)$, the contribution
of $J+N$, and that of all $h$ with $\mathcal H_W(\alpha h^3)\ge Y$,
is $O_K(JHN/Y)\sum|c(a)|^2$.

For each remaining $h$, choose
$\alpha h^3=\gamma_h+\beta_h$ with
$Q_{\gamma_h}A_W(\beta_h)<Y$. The center $\gamma_h/h^3$ is an
approximation to $\alpha$ at scale $V$ with height at most $C_KH^3Y$.
Indeed its conductor norm is at most $|Nh|^3Q_{\gamma_h}$, and
\[
 |Nh|^3 A_V(\beta_h/h^3)
 =\prod_v\max\{|h_v|^3,V^{1/n}|(\beta_h)_v|\}^{d_v}
 \ll_K H^3 A_W(\beta_h).
\]
The upper bound on $Y$ makes the product of this height and $L$
smaller than the separation constant times $V$.
Lemma~\ref{lem:g789-height} therefore gives
$\gamma_h=h^3\gamma$ and $\beta_h=h^3\beta$ exactly.

Let $\mathfrak a$ be the conductor of $\gamma$, $Q=N\mathfrak a$,
and set $A=A_V(\beta)$. The conductor of $h^3\gamma$ is
$\mathfrak a/(\mathfrak a+(h)^3)$. Group the nonzero $h$ according
to the ideal $\mathfrak d=\mathfrak a+(h)^3$.
There are $O_{K,\varepsilon}(V^\varepsilon)$ possible divisor ideals.
For each fixed class put
\[
 \mathfrak b=\prod_{\mathfrak p}
       \mathfrak p^{\lceil v_\mathfrak p(\mathfrak d)/3\rceil},
 \qquad b=N\mathfrak b.
\]
Then $h\in\mathfrak b$ and $N\mathfrak d\le b^3$.
Since $|h_v|\ll_KH^{1/n}$ at every place, one has
\[
 A_W(h^3\beta)\gg_K
       \max\{1,A|Nh|^3/H^3\}.
\]
Consequently the reciprocal-height sum in this class is at most
\[
 \frac{b^3}{Q}
 \sum_{0\ne h\in\mathfrak b\cap\mathcal B(H)}
       \min\{1,C_KH^3/(A|Nh|^3)\}.
\]
Put $T=C_KH/A^{1/3}$, enlarging the fixed constant if necessary.
Lemma~\ref{lem:g789-smallnorm} and dyadic summation bound the inner
sum by
\[
 O_{K,\varepsilon}\left(H^\varepsilon
 \begin{cases}T/b,&T\ge b,\\(T/b)^3,&T<b.\end{cases}\right).
\]
In the first case $b^2T\le T^3$, and in the second case the factor
$b^3$ cancels. Thus either case contributes at most
$O_{K,\varepsilon}(V^\varepsilon H^3/(QA))$.
Summing the divisor classes and recalling $QA=L$ proves
\[
 \sum_{h:\,\mathcal H_W(\alpha h^3)<Y}
       \mathcal H_W(\alpha h^3)^{-1}
 \ll_{K,\varepsilon}V^\varepsilon H^3/L.
\]
Multiplication by $JN\sum|c(a)|^2$ completes the proof.
\end{proof}

\subsection{A second coefficient moment}
\label{sec:k5-coeff2}

All size parameters in this section denote volumes.  The field is fixed,
$n=[K:\mathbb Q]$, and $\mathcal H_T$ denotes the conductor--product height
defined in Section~\ref{sec:preliminaries}.  The notation $\mathcal B(L)$ means a fixed
archimedean box with side lengths comparable to $L^{1/n}$.  The same estimates hold on fixed fractional lattices.

The following is the precise number-field form of the coefficient estimate
needed from Vaughan--Wooley's fifth-power argument \cite[Lemma~4.3]{VW1995}.  It is useful to state
it for the general coefficient degree $2d$.

\begin{proposition}\label{prop:k5-coeff2}
Let $d\ge3$, $J,H,M\ge1$, and $T=JH^3M^{2d}$.  Suppose
\[
 1\le U\le c_K\min\{M,JH^{-3},T^{1/(4d)},
                  (JH^3M^{1-2d})^{1/6}\}.
\]
Let $\mathcal C$ be any subset of a sufficiently small fixed box about a
nonzero point at every embedding, of scale $M^{1/n}$.  For each $h$ one may
include any fixed lower-degree additive twist in the $m$ sum, provided it
is independent of $j$.  Then
\[
 E_1(\alpha):=\sum_{\substack{h\in\mathcal B(H)\\h\ne0}}
 \sum_{j\in\mathcal B(J)}
 \left|\sum_{m\in\mathcal C}e_K(\alpha jh^3m^{2d}+b_h(m))\right|^2
 \ll_{K,d,\epsilon}T^\epsilon JHM^2
       \{\mathcal H_T(\alpha)^{-1/d}+U^{-1}\}.
\]
The same bound holds for smooth weights and for the fixed lattices arising
from class representatives.  No positivity condition at complex places is
being imposed.
\end{proposition}

\paragraph{Conventions for the proof.}
A nonzero algebraic integer has absolute norm at least one.  After clearing
a fixed denominator this also applies, with a fixed positive lower bound,
to every lattice occurring below.  An arbitrary conductor ideal can be
made principal by multiplying it by one of finitely many fixed class
representatives.  Its chosen denominator therefore has norm comparable
to the conductor norm.  In all congruence counts it is the conductor ideal,
not that auxiliary denominator, which is used.

We use anisotropic Dirichlet approximation in the following elementary
form.  Given positive radii $R_v$ with $\prod_vR_v^{d_v}\gg_K1$, Minkowski's
theorem supplies $a\in\mathfrak D^{-1}$ and nonzero $q\in\mathcal O_K$ with
\[
 |q_v|\ll_KR_v,\qquad |(q\beta-a)_v|\ll_K R_v^{-1}.
\]
The constants can be made harmless by reducing the fixed constant $c_K$
in the proposition.  All later identities follow by estimating the norm
of a cross product and showing that it is smaller than its nonzero lower
bound.  This replaces the rational proof's use of a nonzero integer of
absolute value less than one.

A smooth positive majorant of a linear sum over $\mathcal B(J)$ is a
periodic rapidly decreasing kernel $D_J$.  Dyadic boxes for its errors,
and dyadic boxes for all nonzero variables, cost only $T^\epsilon$.
Thus the estimates below may be proved with specified radii at every
embedding.  Products of these radii, with local dimensions $d_v$, give
the volume expressions displayed below.  For example the contribution
where the error kernel has normalized product size at least $U$ is
$O(T^\epsilon JHM^2/U)$.

\begin{proof}
Expanding the square and summing over $j$, using a smooth majorant before
taking absolute values, reduces the assertion to the untwisted positive
kernel sum.  Terms with $m_1=m_2$ contribute $O(JHM)$ and are acceptable.
Because the $m$ box is sufficiently small, $m_1^{2d}=m_2^{2d}$ within this
box implies $m_1=m_2$.  Put $D_0=m_2^{2d}-m_1^{2d}$ for the other pairs.

\emph{First identification: the cubic shift.}
For a fixed pair and an error box of normalized product size at most $U$,
approximate $\alpha D_0$ with denominator radii
$(J_vH_v^3/U_v)^{1/2}$, where $\prod U_v^{d_v}=U$.
If $\alpha D_0h^3$ is in that error box around $a_0\in\mathfrak D^{-1}$,
then the norm of the cross product between this approximation and $a_0/h^3$
is
\[
 \ll_K (UH^3/J)^{1/2}.
\]
It is therefore zero.  Write the resulting common center as $\gamma_{D_0}$.
Its conductor $\mathfrak q_{D_0}$ divides $(h)^3$.
The sum over such $h$, using ideal-lattice packing on dyadic boxes, is at most
\[
 T^\epsilon JH V_{D_0}^{-1/3},\qquad
 V_{D_0}=N\mathfrak q_{D_0}\,
 A_{JH^3}(\alpha D_0-\gamma_{D_0}).
\]
Here is a lattice proof of the summation, including thin boxes.  If
$\mathfrak b=\prod\mathfrak p^{\lceil v_\mathfrak p(\mathfrak q)/3\rceil}$,
then $h\in\mathfrak b$ and $N\mathfrak b\ge(N\mathfrak q)^{1/3}$.
For $A=A_{JH^3}(\beta)$ one has
$A_J(h^3\beta)\gg\max(1,A|Nh|^3/H^3)$.
The elementary ideal count
$\#\{h\in\mathfrak b\cap\mathcal B(H):0<|Nh|\le B\}
\ll H^\epsilon B/N\mathfrak b$, with no points for $B<N\mathfrak b$,
and dyadic summation give
$\sum_h A_J(h^3\beta)^{-1}\ll H^\epsilon H/(A^{1/3}N\mathfrak b)$.
This proves the displayed bound without a boundary or real-positivity assumption.
We may discard the pairs for which $V_{D_0}\ge U^3$ at the required cost.

\emph{Second identification: remove the divided difference.}
Group pairs according to their greatest common divisor ideal
$\mathfrak g=(m_1,m_2)$, of norm $G$.  Make $\mathfrak g$ principal using a
fixed class representative and choose a balanced generator $g$.
Write $m_1=gn$, $m_2=g(n+l)$.  The variables $n,n+l$ belong to fixed fractional
lattices, are primitive up to fixed primes, and have scale $N=M/G$.
The polynomial
\[
 D(n,l)=\frac{(n+l)^{2d}-n^{2d}}{l}
\]
has absolute value comparable to $N_v^{2d-1}$ at every embedding.
Indeed, both arguments are in the same tiny box away from zero, so the
leading derivative $2d n^{2d-1}$ dominates the other terms.
This is the place where localization replaces real positivity.

For fixed $g,l$, approximate $\alpha g^{2d}l$ with denominator radii
$(J_vH_v^3N_v^{2d-1})^{1/2}$, obtaining a center $\gamma_{g,l}$.
A comparison with $\gamma_{D_0}/D(n,l)$ has cross-product norm at most
\[
 \ll_K V_{D_0}
       \left(\frac{N^{2d-1}}{JH^3}\right)^{1/2}
 \le C_K U^3\left(\frac{M^{2d-1}}{JH^3}\right)^{1/2}<1.
\]
Consequently $\gamma_{D_0}=D(n,l)\gamma_{g,l}$.
If $\mathfrak q$ is the conductor of $\gamma_{g,l}$, then
$\mathfrak q_{D_0}=\mathfrak q/\mathfrak q_1$, with
$\mathfrak q_1\mid(D(n,l))$.
Moreover $n(n+l)$ is prime to $\mathfrak q_1$ outside a fixed set of primes.
The number of primitive solutions of $D(n,l)\equiv0\pmod{\mathfrak q_1}$
in residue classes is $O_{K,d,\epsilon}((N\mathfrak q_1)^{\epsilon})$.
To see the uniformity even at primes dividing $2d$, let
$t=v_{\mathfrak p}(l)$.  For a primitive solution modulo $\mathfrak p^E$,
both $n$ and $n+l$ are units.  The map
$n\mapsto r=1+l/n\pmod{\mathfrak p^{E+t}}$ is injective, and
$D(n,l)\equiv0\pmod{\mathfrak p^E}$ implies
$r^{2d}\equiv1\pmod{\mathfrak p^{E+t}}$.
The number of unit $2d$th roots modulo any power of a fixed prime is
bounded in terms of $K,d$: beyond the fixed derivative depth
$v_{\mathfrak p}(2d)$, Hensel lifting gives only a bounded number of
lifts around each root.  At good primes the bound is $2d$.
Multiplying the local bounds gives the stated $T^\epsilon$ factor.
Fixed denominators from class representatives are cleared first; their
fixed valuations only change the constant.

Put
\[
 V_{g,l}=N\mathfrak q\,
 A_{JH^3N^{2d-1}}(\alpha g^{2d}l-\gamma_{g,l}).
\]
The preceding comparison of $D(n,l)$ with its archimedean scale gives
$V_{D_0}\asymp_K V_{g,l}/N\mathfrak q_1$.
Ideal packing gives at most
$T^\epsilon(1+N/N\mathfrak q_1)$ choices for $n$ in each possible class.
The terms $N\mathfrak q_1>N$ contribute $O(T^\epsilon JHM)$ in total:
there are $O(T^\epsilon)$ possible $n$ for each $g,l$ and divisor class,
and their weight is $V_{D_0}^{-1/3}\le1$;
in each dyadic $G$ range the number of $g,l$ pairs is $O(T^\epsilon M)$.
For the rest, $(N\mathfrak q_1)^{1/3}/N\mathfrak q_1\le1$.
Thus the remaining sum is bounded by
\[
 T^\epsilon JHM\sum_gG^{-1}\sum_l V_{g,l}^{-1/3}.
\]
Here and below the summation over $g$ means one balanced generator per
ideal and finitely many fixed class choices; its count in a dyadic range
of $G$ is $O_K(GT^\epsilon)$.
We may further restrict to $V_{g,l}<(U/G)^3$, at a cost
$O(T^\epsilon JHM^2/U)$.

\emph{Third identification: remove the difference.}
For fixed $g$, approximate $\alpha g^{2d}$ with denominator radii
$(J_vH_v^3N_v^{2d})^{1/2}$, obtaining $\gamma_g$.
Comparison with $\gamma_{g,l}/l$ has cross-product norm
\[
 \ll_K\frac{V_{g,l}}{(JH^3)^{1/2}N^{d-1}}<1.
\]
The previous bound on $V_{g,l}$ and the hypotheses on $U$ ensure the last
inequality.  Hence $\gamma_{g,l}=l\gamma_g$ exactly.
Splitting according to the divisor ideal of the conductor that divides $l$,
and summing $l$ on dyadic archimedean boxes, gives
\[
 \sum_l V_{g,l}^{-1/3}\ll T^\epsilon N V_g^{-1/3},\qquad
 V_g=Q_{\gamma_g}A_{JH^3N^{2d}}(\alpha g^{2d}-\gamma_g).
\]
For an explicit verification, put $A=A_{JH^3N^{2d}}(\beta)$.
Then $A_{JH^3N^{2d-1}}(l\beta)\gg A|Nl|/N$.
If $l$ belongs to a divisor ideal $\mathfrak b$, the preceding ideal norm
count gives
$\sum_{0\ne l\in\mathfrak b\cap\mathcal B(N)}|Nl|^{-1/3}
\ll T^\epsilon N^{2/3}/N\mathfrak b$.
Combining this with the conductor factor yields divisor weights
$(N\mathfrak b)^{-2/3}$, whose sum is $O(T^\epsilon)$.
This proves the displayed linear-variable estimate.
We have therefore reduced to
\[
 T^\epsilon JHM^2\sum_gG^{-2}V_g^{-1/3}.
\]
Terms with $V_g\ge(U/G)^3$ are again acceptable, so assume the reverse
inequality.  In particular $G<U$.

\emph{Fourth identification: remove the common divisor.}
Choose one Dirichlet center $\gamma$ for $\alpha$ at scale $T$.
Comparison with $\gamma_g/g^{2d}$ has cross-product norm at most
\[
 \ll_K V_gG^{2d}T^{-1/2}
 \le C_KU^3G^{2d-3}T^{-1/2}
 \le C_KU^{2d}T^{-1/2}<1.
\]
Thus $\gamma_g=g^{2d}\gamma$.  Write $Q=Q_\gamma$ and $A=A_T(\alpha-\gamma)$.
The conductor of $g^{2d}\gamma$ is obtained by dividing the conductor of
$\gamma$ by its greatest common divisor with $(g)^{2d}$.
The chosen generators are balanced, so the archimedean factor in $V_g$
is comparable to $A$.
Finally
\[
 \sum_{\mathfrak g}(N\mathfrak g)^{-2}
 \left(\frac{Q}{N(\mathfrak q+(\mathfrak g)^{2d})}\right)^{-1/3}
 \ll_{K,d,\epsilon}T^\epsilon Q^{-1/d}.
\]
This is an elementary Euler-product calculation: at a prime with
$v_\mathfrak p(\mathfrak q)=e$, the term of valuation $a$ has size
$q^{-2a-(e-\min(e,2da))/3}$, and its largest possible size is
$O(q^{-e/d})$ because $d\ge3$.  Summing the geometric tails and absorbing
the bounded factors per prime into $T^\epsilon$ proves the claim.
Since $A^{-1/3}\le A^{-1/d}$,
\[
 \sum_gG^{-2}V_g^{-1/3}\ll T^\epsilon(QA)^{-1/d}
 \le T^\epsilon\mathcal H_T(\alpha)^{-1/d}.
\]
Combining the four reductions proves the proposition.
\end{proof}


\subsection{Sampled quadratic moments}
\label{sec:quad-sampling-field}

Let $K$ be fixed, $n=[K:\mathbb Q]$, $d_v=[K_v:\mathbb R]$,
$V=K\otimes_{\mathbb Q}\mathbb R$, and
$\mathbb T=V/\mathfrak D^{-1}$. All size parameters in this section
are volumes. Write $\mathcal B(U)$ for a bounded archimedean box with
all radii $O(U^{1/n})$. Fixed dilations and translates are harmless.
For a vector $\mathbf T=(T_v)_{v\mid\infty}$ of positive real numbers put
\[
 A_{\mathbf T}(\beta)=\prod_v\max\{1,T_v|\beta_v|\}^{d_v},\qquad
 \mathcal H_{\mathbf T}(\alpha)=
       \min_{\gamma\in K} Q_\gamma A_{\mathbf T}(\alpha-\gamma),
 \qquad
 Q_\gamma=N\{a\in\mathcal O_K:a\gamma\in\mathfrak D^{-1}\}.
\]
When every $T_v=T^{1/n}$ we abbreviate the height to $\mathcal H_T$.
The minimum includes all lifts of a center. Product Dirichlet approximation
gives $\mathcal H_T(\alpha)\ll_K T^{1/2}$.

\begin{lemma}[Anisotropic bilinear and cluster estimates]
\label{lem:anisotropic-bilinear}
Let $\mathcal U,\mathcal W$ be archimedean boxes of radii $U_v,W_v$,
with volumes $J=\prod_vU_v^{d_v}\ge1$ and
$N=\prod_vW_v^{d_v}\ge1$. Put $T_v=U_vW_v$.
For any coefficients supported on $\mathcal W\cap\mathcal O_K$,
\[
 \sum_{z\in\mathcal U\cap\mathcal O_K}
  \left|\sum_m c_m e_K(\alpha zm)\right|^2
 \ll_K \left(J+N+\frac{JN}{\mathcal H_{\mathbf T}(\alpha)}\right)
               \sum_m|c_m|^2.
\]
Consequently the number of $z$ in $\mathcal U$ for which $\alpha z$
modulo $\mathfrak D^{-1}$ lies in a prescribed box of radii $O(W_v^{-1})$
is
\[
 \ll_K 1+\frac JN+\frac J{\mathcal H_{\mathbf T}(\alpha)}.
\]
The constants are independent of the eccentricities of the two boxes.
\end{lemma}
\begin{proof}
The proof of Lemma~\ref{lem:g789-bilinear} works
with these radii without alteration. We give the relevant uniformity details.
For $J\le N$, subdivide $\mathcal U$ by a fixed number of subdivisions
at each embedding. The product of the sample-box volume and a
frequency-error box volume is then sufficiently small for the affine
field-line argument: the determinant of any two difference pairs has
norm smaller than the fixed lower bound for a nonzero element of
$\mathfrak D^{-1}$. A cluster therefore lies on an affine $K$-line of
slope $\gamma$. Ideal-lattice packing bounds its size by
\[
 1+Q_\gamma^{-1}
       \prod_v\min\{U_v,W_v^{-1}/|\alpha_v-\gamma_v|\}^{d_v}
 =1+\frac J{Q_\gamma A_{\mathbf T}(\alpha-\gamma)}.
\]
The Gaussian large sieve then gives the stated bound, with
$N+JN/\mathcal H_{\mathbf T}(\alpha)$ on the right. When $J>N$,
transpose the bilinear matrix, interchanging the two boxes. Its operator
norm is unchanged and the product radii $T_v$ are unchanged. This proves
the displayed symmetric bound.

For the cluster assertion, put coefficient weight
$\exp(-c\sum_v|m_v/W_v|^2)$ on the dual box. Poisson summation gives a
nonnegative Gram kernel which is $\gg_K N$ for any pair of the selected
sample frequencies. If there are $M$ samples, the corresponding quadratic
form is therefore $\gg_K NM^2$. Its upper bound from the bilinear operator
estimate is $O_K((J+N+JN/\mathcal H_{\mathbf T})M)$.
Gaussian tails are summed by dyadic dilation and rapid decay. Division
gives the claim. The only lattice input throughout is
$\#(\mathfrak a\cap\mathcal R)\ll_K1+\operatorname{vol}(\mathcal R)/N\mathfrak a$
for an arbitrary product box, so no bounded-eccentricity hypothesis is used.
\end{proof}

\begin{proposition}[Distribution of short heights of multiples]
\label{prop:short-height-multiples}
Let $X,H\ge1$, put $V_0=HX^2$, and let
$L=\mathcal H_{V_0}(\alpha)$. For $1\le Y\ll_K X$ one has
\[
 \#\{0\ne h\in\mathcal B(H):\mathcal H_{X^2}(h\alpha)\le Y\}
 \ll_{K,\epsilon} V_0^\epsilon
       \left(\frac{HY^2}{X^2}+\frac{HY}{L}\right).
\]
\end{proposition}
\begin{proof}
Choose an approximation $h\alpha=\gamma+\beta$ of height at most $Y$.
Let $\mathfrak q$ be its conductor and write $Q=N\mathfrak q$.
Minkowski's theorem supplies $0\ne b\in\mathfrak q$ with
$|b_v|\ll_KQ^{1/n}$ at every embedding. Since $(b)\subset\mathfrak q$,
$|Nb|\ge Q$; consequently also $|b_v|\gg_KQ^{1/n}$ at every embedding.
Thus $|Nb|\asymp_K Q$ and $b\gamma\in\mathfrak D^{-1}$.
Put $m=bh$. It lies in a box of volume $O_K(HQ)$ and satisfies
\[
 \alpha m\equiv b\beta\pmod{\mathfrak D^{-1}}.
\]
Each nonzero $m$ has only $V_0^\epsilon$ possible bounded-height factors
$b$, by the ideal-divisor bound and the unit-lattice count.

Decompose $Q$ and
$\max(1,X^{2/n}|\beta_v|)$ into dyadic intervals. Denote the latter
upper endpoints by $A_v\ge1$ and their product by
$A=\prod_v A_v^{d_v}\ll_KY/Q$.
There are only logarithmically many such choices. For a fixed choice,
the residue error $b\beta$ is contained in a box of radii
$O_K(Q^{1/n}X^{-2/n}A_v)$. Apply
Lemma~\ref{lem:anisotropic-bilinear} with
\[
 U_v=(HQ)^{1/n},\qquad
 W_v=(X^2/Q)^{1/n}A_v^{-1},\qquad
 J=HQ,\quad N=\frac{X^2}{QA}.
\]
Here $N\gg_KX\ge1$, after harmless fixed rescaling, and
$T_v=V_0^{1/n}/A_v$.
For every approximation center,
$A_{V_0}(\alpha-\gamma)\le A\,A_{\mathbf T}(\alpha-\gamma)$.
Taking minima gives $L\le A\mathcal H_{\mathbf T}(\alpha)$.
The number of such $m$ is therefore
\[
 \ll_K1+\frac{HQ^2A}{X^2}+\frac{HQA}{L}
 \ll_K1+\frac{HY^2}{X^2}+\frac{HY}{L}.
\]
Finally, the approximation to $h\alpha$ lifts to one to $\alpha$
with center $\gamma/h$. Its conductor has norm at most $Q|Nh|$, and
\[
 |Nh|A_{V_0}(\beta/h)
   =\prod_v\max\{|h_v|,V_0^{1/n}|\beta_v|\}^{d_v}
   \ll_K H A_{X^2}(\beta).
\]
Thus the set being counted is empty unless $L\ll_K HY$.
On its nonempty range this absorbs the constant $1$ above.
Summing the dyadic choices and factor multiplicities proves the result.
\end{proof}

\begin{corollary}[Sampled fourth moment]
\label{cor:quadratic-fourth-field}
Let $Z_t(\beta)$ be a quadratic exponential sum on a box of root volume
$X$, with arbitrary linear modulation $t$. For arbitrary modulations $t_h$,
\[
 \sum_{0\ne h\in\mathcal B(H)}|Z_{t_h}(h\alpha)|^4
 \ll_{K,\epsilon}(HX^2)^\epsilon
 \left\{HX^2+X^4\min\left(\frac HL,\frac{H^2}{L^2}\right)\right\},
 \qquad L=\mathcal H_{HX^2}(\alpha).
\]
The assertion is uniform for fixed nonzero integral multipliers of $\alpha$.
\end{corollary}
\begin{proof}
One ordinary difference removes the linear modulation. The positive
linear-kernel estimate, proved by the bilinear Gram argument above, gives
\[
 |Z_t(\beta)|^2\ll X^\epsilon
       \left(X+\frac{X^2}{B}\right),\qquad
 B=\mathcal H_{X^2}(\beta)\ll_K X.
\]
Sharp fixed boxes can be completed with a logarithmic loss; smooth majorants
give the same estimate directly. In particular the bound is uniform in $t$.
Hence $|Z_t(\beta)|^4\ll X^\epsilon(X^2+X^4B^{-2})$.
Proposition~\ref{prop:short-height-multiples}, summed dyadically from
$Y\asymp_K\max(1,L/H)$ to $O_K(X)$, shows that
\[
 \sum_{0\ne h\in\mathcal B(H)}\mathcal H_{X^2}(h\alpha)^{-2}
 \ll (HX^2)^\epsilon
       \left\{\frac H{X^2}
        +\min\left(\frac HL,\frac{H^2}{L^2}\right)\right\}.
\]
Combining the last two bounds proves the assertion. A fixed multiplier
changes conductor heights only by fixed factors.
\end{proof}

In the eighth-power terminal, $H=X^h$ and $L\asymp X^\ell$.
Normalization by $HX^4$ gives
\[
 z_4(\ell)=\min\{2,\max(\ell,2\ell-h)\}.
\]
Both resources used by the exact policy,
$\min\{2,2\ell-h\}$ and $\min\{2,\ell,h\}$, are weaker and therefore valid.
Products of a fixed number of shifts may replace $h$, at the cost of the
same divisor loss. The quadratic second-moment resource remains the
previously proved bound with normalized saving $1$ when $H\ge X$ and
$L\gg X^{1-O(\eta)}$.

\section{Restricted-height smooth means}
\label{sec:restricted-field}

The following field version of Br\"udern--Wooley's restricted-height
mean theorem is needed by the eighth- and ninth-power routes.  The use of
ideals in its proof is particularly useful: an extracted ideal factor
need not be principal.  The quotient variables instead lie in one of
finitely many fixed integral ideals.

Write $V=K\otimes\mathbb R$, $n=[K:\mathbb Q]$, and $d_v=[K_v:\mathbb R]$.
For $T\ge1$ put
\[
 A_T(\beta)=\prod_{v\mid\infty}
                 \max(1,T^{1/n}|\beta_v|)^{d_v}.
\]
For a fixed integral ideal $C$, use the coefficient torus
$\mathbb T_C=V/(C^{-k}\mathfrak D^{-1})$ with probability Haar measure.
The conductor of $\gamma\in K/(C^{-k}\mathfrak D^{-1})$ is
\[
 \mathfrak a_C(\gamma)
       =\{a\in\mathcal O_K:a\gamma C^k\subseteq\mathfrak D^{-1}\}.
\]
Let $\mathfrak M_C(Z;Y)$ be the union of the product arcs
\[
 (N\mathfrak a_C(\gamma)) A_{Y^k}(\alpha-\gamma)\le Z.
\]
All lifts of a center are included.  A subscript $\mathfrak a$ on this
set restricts to centers of exact conductor $\mathfrak a$.
For $C=\mathcal O_K$, suppress the subscript $C$.

Let $g(\alpha;X,R)$ be a sum over the nonzero algebraic integers in a fixed
bounded box of scale $X^{1/n}$ whose prime ideal factors have norm at most
$R$.  Nonnegative bounded weights and smaller fixed boxes are allowed.

\begin{proposition}[Field restricted-height mean]
\label{prop:restricted-field}
Let $s\ge2$ be even.  Suppose $\lambda=s-k+\Delta$ is a permissible
full smooth-moment exponent, with
$\lambda\ge\max(s/2,s-k)$ and $\Delta\ge0$.  Then for every
$\epsilon>0$ a sufficiently small smoothness exponent gives
\[
 \int_{\mathfrak M(Z;X)}|g(\alpha;X,R)|^s\,d\alpha
       \ll_{K,k,s,\epsilon} X^{s-k+\epsilon}Z^{2\Delta/k}
       \qquad(1\le Z\le X^{k/2}).
\]
The same conclusion holds at intermediate real orders when their full
exponents are obtained by H\"older interpolation between such even orders.
\end{proposition}

\subsection{Fixed root lattices and multiplication of arcs}

Choose integral ideal representatives $C_1,\ldots,C_h$ for the ideal
classes.  For every nonzero integral ideal $\mathfrak b$ in the class of
$C_i$, choose $b=b(\mathfrak b)\in K^\times$ with
\[
 \mathfrak b=b C_i,\qquad
 |b_v|\asymp_K(N\mathfrak b)^{1/n}\quad(v\mid\infty).
\]
Dirichlet's unit theorem gives the second condition after choosing one
representative in a fundamental domain for the unit logarithms.  This
choice is made once for each ideal.  Multiplication by $b$ maps $C_i$
into $\mathcal O_K$.

If $(\mathfrak a,\mathfrak b)=1$, multiplication by $b^k$ maps centers
of exact conductor $\mathfrak a$ on $\mathbb T_{\mathcal O_K}$ bijectively
to those of exact conductor $\mathfrak a$ on $\mathbb T_{C_i}$.
Indeed the conductor calculation is local, and
$b^kC_i^k=\mathfrak b^k$ is a unit ideal at every prime dividing
$\mathfrak a$.  Equivalently, multiplication by $b^k$ is an isomorphism
on these finite conductor-character groups.
Its archimedean Jacobian, with the indicated probability measures, is
$(N\mathfrak b)^k$.  To check the normalization, ordinary Lebesgue measure
has Jacobian $|Nb|^k=(N\mathfrak b/NC_i)^k$, while the volume of
$\mathbb T_{C_i}$ is $(NC_i)^{-k}$ times that of $\mathbb T_{\mathcal O_K}$.

Consequently, whenever the image arcs are disjoint, for any nonnegative
integrand $F$ one has
\begin{equation}\label{eq:field-arc-scaling}
 \int_{\mathfrak M_{\mathfrak a}(Z;X)}F(b^k\alpha)\,d\alpha
 \ll_K (N\mathfrak b)^{-k}
 \int_{\mathfrak M_{C_i,\mathfrak a}(c_1Z;X/N\mathfrak b)}
                         F(\theta)\,d\theta .
\end{equation}
Here $c_1$ is a fixed constant, independent of the ideal and its class.
Balancing $b$ makes the archimedean scaling comparable at every place.
The requisite disjointness holds if
$c_1Z\le c_2Y^{k/2}$, with a sufficiently small $c_2=c_2(K,k)$,
where $Y=X/N\mathfrak b$.
For distinct lifted centers the nonzero norm of their difference is
$\gg_K(QQ')^{-1}$, while overlapping product arcs would give an upper
bound $\ll_KY^{-k}Z^2/(QQ')$.  This also excludes overlap of distinct
lifts of the same center.  There are only finitely many root lattices
in use, so one constant suffices for all of them.

\subsection{The ideal-factor decomposition}

Fix a total ordering of the prime ideals which refines their norm ordering.
Given a conductor ideal $\mathfrak a$, factor
\[
 (x)=\mathfrak v\mathfrak u,
 \qquad (\mathfrak v,\mathfrak a)=1,
 \qquad \mathfrak u\text{ supported on primes dividing }\mathfrak a.
\]
Call a root bad when $N\mathfrak v\le M$.  The supported-ideal counting
lemma from the fundamental differencing argument, followed by bounded-height
unit counting, gives at most $X^\epsilon M$ such roots, uniformly for
$N\mathfrak a\le X^{k/2}$.  This remains true when the variable lies in
a smaller box or has a nonnegative weight bounded by a fixed constant.

For a good root, multiply the prime factors of $\mathfrak v$ in decreasing
order until their product first has norm exceeding $M$.  Denote the
resulting ideal by $\mathfrak m$ and its last prime factor by $\mathfrak p$.
This decomposition is unique, and
\[
 M<N\mathfrak m\le MR,\qquad
 (\mathfrak m,\mathfrak a)=1.
\]
Every prime of the remaining ideal $\mathfrak w=\mathfrak v/\mathfrak m$
is at most $\mathfrak p$ in the fixed ordering.  The condition on
$\mathfrak m$ is exactly that its factors are at least $\mathfrak p$ and
$N(\mathfrak m/\mathfrak p)\le M<N\mathfrak m$.
Thus for each outer $\mathfrak m$, the remaining $\mathfrak w$ runs over
the complete set of $\mathfrak p$-smooth ideals prime to $\mathfrak a$.

Write $\mathfrak m=m C_i$ using the balanced choice above.  Then $x=mr$
with $r\in C_i$ and
$(r)C_i^{-1}=\mathfrak w\mathfrak u$.
Split $\mathfrak u$ into its prime factors at most $\mathfrak p$ and those
larger than $\mathfrak p$; call the latter product $\mathfrak z$.
The smaller part combines with $\mathfrak w$ to give every
$\mathfrak p$-smooth integral ideal, without a coprimality restriction.
Writing $\mathfrak z=z C_j$, one obtains $r=zy$ with
\[
 y\in C_iC_j,
 \qquad (y)(C_iC_j)^{-1}\text{ a }\mathfrak p\text{-smooth integral ideal}.
\]
The original box condition is simply its original weight evaluated at
$mzy$.  For fixed $\mathfrak m,\mathfrak z$ this inner sum is independent
of the conductor ideal.  All root lattices $C_i$ and $C_iC_j$ are fixed;
there are only finitely many of them.  This is the exact replacement for
the rational factor decomposition in \cite[Sections~3--4]{BW2023}.

\subsection{Integration and the exponent}

Choose a sufficiently small fixed $c>0$ and put
\[
                  M=cX Z^{-2/k}R^{-1}.
\]
If $M<R$, the desired assertion follows directly from the full moment:
$Z\gg_K X^{k/2}R^{-k}$, so the loss is at most $R^{2\Delta}$,
absorbed by choosing the smoothness exponent small enough.
We may therefore assume $M\ge R$.
There are $O_K(MR)$ choices of the ideal $\mathfrak m$.
H\"older's inequality applied to the preceding exact decomposition gives
a factor $(MR)^{s-1}$ and a sum of $s$th moments of the inner sums.
The bad part contributes at most
\begin{equation}\label{eq:restricted-bad}
                       X^{-k+\epsilon}Z^2M^s.
\end{equation}
Here the total measure of the product arcs is
$\ll_KX^{-k}Z^2(1+\log X)^{O_K(1)}$: a conductor of norm $q$ has at
most $q$ centers, each error region has volume
$\ll_K X^{-k}(Z/q)(1+\log Z)^{O_K(1)}$, and there are $O_K(Z)$ ideals
of norm at most $Z$.

For each fixed $\mathfrak m$, apply \eqref{eq:field-arc-scaling}.
Our choice of $c$ ensures disjointness on the quotient root scale because
$N\mathfrak m\le MR=cXZ^{-2/k}$.
There are only $X^\epsilon$ ideals $\mathfrak z$ supported on a given
conductor and of the relevant norm.  Another application of H\"older
therefore costs $X^\epsilon$.
Now extend the $\mathfrak z$ summation to all integral ideals in its
allowed size range.  For fixed $\mathfrak m,\mathfrak z$, the inner
integrand no longer depends on $\mathfrak a$.  The sum over the disjoint
conductor arcs is consequently at most a full torus integral.
Multiplication by $z^k$ maps $\mathbb T_{C_i}$ onto
$\mathbb T_{C_iC_j}$ and preserves probability Haar integrals.
The remaining full moment is bounded by
\[
 X^\epsilon\left(\frac{X}{N\mathfrak m N\mathfrak z}\right)^\lambda.
\]
Indeed its weights are nonnegative and bounded, its support lies in a
fixed box of that volume scale, and its roots lie in a fixed integral
ideal.  At even order, orthogonality bounds it by a constant times the
complete smooth moment on $\mathcal O_K$, after allowing the finitely
many fixed prime factors of $C_iC_j$.
This use of coefficient domination is only on a full torus integral.

The permissible bound is uniform for these smaller scales: for scales
$Y\ge X^\delta$, choose the original smoothness exponent at most
$\delta$ times a permissible one; for $Y<X^\delta$, use the trivial
bound and choose $\delta$ sufficiently small in terms of the allocated
$\epsilon$.  Bounded scales and the finitely many fixed primes cause
only a constant or $X^\epsilon$ factor.

Since $\lambda\ge s/2\ge1$, ideal counting gives
$\sum_{N\mathfrak z\le X}(N\mathfrak z)^{-\lambda}\ll_K X^\epsilon$,
including the logarithmic endpoint $\lambda=1$.
The good contribution is therefore at most
\[
 X^\epsilon(MR)^{s-1}
 \sum_{M<N\mathfrak m\le MR}(N\mathfrak m)^{-k}
             \left(\frac{X}{N\mathfrak m}\right)^\lambda
 \ll X^{s-k+\epsilon}(X/M)^\Delta.
\]
Inserting $M$ gives the asserted main bound.
The bad contribution \eqref{eq:restricted-bad} is
$\ll X^{s-k+\epsilon}Z^{2-2s/k}$, which is no larger because
$\Delta\ge k-s$ and $Z\ge1$.
H\"older interpolation proves the final assertion at real orders.
This completes the proof of Proposition~\ref{prop:restricted-field}.

\section{The cubic terminal}
\label{sec:cubic-terminal}
\label{sec:g7-cubic}
All parameters are volumes, and the conductor--product height is the one fixed in Section~\ref{sec:preliminaries}.
\subsection{A linear-kernel consequence of the bilinear large sieve}

The elementary bilinear sampling bound is
\begin{equation}\label{eq:g7-bilinear}
 \sum_{z\in\mathcal B(J)}\left|\sum_{a\in\mathcal B(N)}c(a)e_K(\theta za)\right|^2
 \ll_K\left(J+N+\frac{JN}{\mathcal H_{JN}(\theta)}\right)\sum_a|c(a)|^2.
\end{equation}
This is Lemma~\ref{lem:g789-bilinear}. The same bound holds for fixed translates, subsets, and fixed dilations of the boxes.

Here is the corresponding positive kernel estimate.  Let
\[
 D_N(t)=\sum_{z\in\mathcal O_K}
      \exp(-c|z|^2/N^{2/n})e_K(tz).
\]
Choose the Euclidean norm compatibly with the trace pairing.  Poisson
summation writes this as a positive sum of Gaussians centered on
$\mathfrak D^{-1}$, each of height $O_K(N)$ and width $N^{-1/n}$.
Then
\begin{equation}\label{eq:g7-kernel}
 \sum_{h\in\mathcal B(J)}D_N(\theta h)
 \ll_K J+N+\frac{JN}{\mathcal H_{JN}(\theta)}.
\end{equation}
Indeed, enlarge the $h$ box by a fixed factor.  Every difference in the
original box has $\gg_K J$ representations as $r-s$ in the enlarged box.
Because $D_N\ge0$, its sum over the original box is at most a constant times
$J^{-1}\sum_{r,s}D_N(\theta(r-s))$.  The latter double sum is a quadratic
form in the Gram matrix of $(e_K(\theta rz))_{r,z}$.  Its norm is bounded by
\eqref{eq:g7-bilinear}.  Gaussian tails in $z$ are handled by dyadic boxes;
monotonicity of $\mathcal H_T$ in $T$ and Gaussian decay make their total
convergent.  This proves \eqref{eq:g7-kernel}.
For Fourier transforms of smooth compactly supported weights, replace the
Gaussian kernel by a rapidly decreasing positive majorant; a dyadic
superposition of Gaussian kernels proves the same bound.  When the kernel
parameter decreases from $N$ to $N/2^r$, use that
$\mathcal H_t(\theta)/t$ is nonincreasing in $t$, directly from its
definition; consequently $(JN/2^r)/\mathcal H_{JN/2^r}(\theta)$
is at most $JN/\mathcal H_{JN}(\theta)$.  The rapid tail weights
absorb the corresponding powers of $2^r$.

\subsection{The cubic family}

Let $j=k-3$, and put
\[
 M_i=X^{\phi_i},\quad H_i=XM_i^{-k},\quad
 M=\prod_{i\le j}M_i,\quad H=\prod_{i\le j}H_i,\quad
 Q=X/M,\quad T=Q^k=X^3H.
\]
We require $T\gg_KX^2$, which is checked in each application.
The relevant auxiliary polynomial has the form
\[
 \Psi_j(z;\mathbf h;\mathbf m)
   =a(\mathbf h)\{c_k z^3+B(\mathbf h,\mathbf m)z\},\qquad
 a(\mathbf h)=h_1\cdots h_j,
\]
where $c_k\ne0$ is fixed and integral, and $B$ is integral after multiplication
by a fixed denominator depending only on $k$.  Such fixed denominators can be
cleared once, changing only finitely many bad prime ideals and fixed lattices.
We sum over nonzero $h_i$ in boxes of volume $O(H_i)$ and over $m_i$ in boxes
of volume $O(M_iX^\epsilon)$.  Retain the actual small-support restrictions
from efficient differencing.  In particular $z$ lies in a fixed small box
around a nonzero archimedean center of size $X^{1/n}$, and
$|B_v|\ll_KX^{2/n}$ uniformly.  The Hessian of the cubic phase in $z$ is
nonzero on that box at every archimedean place.

Let $F_j(\alpha)$ be this sum, allowing a bounded smooth weight in $z$ with
uniformly bounded scaled derivatives and any added linear archimedean phase.
To arrange this smoothness without a boundary hypothesis, start with a
nonnegative $C^\infty$ weight that is at least one on the distinguished
ordinary root box and supported in a slightly larger tiny box.  The initial
solution count is thereby enlarged.  Every Cauchy or differencing step
produces products of translates of this weight.  The affine changes in $z$
have coefficients bounded in terms of the fixed differencing depth, so
all derivatives remain uniformly bounded at scale $X^{1/n}$.  Its support
retains the small-increment restrictions and the nonzero center.  Extra
indicators may be enlarged to such products before the last absolute-value
estimate, since the underlying orthogonality count has nonnegative weights.
Alternatively, ordinary Fourier completion of sharp box intersections costs
only $X^\epsilon$; the estimates below are uniform in its linear frequencies.

Define
\[
 \mathfrak M=\{\alpha:\mathcal H_T(\alpha)\le C_KX\},\qquad
 \mathfrak m=\mathbb T_K\setminus\mathfrak M.
\]
Take a fixed sufficiently small $C_K$ when disjointness is wanted; changing
it by a fixed factor changes no bound below.

\begin{proposition}[Cubic auxiliary major and minor estimates]
\label{prop:g7-cubic}
Under the preceding assumptions,
\begin{align}
 \sup_{\alpha\in\mathfrak m}|F_j(\alpha)|
       &\ll_{K,\epsilon}X^{3/4+\epsilon}HM,
       \label{eq:g7-cubic-minor}\\
 \int_{\mathfrak M}|F_j(\alpha)|^4\,d\alpha
       &\ll_{K,\epsilon}(XHM)^4T^{-1}X^\epsilon.
       \label{eq:g7-cubic-major}
\end{align}
\end{proposition}

\begin{proof}[Minor-arc estimate]
Apply two ordinary differences in $z$ to a cubic summand and use H\"older
for the $h,m$ sums.  Zero differences contribute $O(X^3H)$ before the final
H\"older factor.  Every remaining linear coefficient is a fixed nonzero
multiple of $h_1\cdots h_j u_1u_2$.  For any prescribed nonzero value,
the number of tuples $(\mathbf h,u_1,u_2)$ producing it is $O(X^\epsilon)$,
by ideal divisors followed by the bounded-height unit count.  Such values
lie in a box of volume $O(HX^2)=O(T/X)$.
The linear sums are bounded by positive kernels of width $X^{-1/n}$.
Equation~\eqref{eq:g7-kernel}, with $J=T/X$ and $N=X$, bounds their total by
\[
 X^\epsilon\left(T/X+X+T/\mathcal H_T(\alpha)\right)
       \ll_K X^\epsilon T/X
 \qquad(\alpha\in\mathfrak m).
\]
Consequently the sum over $h$ of the fourth powers of the cubic sums is
$O(X^{3+\epsilon}H)$, uniformly in $m$.  H\"older over the $HM$ parameter
choices now gives $|F_j|^4\ll X^{3+\epsilon}(HM)^4$, as asserted.
\end{proof}

\subsection{Cubic complete sums and the major arcs}

For a good prime ideal $\mathfrak p$, $q=N\mathfrak p$, define
\[
 w_3(\mathfrak p^{3u+1})=q^{-u-1/2},\qquad
 w_3(\mathfrak p^{3u+2})=w_3(\mathfrak p^{3u+3})=q^{-u-1}
 \quad(u\ge0).
\]
At the finitely many bad primes, including those above $6c_k$, take
$w_3(\mathfrak p^e)=q^{-e/3}$.  Extend multiplicatively and put $w_3(1)=1$.
For any cubic polynomial with fixed unit leading coefficient at good primes,
and any primitive additive character modulo an ideal $\mathfrak q$,
its normalized complete sum is
\begin{equation}\label{eq:g7-complete}
 \ll_{K,\epsilon}(N\mathfrak q)^\epsilon w_3(\mathfrak q),
\end{equation}
uniformly in all integral lower coefficients, including a linear twist.
For completeness, at good primes translate to $az^3+bz$.
At modulus $\mathfrak p$ the finite-field Weil bound gives $2q^{-1/2}$.
For higher powers, stationary residue classes satisfy $3az^2+b=0$.
Simple classes give quadratic Gauss bounds $O(q^{-e/2})$.
A multiple class is $z\equiv0$, with $\mathfrak p\mid b$.
If $\mathfrak p\parallel b$, its contribution is zero for $e\ge3$;
if $\mathfrak p^2\mid b$, substitution $z=\pi y$ gives precisely the
factor $q^{-1}$ times the normalized sum at level $e-3$.
The base cases $e=2,3$ give $q^{-1}$.  These recurrences prove the displayed
weight, with an absolute factor per prime absorbed by $(N\mathfrak q)^\epsilon$.
At bad primes one may apply the local-field van der Corput estimate to the
third derivative.  Proposition~3.3 of Cluckers \cite{Cluckers2011}, \emph{Expositiones Mathematicae}
29 (2011), 371--386, gives $O_K(q^{-e/3})$ uniformly here: after division by
$6c_k$ the third derivative has fixed absolute value, and all other coefficients
have bounded local Gauss norm.  The cited local-field estimate is uniform under this normalization.

Two consequences, both elementary Euler-product calculations, are
\begin{align}
 \sum_{N\mathfrak q\le Y}(N\mathfrak q)w_3(\mathfrak q)^4
       &\ll_{K,\epsilon}Y^\epsilon,\label{eq:g7-wsum}\\
 \sum_{\mathfrak d\mid\mathfrak q}
       \frac{w_3(\mathfrak q/\mathfrak d)}{N\mathfrak d}
       &\ll_{K,\epsilon}(N\mathfrak q)^\epsilon w_3(\mathfrak q).
       \label{eq:g7-wconv}
\end{align}
For the first, the good-prime Euler factor is $1+O(q^{-1})$ and the remaining
terms are summable; for the second the local ratios form a geometrically
convergent sum up to a bounded factor per prime.

\begin{proof}[Major-arc estimate]
Write $\alpha=\gamma+\beta$ on an arc with conductor $\mathfrak q$ and
$(N\mathfrak q)A_T(\beta)\le C_KX$.
For fixed $h,m$, put $a=h_1\cdots h_j$ and
\[
 \mathfrak q_a=\mathfrak q/(\mathfrak q+(a)).
\]
This is exactly the conductor of $a\gamma$.
Poisson summation modulo $\mathfrak q_a$, followed by
\eqref{eq:g7-complete}, bounds the cubic sum by
\begin{equation}\label{eq:g7-poisson}
 \ll X^{1+\epsilon}w_3(\mathfrak q_a)
       A_{X^3}(a\beta)^{-1/2}.
\end{equation}
Here are the archimedean details.  At a real place, the second derivative is
bounded away from zero after division by its scale, and the second derivative
test gives $(1+X^{3/n}|a_v\beta_v|)^{-1/2}$.
At a complex place the real Hessian of the real part of the holomorphic cubic
has determinant a nonzero constant times $|a_v\beta_v z_v|^2$;
two-dimensional stationary phase gives the square of this factor.
These estimates are uniform in the linear Poisson and completion frequencies.
Outside the natural gradient box repeated integration by parts gives arbitrary
decay.  Ideal packing bounds the total number of relevant dual frequencies by
\[
 O_K\left(1+\frac{N\mathfrak q_a}{X}A_{X^3}(a\beta)\right)=O_K(1),
\]
because $|a_v|\ll H^{1/n}$ and the product-arc inequality holds.
Dyadic enlargement of the gradient box, with enough integrations by parts,
proves the same bound for the tails.  This proves \eqref{eq:g7-poisson} for
the complete sum, including all Poisson modes.

It remains to average over $h$.  On an archimedean dyadic box with radii
$r_{i,v}$ for $h_i$, let $U_i=\prod_v r_{i,v}^{d_v}$.
There are no nonzero $h_i\in\mathfrak b$ in such a box unless
$U_i\gg_KN\mathfrak b$; when nonempty, ideal packing gives
$O_K(U_i/N\mathfrak b)$ choices.
Split the condition $\mathfrak d\mid(h_1\cdots h_j)$ among divisor ideals
of the individual $h_i$.  The number of splittings is $O(X^\epsilon)$.
Thus summing \eqref{eq:g7-poisson} on this dyadic box, and using
\eqref{eq:g7-wconv}, gives its trivial cardinality times
$X^\epsilon w_3(\mathfrak q)$ times the archimedean decay.
The sum over dyadic radii converges toward their largest values: at each
place the counting factor has degree $d_v$ whereas the decay removes only
$d_v/2$.  Therefore
\[
 |F_j(\gamma+\beta)|
 \ll X^{1+\epsilon}HM\,w_3(\mathfrak q) A_T(\beta)^{-1/2}.
\]
Finally $\int_{V}A_T(\beta)^{-2}\,d\beta\ll_KT^{-1}$,
and there are at most $N\mathfrak q$ center classes of conductor
$\mathfrak q$.  Sum the fourth power of the preceding bound using
\eqref{eq:g7-wsum}.  Overlaps, if any, only enlarge this upper bound.
This proves \eqref{eq:g7-cubic-major}.
\end{proof}

\subsection{The pure smooth terminal process}

Let $\lambda_t$ be any permissible smooth $2t$-moment exponent, extended to
real $t\ge0$ by H\"older interpolation.  For every $u\ge0$, Proposition
\ref{prop:g7-cubic} implies
\begin{equation}\label{eq:g7-M1-terminal}
 \int|F_j(\alpha)g_Q(\alpha)^{2u}|\,d\alpha
 \ll X^{1+\epsilon}HM
 \left\{X^{-1/4}Q^{\lambda_u}
       +Q^{-k/4+(3/4)\lambda_{4u/3}}\right\}.
\end{equation}
The minor term is the supremum estimate times the full smooth moment.
The major term is H\"older with exponents $4$ and $4/3$, using
$T=Q^k$.  This is the pure smooth $M_1$ process; it needs no extra ordinary
variables inside the terminal integral and no smooth major-arc asymptotic.

\section{Auxiliary quadratic terminals}
\label{sec:quadratic-terminals}
\label{sec:eighth-quadratic-terminal}
\label{sec:ninth-terminal}

The residual degrees two, three, and four in efficient differencing
admit a common treatment. Ordinary differences reduce them to a
quadratic sum and separate coefficient sums of degree $2k$.
We combine their moment estimates on one shift space, then integrate
the resulting estimate against a restricted-height smooth mean.

Fix $k\ge4$ and a depth $j$ such that $d=k-j\in\{2,3,4\}$.
Put $r=d-2$ and
\[
 \begin{gathered}
 M_i=X^{\phi_i},\quad H_i=XM_i^{-k},\quad
 \Phi_i=\sum_{a\le i}\phi_a,\quad q_i=1-\Phi_i,
 \qquad0\le\phi_i<1/k,\\
 M=\prod_{i\le j}M_i,\quad H=\prod_{i\le j}H_i,\quad
 Q_0=X/M,\quad T=Q_0^k=X^dH,\\
 A=XHM=X^{a_0},\quad a_0=1+j-(k-1)\Phi_j,\qquad
 H_*=X^rH=X^{h_*},\quad h_*=k-2-k\Phi_j.
 \end{gathered}
\]
All these scales are volumes. The conductor height is normalized
as in Section~\ref{sec:coefficient-sampling-common}, and
$\mathcal H_T(\alpha)\ll_KT^{1/2}$. The parameters are fixed before
the smoothness exponent is chosen; factors $X^{O_K(\eta)}$ are
absorbed in the final $X^\epsilon$.

\subsection{Supports and a common quadratic completion}

Partition the balanced principal divisors into dyadic norm annuli
and finitely many small archimedean boxes. On one piece each
$|m_{i,v}|$ has a prescribed radius up to fixed factors.
While the expressions are nonnegative orthogonality counts,
enlarge the ordinary and shift ranges to independent Cartesian
boxes. The constraints $|h_{i,v}m_{i,v}^k|\ll\delta X^{1/n}$
change by fixed factors only. Taking the original relative
diameter small enough keeps all ordinary vertices in a common
box about a nonzero center at every embedding, through all
subsequent ordinary differences. Centering congruences can be
discarded after their fixed denominators are cleared.
The number of pieces is $X^\epsilon$.

In a fixed integral basis take the ordinary box to be a
parallelepiped. Its intersections with translates after ordinary
differencing depend only on the ordinary shifts; the coefficient
sets and efficient-shift boxes remain independent. These
enlargements were made in counts before a terminal absolute value
was introduced.

For the centered efficient-difference polynomial one has the generating
expansion
\[
 \Psi_j(z;\mathbf h,\mathbf m)
 =2^jk!\Bigl(\prod_{i\le j}h_i\Bigr)
 \left\{\frac{z^d}{d!}
   +\frac{z^{d-2}}{6(d-2)!}\sum_{i\le j}h_i^2m_i^{2k}
   +\text{terms of degree at most }d-4\right\}.
\]
Apply $r=d-2$ ordinary differences with nonzero shifts $u_1,\ldots,u_r$.
With
\[
 \xi=2z+u_1+\cdots+u_r,\qquad
 \Gamma=\Bigl(\prod_{i\le j}h_i\Bigr)\Bigl(\prod_{a\le r}u_a\Bigr),
\]
the resulting polynomial is exactly
\begin{equation}\label{eq:quadratic-terminal-exact}
 \frac{2^jk!}{24}\Gamma
 \left(3\xi^2+\sum_{a\le r}u_a^2
                         +4\sum_{i\le j}h_i^2m_i^{2k}\right).
\end{equation}
This formula also covers $r=0$, in which case $\xi=2z$.
All constants are fixed; denominators can be cleared once.

Let $F_j$ be the original auxiliary sum and $\mathcal G_r$ its sum after
these ordinary differences, including the ordinary shifts. Cauchy's
inequality over the existing parameter tuple, before separating the
coefficient sums, gives
\begin{equation}\label{eq:quadratic-terminal-unseparated}
 |F_j|^{2^r}\ll
       X^{2^r-1}(HM)^{2^r}
       +X^{2^r-r-1}(HM)^{2^r-1}|\mathcal G_r|.
\end{equation}
Only $r=0,1,2$ is needed. For $r=1$, the zero ordinary shift contributes
$X(HM)^2$ and the remaining term is $HM|\mathcal G_1|$.
For $r=2$, Cauchy on the original tuple and the first shift gives
$|\mathcal G_1|^2\ll X^3(HM)^2+XHM|\mathcal G_2|$; substitution gives
\eqref{eq:quadratic-terminal-unseparated}. For $r=0$, it is the harmless bound
$|F_j|\le HM+|\mathcal G_0|$.

For each ordinary shift tuple, the remaining $z$-domain is an intersection
of translates of the same parallelepiped. Fourier completion in the fixed
integral basis gives an integrable kernel dominated by a common majorant
of $L^1$ norm $O((\log X)^n)$. The kernel may depend on $\mathbf u$, but
is independent of $\mathbf h,\mathbf m$. The condition
$\xi\equiv\sum_a u_a\pmod{2\mathcal O_K}$ is imposed by finitely many
additive characters, again independent of $\mathbf m$. Thus after one
common integration in a linear modulation $t$, the absolute value of
$\mathcal G_r$ is bounded by $X^\epsilon$ times
\begin{equation}\label{eq:quadratic-terminal-completion}
 \sup_t\sum_{\mathbf h,\mathbf u}
       |Z_t(c_0\alpha\Gamma)|
       \prod_{i=1}^j|L_i(c_i\alpha\Gamma h_i^2)|,
\end{equation}
uniformly in $t$, where
\[
 Z_t(\beta)=\sum_{z\in\mathcal B(X)}e_K(\beta z^2+tz),\qquad
 L_i(\beta)=\sum_{m\in\mathcal C_i}e_K(\beta m^{2k}),
\]
and $c_0,c_i$ are fixed nonzero integers. Smooth completion gives the same
formula with bounded weights and the same coefficient tuple throughout.

\subsection{Moment resources on the shift space}

On a minor height shell write
\[
 \mathcal H_T(\alpha)\asymp X^\ell,
 \qquad 1\le\ell\le\frac{kq_j}{2}.
\]
The fixed-factor endpoint changes and smoothness perturbations are
$O(1/\log X)+O_K(\eta)$ and are absorbed in the final error.
All products of a fixed number of nonzero algebraic integers have only
$X^\epsilon$ representations in the relevant boxes, by ideal divisors
and bounded-height unit counting. Thus $\Gamma$ may be grouped as a single
variable in $\mathcal B(H_*)$ for nonnegative moment estimates.

Majorize the shift box by a positive Gaussian and expand
the quadratic second moment. The diagonal $z^2=w^2$ contributes
$O(XH_*)$, and each nonzero value of $(z-w)(z+w)$ has
$X^\epsilon$ preimages. After the Gaussian sum in the shift
variable is evaluated as a nonnegative periodized kernel, the
common modulation can be removed in absolute value.
The linear-kernel estimate gives
\[
 \sum_a|Z_t(c\alpha a)|^2
 \ll X^\epsilon\left(XH_*+X^2+H_*+
                \frac{H_*X^2}{\mathcal H_{H_*X^2}(\alpha)}\right).
\]
Here $H_*X^2=T$. Dividing by the trivial size $H_*X^2$
gives the saving $\min\{1,h_*,\ell\}$, hence on the minor shells
\begin{equation}\label{eq:quadratic-terminal-z2}
 z_2=\min\{1,h_*\}.
\end{equation}
In particular $z_2=1$ when $H_*\ge X$.
The sampled fourth moment of Section~\ref{sec:quad-sampling-field}
gives the additional resources
\[
 z_4=\min\{2,\ell,h_*\},\qquad
 z_4^{\rm ref}=\min\{2,2\ell-h_*\}.
\]
Each is used only on intervals where its hypotheses hold.
The trivial bound is also available, with saving zero.

For coefficient block $i$, group all shifts except $h_i$ into a variable
of volume
\[
 J_i=X^r\prod_{a\ne i}H_a,
 \qquad \log_X J_i=k-3-k\Phi_j+k\phi_i.
\]
The conductor identity is
$J_iH_i^3M_i^{2k}=T$.
For an even coefficient order $p$, raise $L_i$ to the integer power $p/2$.
Its Fourier coefficients have support in a box of volume
$N_i\ll X^\epsilon M_i^{2k}$, and their squared norm is at most
\[
 X^\epsilon M_i^{p-2k+\Delta_p^{(2k)}}.
\]
This follows by orthogonality and coefficient domination by the full
$2k$th-power smooth moment. The exponent $\Delta_p^{(2k)}$ is obtained
from the independently proved field seed curve. The refined coefficient
sampling estimate, with
$Y\le c_K\min\{N_i,J_i,(J_iN_iH_i^{-3})^{1/2}\}$, gives saving
\begin{align}
 b_{i,p}^{\rm ref}(\ell)
 &=\min\left\{2k\phi_i,\ k-3-k\Phi_j+k\phi_i,
       \frac{k-6-k\Phi_j+6k\phi_i}{2},\ \ell-2+2k\phi_i\right\}
       -\Delta_p^{(2k)}\phi_i.\label{eq:quadratic-terminal-refined}
\end{align}
The cubic-height sampling estimate \eqref{eq:k5-sampling-strong} gives the alternative
\begin{align}
 b_{i,p}^{\rm cube}(\ell)
 &=\min\left\{2k\phi_i,\ k-3-k\Phi_j+k\phi_i,
       \frac{k-6-k\Phi_j+6k\phi_i}{2},\ \ell/3\right\}
       -\Delta_p^{(2k)}\phi_i.\label{eq:quadratic-terminal-cube}
\end{align}
The coefficient normalization is exact: sampling produces
$J_iH_i M_i^{p+\Delta_p^{(2k)}}$ times the reciprocal saving, whereas the
trivial normalized size is $J_iH_iM_i^p$.

Proposition~\ref{prop:k5-coeff2}, with $d=k$, gives the basic second coefficient resource
\begin{equation}\label{eq:quadratic-terminal-basic}
 b_{i,2}^{\rm basic}=
 \min\left\{\phi_i,\ k-6-k\Phi_j+4k\phi_i,
               q_j/4,\ \frac{kq_j-(4k-1)\phi_i}{6}\right\}.
\end{equation}
Put $b=b_{i,2}^{\rm basic}$. On the present minor shells, the lemma's other term
$\mathcal H_T(\alpha)^{-1/k}$ is no larger than $X^{-b}$, since
$b\le\phi_i<1/k$ and $\ell\ge1$.
A block with $\phi_i=0$ may be bounded trivially, including its fixed smoothness factors.

Each displayed coefficient resource bounds the sum of $|L_i|^p$, divided
by the common shift volume $H_*=X^rH$, over the \emph{same} shift tuple in \eqref{eq:quadratic-terminal-completion}
by $X^\epsilon M_i^pX^{-b_{i,p}(\ell)}$.
The quadratic resources likewise bound these normalized sums of $|Z_t|^p$ by
$X^\epsilon X^pX^{-z_p(\ell)}$.
Grouping products only enlarges nonnegative sums, so no independence between
these factors is needed.

\subsection{A combined minor-arc terminal}

Choose nonnegative resource weights $w_R$ on a fixed height interval.
If a resource has order $p_R$ and belongs to factor $i$, impose
\[
 \sum_{R\text{ on factor }i}p_Rw_R\le1
       \quad(i=0,1,\ldots,j),\qquad \sum_R w_R\le1.
\]
Factor $0$ denotes the quadratic sum. Divide each factor by its trivial
size, $X$ for $Z_t$ and $M_i$ for $L_i$. Its unused power is bounded
trivially, and the unused H\"older mass is assigned to the constant
function $1$. Generalized H\"older on the common shift space then gives
\[
 |\mathcal G_r|\ll X^{r+1+\epsilon}HM\,X^{-S(\ell)},\qquad
 S(\ell)=\sum_Rw_Rb_R(\ell).
\]
Combining this with \eqref{eq:quadratic-terminal-unseparated} proves
\begin{equation}\label{eq:quadratic-terminal-pointwise}
 |F_j(\alpha)|\ll A X^{\epsilon-
               2^{-(k-j-2)}\min\{1,S(\ell)\}}.
\end{equation}
In particular, the Weyl root is applied to the entire combined saving
$S(\ell)$, after all resources have been coupled by one H\"older inequality.

Let the remaining smooth order be $2u$. Choose nonnegative
$a,c,\gamma$ and an even order $v$ such that
\[
 a+c=1,\qquad \gamma+2a=1,\qquad cv=2u.
\]
If $a>0$, assume the auxiliary second-moment estimate
$\int|F_j|^2\ll X^\epsilon A$, supplied by
Lemma~\ref{lem:field-general-seed-aux} when its hypotheses hold.
The usual choice is $a=0$, $c=\gamma=1$, and $v=2u$.
On a height shell, H\"older gives
\[
 \int_{\rm shell}|F_j|\,|g_{Q_0}|^{2u}
 \le(\sup_{\rm shell}|F_j|)^\gamma
       \left(\int|F_j|^2\right)^a
       \left(\int_{\rm shell}|g_{Q_0}|^v\right)^c.
\]
The restricted-height smooth theorem and the pointwise estimate
therefore give the exponent
\begin{equation}\label{eq:quadratic-terminal-shell}
 C(\ell)=(\gamma+a)a_0+c(v-k)q_j+
       \frac{2c\Delta_v^{(k)}}k\ell
       -\gamma2^{-r}\min\{1,S(\ell)\}.
\end{equation}
Each resource is a minimum of affine functions of $\ell$,
hence concave. Its nonnegative weighted sum and the minimum
of that sum with one are also concave. Thus $C(\ell)$ is convex
on every interval with fixed resource weights, and its maximum
is attained at an endpoint.

\subsection{Uniform auxiliary major arcs}
\label{sec:aux-major-24}

For each auxiliary tuple write $a=h_1\cdots h_j\ne0$ and
\[
 \Psi_j(z;\mathbf h;\mathbf m)=aP_{\mathbf h,\mathbf m}(z),
 \qquad P_{\mathbf h,\mathbf m}(z)=c_{k,j}z^d+\cdots,
 \qquad c_{k,j}\ne0.
\]
The leading coefficient is a fixed integer, and all other coefficients
are algebraic integers after clearing one fixed denominator if necessary.
The ordinary variable has a nonnegative smooth weight in the small box
retained by efficient differencing.  All scaled weight derivatives are
uniformly bounded.  Linear completion twists may be arbitrary, independently
for every auxiliary tuple.

\begin{proposition}\label{prop:aux-major-24}
Let $F_j$ be the resulting auxiliary sum.  On every broad arc
\[
 \alpha=\gamma+\beta,\qquad Q_\gamma A_T(\beta)\le C_KX,
\]
one has, for every $\epsilon>0$,
\begin{equation}\label{eq:aux-major-strong}
 |F_j(\gamma+\beta)|
 \ll_{K,\epsilon}X^{1+\epsilon}HM\,
            Q_\gamma^{-1/d}A_T(\beta)^{-1/2}.
\end{equation}
Consequently, with $L=Q_\gamma A_T(\beta)$,
\begin{equation}\label{eq:aux-major-height}
 |F_j(\gamma+\beta)|\ll X^{1+\epsilon}HM L^{-1/d}.
\end{equation}
The bound includes every Poisson mode, also when $d=2$.
\end{proposition}

\subsubsection{Uniform mixed complete sums}

The required complete-sum bound is slightly more general than the ordinary
monomial-plus-linear estimate.  If $P\in\mathcal O_K[z]$ has degree
$d\le4$ and fixed nonzero leading coefficient $c$, then for a primitive
character $\psi$ modulo $\mathfrak q$ and every integral $b$,
\begin{equation}\label{eq:aux-complete-general}
 (N\mathfrak q)^{-1}
 \left|\sum_{z\bmod\mathfrak q}\psi(P(z)+bz)\right|
 \ll_{K,c,d,\epsilon}(N\mathfrak q)^{-1/d+\epsilon},
\end{equation}
uniformly in the lower coefficients.  Those coefficients are allowed to
vary with $\mathbf h,\mathbf m$ and with the Poisson frequency.

Here is a local proof recording uniformity in the prime.  At a prime of
residue characteristic greater than $d$ and not dividing $c$, consider
an integral polynomial whose reduction has positive degree $r\le d$.
The normalized complete sum of conductor $\mathfrak p^E$ is
$O_d(q^{-E/r})$, where $q=N\mathfrak p$.
For $E=1$ this is the finite-field Weil bound (the linear case is zero).
For $E\ge2$, nonstationary residue classes contribute zero.
At a stationary residue class of multiplicity $m$ as a root of the reduced
derivative, translate to that class and substitute $z=z_0+\pi y$.
Let $s$ be the minimum valuation of the resulting nonconstant Taylor
coefficients.  Then $2\le s\le m+1\le r$; after dividing by $\pi^s$,
the new reduction degree $r'$ satisfies $1\le r'\le s$.
The class has measure $q^{-1}$, and induction supplies
\[
 q^{-1}O_d(q^{-(E-s)/r'})=O_d(q^{-E/r})
 \qquad(E>s).
\]
When $E\le s$, its trivial measure $q^{-1}$ is already at most $q^{-E/r}$.
A child with $r'=r$ can occur only when the reduced derivative has a single
root of multiplicity $r-1$.  Thus a same-degree recursion has only one
child and does not accumulate constants.  Every other branch lowers the
reduction degree; at most $r-1$ stationary residue classes occur.
Induction on reduction degree and conductor gives a constant depending
only on $d$ (one may take constants bounded by $(d-1)!$ for $d\le4$).
This proof applies over every finite residue extension and is unchanged
by ramification at such good primes.

At the finitely many remaining primes, the $d$th derivative is the fixed
nonzero constant $d!c$.  Cluckers' local-field van der Corput theorem
\cite[Proposition~3.3]{Cluckers2011} applies after putting
$f=(P+bz)/(d!c)$: then $f^{(d)}=1$, while the Gauss norm of
$f-f(0)$ is at most $|d!c|_\mathfrak p^{-1}$ uniformly in the integral
lower coefficients and in $b$.  Its permitted dependence on the fixed
local field, derivative order and Gauss norm consequently gives
$O_{K_\mathfrak p,c,d}(q^{-E/d})$ uniformly in these coefficients.
The fixed denominators from centering are absorbed into fixed fractional
lattices and fixed conductor factors.  Chinese remaindering and the
ideal-divisor bound absorb
products of the good-prime constants into $(N\mathfrak q)^\epsilon$.
This proves \eqref{eq:aux-complete-general}.

\subsubsection{Archimedean curvature and all Poisson modes}

For a fixed auxiliary tuple, put $P=X^{1/n}$ and write
$p_v(u)=P^{-d}P_{\mathbf h,\mathbf m,v}(Pu)$.
On the retained tiny support all increments $h_im_i^k$ are small relative
to the nonzero ordinary center.  The centered difference formula therefore
gives
\[
 p_v(u)=c_{k,j}u^d+O(\delta^2)
\]
in bounded $C^d$ norms, with $|u|\asymp_K1$.  Taking $\delta$ sufficiently
small ensures $|p_v''(u)|\gg_K1$ uniformly.  For $d=2$ this is immediate
from its fixed leading coefficient.
At a real place, the second-derivative test is uniform in any added linear
phase.  At a complex place, the Hessian determinant of
$\operatorname{Re}(t p_v(u))$ is a nonzero constant times
$-|t p_v''(u)|^2$, so two-dimensional stationary phase gives the corresponding
square decay.  After a fixed partition of the tiny support, the standard
uniform stationary-phase theorem \cite[Section~7.7]{Hormander2003} and integration
by parts give
\begin{equation}\label{eq:aux-arch-tails}
 \left|\int w_v(u)e_{K_v}(t p_v(u)-\zeta u)\,du\right|
 \ll_{K,B}(1+|t|)^{-d_v/2}
       \left(1+\frac{|\zeta|}{1+|t|}\right)^{-B}
\end{equation}
for every prescribed $B$.  A linear completion twist simply translates
$\zeta$; the constants are independent of its size.  The same conclusion
holds for products of translated smooth vertex weights, since their
scaled derivatives remain uniformly bounded.  Sharp completion, if used,
only contributes an arbitrary $X^\epsilon$ factor.

Let $\mathfrak q$ be the conductor of $\gamma$ and
\[
 \mathfrak q_a=\mathfrak q/(\mathfrak q+(a)),\qquad q_a=N\mathfrak q_a.
\]
This is the conductor of $a\gamma$.
Poisson summation of the fixed-tuple sum modulo $\mathfrak q_a$, followed
by \eqref{eq:aux-complete-general}, bounds it by
\[
 Xq_a^{-1/d+\epsilon}A_{X^d}(a\beta)^{-1/2}
 \sum_{\xi\in\mathfrak q_a^{-1}\mathfrak D^{-1}}
 \prod_v\left(1+
   \frac{P|\xi_v-\xi_{0,v}|}{1+P^d|a_v\beta_v|}\right)^{-B},
\]
where $\xi_0$ is the arbitrary shift from completion.
The normalizing factor $1/q_a$ from Poisson has already been combined
with the complete-sum bound.

Ideal packing in a translated product box bounds the number of relevant
dual points by
\[
 O_K\left(1+\frac{q_a}{X}A_{X^d}(a\beta)\right).
\]
Indeed the radii are $(1+P^d|a_v\beta_v|)/P$ and their product volume is
$A_{X^d}(a\beta)/X$.  Since $|a_v|\ll H^{1/n}$,
\[
 q_aA_{X^d}(a\beta)\ll_K Q_\gamma A_T(\beta)\ll_KX.
\]
For completeness this packing also sums all tails.  Enlarge each radius
by an independent factor $2^{r_v}$; the number of dual points grows by at
most $\prod_v2^{d_vr_v}$, while \eqref{eq:aux-arch-tails} supplies
$\prod_v2^{-Br_v}$.  Choosing $B>n+2$ makes the resulting series convergent,
uniformly in $\xi_0$ and in the eccentricity of the gradient box.
Thus the whole dual sum is $O_K(1)$, proving the individual bound
\begin{equation}\label{eq:aux-fixed-tuple}
 |S_{\mathbf h,\mathbf m}(\gamma+\beta)|
 \ll X^{1+\epsilon}q_a^{-1/d}A_{X^d}(a\beta)^{-1/2}.
\end{equation}

\subsubsection{Averaging the conductor greatest common divisors}

There are $O(MX^\epsilon)$ choices of $\mathbf m$.
For prescribed $a\ne0$, the number of bounded-height tuples
$h_1\cdots h_j=a$ is $O(X^\epsilon)$, by ideal divisors and the unit
lattice.  We may therefore sum over $a\in\mathcal B(H)\setminus\{0\}$.
Write $Q=Q_\gamma$ and $A_\beta=A_T(\beta)$. The elementary inequality
\[
 A_{X^d}(a\beta)\gg_K A_\beta|Na|/H
\]
gives
\[
 A_{X^d}(a\beta)^{-1/2}
 \ll_K\min\{1,(H/(A_\beta|Na|))^{1/2}\}.
\]
For an ideal $\mathfrak b$ of norm $b$, the small-norm ideal count
\[
 \#\{a\in\mathfrak b\cap\mathcal B(H):0<|Na|\le Y\}
       \ll H^\epsilon Y/b
\]
has no points for $Y<b$.  Dyadic summation consequently gives, including
thin scales,
\[
 \sum_{0\ne a\in\mathfrak b\cap\mathcal B(H)}
       \min\{1,(H/(A_\beta|Na|))^{1/2}\}
 \ll H^\epsilon\frac{H}{bA_\beta^{1/2}}.
\]
Split according to $\mathfrak b=\mathfrak q+(a)$, and enlarge each class
to all nonzero $a\in\mathfrak b$.  Equation~\eqref{eq:aux-fixed-tuple} then
yields
\begin{align*}
 |F_j(\gamma+\beta)|
 &\ll X^{1+\epsilon}M Q^{-1/d}
       \sum_{\mathfrak b\mid\mathfrak q}(N\mathfrak b)^{1/d}
         \frac{H}{(N\mathfrak b)A_\beta^{1/2}}\\
 &\ll X^{1+\epsilon}HM Q^{-1/d}A_\beta^{-1/2},
\end{align*}
since $\sum_{\mathfrak b\mid\mathfrak q}(N\mathfrak b)^{-1+1/d}
\ll Q^\epsilon$.  This proves \eqref{eq:aux-major-strong}.
Because $A_\beta\ge1$ and $d\ge2$, it implies \eqref{eq:aux-major-height}.

\subsubsection{The major terminal}

Let $v$ be an even smooth order with permissible excess $\Delta_v^{(k)}$.
The restricted-height theorem bounds its mean on the union of arcs
of height at most $L$ by
\[
 Q_0^{v-k+\epsilon}L^{2\Delta_v^{(k)}/k}.
\]
Summing dyadic shells $1\le L\le C_KX$ with
\eqref{eq:aux-major-height} gives, provided $X\le Q_0^{k/2}$,
\begin{equation}\label{eq:quadratic-terminal-major-integral}
 \int_{\mathfrak M}|F_jg_{Q_0}^v|
 \ll X^\epsilon A Q_0^{v-k}
       X^{\max\{0,2\Delta_v^{(k)}/k-1/d\}}.
\end{equation}
Here $h_*=kq_j-2>0$ implies the required scale inequality.
An equality in the exponent contributes only a logarithm.

For the counted recurrence use the original smooth order $2u$,
without the optional auxiliary redistribution. Its major terminal is
\begin{equation}\label{eq:quadratic-terminal-major}
 C_{\rm maj}=a_0+(2u-k)q_j+
              \max\{0,2\Delta_{2u}^{(k)}/k-1/(k-j)\}.
\end{equation}

\subsection{The counted recurrence}

Let $\lambda_t$ denote the current full exponent at smooth
moment order $2t$. Choose integer half-orders
$r_1=s-1,\ldots,r_j=u$ with $0<r_{i+1}<2r_i$.
Start $C$ with the maximum of
\eqref{eq:quadratic-terminal-major} and every endpoint value of
\eqref{eq:quadratic-terminal-shell} in the assigned intervals.
The principal-block counted recurrence gives, for $i=j-1,\ldots,1$,
\[
 \begin{split}
 D_i&=1+i-k\Phi_i+\Phi_{i+1}+q_{i+1}\lambda_{r_{i+1}},\\
 B_i&=\tfrac12\{q_i\lambda_{2r_i-r_{i+1}}
                   +i-(k-1)\Phi_i+(2r_{i+1}-1)\phi_{i+1}\},\\
 C&\longleftarrow B_i+\tfrac12\max(D_i,C).
 \end{split}
\]
The first-stage exponent is then
\[
 \lambda_s^{\rm new}=(2s-3)\phi_1+
           \max\{1+\phi_1+q_1\lambda_{s-1},C\}.
\]
Every subsisting half-order is integral, every backward stage satisfies
$0<2r_i-r_{i+1}<2r_i$, and each round reads only the exponents available
at its start. More precisely, fix a sufficiently small ordinary root box
$\mathcal B$ away from zero at every archimedean place, with a contained
ideal-smooth root set $\mathcal A$. Write $f_{\mathcal B}$ and
$g_{\mathcal A}$ for their unweighted exponential sums. The first-stage
formula bounds the localized joint moment
$\int|f_{\mathcal B}|^2|g_{\mathcal A}|^{2s-2}$ itself, with the support
enlargements described at the start of this section. Coefficient domination
in its nonnegative orthogonality count gives
\[
 \int|g_{\mathcal A}|^{2s}
 \le \int|f_{\mathcal B}|^2|g_{\mathcal A}|^{2s-2}.
\]
A fixed compact root annulus can be covered by finitely many such boxes;
$|\sum_{a=1}^N g_a|^{2s}\le N^{2s-1}\sum_{a=1}^N|g_a|^{2s}$ then gives
the required pure smooth moment on that annulus. Bounded nonnegative
weights are handled by the same coefficient domination, with constants
depending on their fixed bounds. The formulas give the field estimates evaluated by the exact parameter certificates.

\section{Fourth powers}
\label{sec:fourth-powers}\label{sec:q12-analytic}
The quartic argument needs two seed moments and three applications
of the cubic terminal. Write $b_s$ for the exponent of the smooth
moment of order $2s$.

\subsection{The two seed moments}
At depth one, the A processes of orders two and four give
\begin{align}
 T_3(b_4)&=\max\{3+2\phi_3,
      1+\tfrac32\phi_3+\tfrac12(1-\phi_3)b_4\},
      \label{eq:q12-T3}\\
 T_4(b_3,b_4)&=\max\{1+b_3+(6-b_3)\phi_4,
      \tfrac54+\tfrac{11}{4}\phi_4+
                        \tfrac34(1-\phi_4)b_4\}.
      \label{eq:q12-T4}
\end{align}
These are \eqref{eq:seed-A-terminal} and
\eqref{eq:seed-initial-output} with $k=4$. The auxiliary moments
follow from Lemma~\ref{lem:field-general-seed-aux}.
Choose
\[
 \phi_3=\frac{109}{1000},\qquad\phi_4=\frac7{50},\qquad
 B_3=\frac{161}{50},\qquad B_4=\frac{461}{100}.
\]
The branches have the exact values
\[
\begin{array}{c|cc|c}
 &\text{diagonal}&\text{auxiliary}&\text{target}\\\hline
 s=3&3.218&3.217255&3.22\\
 s=4&4.6092&4.60845&4.61
\end{array}
\]
The coefficient sums are at most $43/50$.
Lemma~\ref{lem:supersolution} therefore proves the permissible
exponents $b_3=161/50$ and $b_4=461/100$.

\subsection{The tenth smooth moment}
\begin{proposition}\label{prop:quartic-tenth}
For every fixed number field $K$ and $\epsilon>0$, the smooth
quartic sum satisfies
\begin{equation}\label{eq:q12-moment}
 \int_{\mathbb T_K}|g_4|^{10}\ll X^{9993/1600+\epsilon}.
\end{equation}
\end{proposition}
\begin{proof}
Use the cubic terminal of Section~\ref{sec:cubic-terminal} with
$k=4$, $j=1$, $s=5$, and $\phi_1=3/16$. Then
\[
 H=X^{1/4},\quad M=X^{3/16},\quad Q=X^{13/16},
 \qquad T=Q^4=X^{13/4}>X^2.
\]
The diagonal and cubic minor branches agree and give
\[
 \frac52+\frac{13}{16}b_4=\frac{9993}{1600}.
\]
In the cubic major branch interpolate
$b_{16/3}\le(2b_5+b_6)/3$ and use the Hua value $b_6=17/2$.
The resulting transformation of $c=b_5$ is
\[
 c\longmapsto\max\left\{\frac{9993}{1600},
       \frac{31}{16}+\frac{13}{32}c+
                           \frac{13}{64}\frac{17}{2}\right\}.
\]
Starting with the Hua value $27/4$, three steps give
\[
 \frac{27}{4}\longmapsto\frac{205}{32}
 \longmapsto\frac{6417}{1024}\longmapsto\frac{9993}{1600}.
\]
The exact calculation is reproduced in
Appendix~\ref{app:quartic-exact-certificate}.
\end{proof}

\subsection{The global quartic bounds}
\begin{proof}[Proof of Theorem~\ref{thm:quartic-refinements}]
The excess at order ten is $\Delta_{10}=393/1600$.
With $(w,S,u)=(20,12,10)$ the absorption margin is
\[
 D=2-8\Delta_{10}=\frac7{200}>0.
\]
Thus $U_{12}\ll X^8$, and the minor-arc saving in
Corollary~\ref{cor:even-transfer} is $D/20=7/4000$.
That corollary, with the local requirement retained, gives
\[
 G_{\mathbb K}(4)\le\max\{12,H_{\mathbb K}(4)\}.
\]
Theorem~\ref{thm:local-table} gives the universal upper bound
sixteen, the upper bound fifteen when no dyadic completion is
$\mathbb Q_2$, and the upper bound twelve when every dyadic residue
degree is at least two.

For $K\supset\mathbb Q(\sqrt5)$, the polynomial $t^2-t-1$ is
irreducible over $\mathbb F_2$. Every dyadic completion of $K$
therefore contains the unramified quadratic extension of
$\mathbb Q_2$, and its residue degree is even. Hence
$G_{\mathbb K}(4)\le12$.
\end{proof}

\section{Fifth powers}
\label{sec:k5-terminal}

Write $M_i=X^{\phi_i}$, $H_i=XM_i^{-5}$, $H=H_1H_2$,
$M=M_1M_2$, $Q=X/M$, and $T=Q^5=X^3H$.
The centered second difference for fifth powers is a cubic in $z$.
We adapt the sharper dissection of \cite[Lemma~6.1]{VW1995}.
After one ordinary difference in $z$, its quadratic part and its two
coefficient parts separate as
\[
 K(\alpha;h,h_1,h_2)
 L_1(\alpha;h,h_1,h_2)L_2(\alpha;h,h_1,h_2),
\]
where, up to fixed nonzero constants in $K$,
\[
 K=\sum_z e_K(\alpha hh_1h_2 z^2+\xi z),\qquad
 L_i=\sum_{m_i}e_K(c\alpha hh_1h_2h_i^2m_i^{10}+b_i(m_i)).
\]
The ordinary variable $h$ has volume $X$.

The support issue deserves to be explicit.  Partition each extracted
principal divisor into a dyadic norm annulus and then into finitely many
tiny archimedean boxes.  Its generator is balanced, so on one piece all
$|m_{i,v}|$ are comparable to $M_i^{1/n}$, up to the harmless original
smoothness loss.  The actual constraints
$|h_{i,v}m_{i,v}^5|\ll\delta X^{1/n}$ are then enlarged to independent
small boxes for $h_i$.  Every enlarged vertex remains in a slightly larger
tiny box about the same nonzero center.  This enlargement is made in the
nonnegative orthogonality count before an absolute value is introduced.
Consequently it preserves the nonvanishing of the differenced polynomial
and its derivative and yields the Cartesian auxiliary sums used here.
Smooth completion can alternatively be applied to the original vertex
weights.  Fourier inversion gives a linear twist in $z$ and twists of
degree five in $m_i$; the coefficient of the latter depends on $h_i$ and the
completion frequencies but is independent of the linear sampling variable
$hh_{3-i}$.  The estimates below are uniform in precisely these twists.

\subsection{The higher coefficient moment}
The coefficient-sampling proof in Section~\ref{sec:coefficient-sampling-common} has the
following strengthening.  If $V=H_0^3JN$ and
\[
 Y\le c_K\min\{J,N,(JN/H_0^3)^{1/2}\},
\]
then
\begin{equation}\label{eq:k5-sampling-strong}
 \sum_{0\ne h\in\mathcal B(H_0)}\sum_{j\in\mathcal B(J)}
 \left|\sum_a c_h(a)e_K(\alpha h^3ja)\right|^2
 \ll V^\epsilon JH_0N\{Y^{-1}+\mathcal H_V(\alpha)^{-1/3}\}
                      \sup_h\sum_a|c_h(a)|^2.
\end{equation}
The coefficients may depend on $h$ because the elementary bilinear
large-sieve bound is applied separately for each $h$.
Here are the additional details proving the exponent $-1/3$.
The small-height centers are identified as $h^3\gamma$, where
$L=Q_\gamma A_V(\alpha-\gamma)\ll V^{1/2}$, by the conductor separation
lemma.  Put $\mathfrak d=\mathfrak q+(h)^3$ and
$\mathfrak b=\prod\mathfrak p^{\lceil v_\mathfrak p(\mathfrak d)/3\rceil}$.
For $A=A_V(\alpha-\gamma)$ the norm-count argument bounds the contribution
of this divisor class by
\[
 V^\epsilon\frac{N\mathfrak d}{Q_\gamma}
       \frac{H_0}{A^{1/3}N\mathfrak b}
 \le V^\epsilon H_0L^{-1/3},
\]
since $N\mathfrak b\ge(N\mathfrak d)^{1/3}$ and
$N\mathfrak d\le Q_\gamma$.
There are only $V^\epsilon$ divisor classes.  The large-height classes and
the $J+N$ terms of the bilinear bound contribute $JH_0N/Y$.
This proves \eqref{eq:k5-sampling-strong}.

Apply this with $c_h(a)$ the coefficients of a $t$th power of a decic smooth
sum of volume $M_1$.  Its squared coefficient norm is bounded by the decic
$2t$th moment, even in the presence of the lower-degree twists: by
orthogonality and the triangle inequality, its absolute coefficient count
is dominated by the untwisted decic count.  Thus if $\mu_t$ is a permissible
decic $2t$-moment exponent, then
\[
 E_t(\alpha;J,H_0,M_1)
 \ll X^\epsilon JH_0M_1^{\mu_t+10}
             \{Y^{-1}+\mathcal H_{JH_0^3M_1^{10}}(\alpha)^{-1/3}\}.
\]

\subsection{The terminal estimate}
Let $\theta=\phi_1$, $\phi=\phi_2$, and write $U=X^u$.  Assume
\begin{gather}
 1/15\le\theta\le1/5,\quad 0<\phi\le\theta,\quad
 \phi<5\theta-1/3,\label{eq:k5-parameter1}\\
 0<u<\phi,\quad u<-1-5\theta+15\phi,\quad
 4u+\theta+\phi<1,\quad 6u+5\theta+24\phi<5.
 \label{eq:k5-parameter2}
\end{gather}
The strict inequalities absorb the small powers arising from principal
blocks and from dyadic completion.

Let $t\ge3$, and let $\mu_t$ be the decic $2t$-moment exponent.  Put
\[
 Z=XU^{1-1/t}\left(X^{1/3}M_1^{2t-10-\mu_t}\right)^{1/t},
 \qquad \rho=\frac{\log_X Z}{4}.
\]
On the cubic minor arcs $\mathcal H_T(\alpha)>C_KX$, ordinary differencing
and H\"older give
\[
 |F_2|^2\ll X(HM)^2+HM\,D^{1/2}E_1^{1/(2t)}E_2^{1/2-1/(2t)}.
\]
Here $D$ is the sum of $|K|^2$, $E_1$ the sum of $|L_1|^{2t}$, and
$E_2$ the sum of $|L_2|^{2t/(t-1)}$, over $h,h_1,h_2$.
The quadratic linear-kernel calculation gives
\[
 D\ll X^{2+\epsilon}H.
\]
In detail, after expanding $|K|^2$, the zero difference contributes $X^2H$;
for the other terms factor $z_1^2-z_2^2=(z_1-z_2)(z_1+z_2)$, group the
nonzero coefficient $hh_1h_2(z_1-z_2)$ by ideal divisors, and apply the
positive linear-kernel estimate at scale $T$.  Its minor-arc bound is
$T/X=X^2H$.

Group the product $hh_2$ in $E_1$, at a cost $X^\epsilon$.
Equation~\eqref{eq:k5-sampling-strong}, with
$J=XH_2$, $H_0=H_1$, $N=M_1^{10}$ and $Y=X^{1/3}$, is applicable by
\eqref{eq:k5-parameter1}.  It gives
\[
 E_1\ll X^{2/3+\epsilon}H M_1^{\mu_t+10}.
\]
For $E_2$, first remove the power exceeding two using $|L_2|\ll M_2$.
Proposition~\ref{prop:k5-coeff2}, with $J=XH_1$, $H_0=H_2$, $M=M_2$ and
$d=5$, applies by \eqref{eq:k5-parameter2}.  Since $U<X^{1/5}$, it yields
\[
 E_2\ll X^{1+\epsilon}H M_2^{2t/(t-1)}U^{-1}.
\]
Consequently, provided $\rho\le1/2$,
\begin{equation}\label{eq:k5-sharp-minor}
 \sup_{\mathfrak m}|F_2|\ll X^{1+\epsilon}HM Z^{-1/4}.
\end{equation}

The already-proved cubic complete-sum and Poisson estimates give
\[
 \int_{\mathfrak M}|F_2|^4\ll X^\epsilon(XHM)^4Q^{-5}.
\]
They apply uniformly to the completed support and linear twists just
specified, including all complex places and all nonzero Poisson modes.
For any permissible quintic smooth moments $\lambda_r$ and
$\lambda_{4r/3}$, H\"older on the major arcs and
\eqref{eq:k5-sharp-minor} on the minor arcs therefore prove
\begin{equation}\label{eq:k5-enhanced-terminal}
 \int|F_2g_Q^{2r}|
 \ll X^{1+\epsilon}HM
 \left\{Z^{-1/4}Q^{\lambda_r}
       +Q^{-5/4+(3/4)\lambda_{4r/3}}\right\}.
\end{equation}
This proves the sharper fifth-power terminal, using the degree-ten
moment and the coefficient estimates of Section~\ref{sec:coefficient-sampling-common}.


\subsection{The joint-moment certificate}
\label{sec:k5-certificate}

The principal-block inequality gives both smooth moments and joint
moments with two distinguished ordinary roots. The auxiliary estimates
at depths one and two follow from Lemma~\ref{lem:field-general-seed-aux}.
Start with $\lambda_0=0$, $\lambda_1=1$, $\lambda_2=2$, and
$\lambda_s=2s$ for $3\le s\le8$.

\subsubsection{The map being certified}
The A rows use \eqref{eq:seed-A-terminal} and the common
backwards recurrence. For the enhanced rows, the output is a joint
moment with two ordinary roots.

For the two enhanced rows take $r=6$, $t=2(S-1)-r$, and
\[
 \theta=\phi_1,\quad\phi=\phi_2,\quad q=1-\theta-\phi,
 \quad
 \rho=\frac14\left\{1+\frac{21u}{22}
                  +\frac{1/3+(34-34.227)\theta}{22}\right\}.
\]
The degree-ten exponent $34.227$ is supplied by Corollary~\ref{prop:tenth44}.
Put
\begin{align*}
 B&=\tfrac12\{(1-\theta)\lambda_t+1-4\theta+(2r-1)\phi\},\\
 D&=2-4\theta+\phi+q\lambda_r,\\
 C&=3-4(\theta+\phi)
        +\max\{-\rho+q\lambda_r,
                 q(-5/4+(3/4)\lambda_{4r/3})\}.
\end{align*}
The enhanced output is
\[
 T_S(\lambda)=\max\{1+\lambda_{S-1}
             +(2S-2-\lambda_{S-1})\theta,
                         (2S-3)\theta+B+\tfrac12\max(D,C)\}.
\]
Equation~\eqref{eq:k5-enhanced-terminal} proves this formula, including its
joint-moment interpretation.

\subsubsection{The exact supersolution}
All entries in the next table are rational numbers; the displayed integers
in the parameter columns are divided by $10^6$.
\[
\begin{array}{c|c|c|c|r|r|r|c}
 S&\text{row}&j&\ell&10^6\theta&10^6\phi&10^6u&b_S\\\hline
3&A&1&1&68130&&&3.1364\\
4&A&1&2&105600&&&4.4389\\
5&A&2&2&136590&72270&&5.9254\\
6&A&2&2&151340&113110&&7.5422\\
7&\text{enhanced}&2&-&163790&139365&139360&9.2726\\
8&\text{enhanced}&2&-&169960&138337&138330&11.0764
\end{array}
\]
The six strict differences $b_S-T_S(b)$ are, exactly,
\[
 \frac{62257}{2000000000},\quad
 \frac{2147}{25000000},\quad
 \frac{89351}{1000000000},\quad
 \frac{37509}{250000000},\quad
 \frac{128469}{500000000},\quad
 \frac{41387}{125000000}.
\]
The program included below verifies all parameter restrictions, interpolated
indices, and these rational inequalities without floating-point arithmetic.

Every branch is monotone in the input exponents, with coefficient
sum at most $47/50$. Lemma~\ref{lem:supersolution} therefore proves
\[
 \int_{\mathbb T_K}|f_5(\alpha)|^2|g_5(\alpha)|^{14}\,d\alpha
 \ll X^{11.0764+\epsilon}.
\]

\subsection{The global bound}
\label{sec:k5-global}
The preceding joint moment has excess
\[
 \Delta=\frac{191}{2500}<\frac1{13}.
\]
The ordinary minor-arc thirtieth moment and H\"older with weights
$1/14,13/14$ give
\[
 \int_{\mathfrak m_*}|f_5|^4|h_5|^{13}
 \ll X^{12-17/35000+\epsilon}.
\]
Here $h_5$ is the finite-product root sum in Section~\ref{sec:product-endgame}, so its
joint moments are dominated by the smooth moments up to an arbitrarily
small independent divisor loss.  Coefficientwise domination gives the
pure sixteenth moment first, and the loss-free absorption lemma gives
$U_{18}\ll X^{13}$ because
\[
 2-12\Delta=\frac{677}{625}>0.
\]
The broad major-arc annulus is controlled with ordinary moment order
$72/5>10$ and product moment order $18$, with H\"older weights $5/18,13/18$.
The finite-product core has a positive main term by the Poisson and
local-density argument in Section~\ref{sec:product-endgame}.  Primitive local width is at most ten
for fifth powers, and hence is below seventeen at every finite completion.
Totally positive unit balancing supplies the same positivity and norm
convention as in Section~\ref{sec:introduction}.  We conclude
\[
 \boxed{G_{\mathbb K}(5)\le17\quad\hbox{for every number field }K.}
\]

\section{Sixth powers}
\label{sec:sixth-powers}
The sixth-power certificate uses one auxiliary fourth moment at the
linear boundary of ordinary differencing. We record its extra factor.

\begin{lemma}\label{lem:sixth-linear-boundary}
For the degree-six auxiliary sums at depths $j\le4$, the second and
fourth moments used below have exponents
\[
 A_{j,\ell}=2^\ell-\ell+(2^\ell-1)(j-5\Phi_j)
                  +\mathbf1_{j+\ell=6}\Phi_j,
 \qquad \ell\in\{1,2\}.
\]
The parameters in the certificate satisfy the chronological divisor
conditions whenever they are needed.
\end{lemma}
\begin{proof}
For $j+\ell\le5$, apply Lemma~\ref{lem:field-general-seed-aux}.
Only the case $(j,\ell)=(4,2)$ remains. Its second moment is
$O(X^{1+\epsilon}U_4)$. After one ordinary difference the phase is
a nonzero constant times $zh_1\cdots h_5$, independent of the
four efficient divisors. Divisor counting in its nonzero value
gives an appended second moment
$O(X^{2+\epsilon}\prod_{i\le4}H_iM_i^2)$.
The padded-moment induction consequently gives
\[
 \int|F_4|^4\ll X^{2+\epsilon}U_4^3\prod_{i\le4}M_i.
\]
This is the stated extra factor $X^{\Phi_4}$.
\end{proof}

\subsection{The moment certificate}
\label{sec:g6-certificate}

Write $\mu_r=\lambda_{r/2}$ for an exponent of the full $r$th moment.
We fix $\mu_4=2$. All decimals and parameter numerators in this section
are exact rational numbers. The following eight rows suffice:
\begin{center}
\small
\begin{tabular}{@{}r l c c l r@{}}
\hline
$S$ & Cauchy path & $(\ell_1,\ell_2)$ & $u$ & $10^5(\phi_1,\ldots,\phi_d)$ & $b_{2S}$\\
\hline
3 & empty & $(1,1)$ & $8$ & $(4545)$ & 3.0915\\
4 & empty & $(2,2)$ & $8$ & $(8333)$ & 4.334\\
5 & $(5)$ & $(1,2)$ & $10$ & $(10673,5080)$ & 5.726\\
6 & $(6)$ & $(2,2)$ & $32/3$ & $(11855,8751)$ & 7.233\\
7 & $(7,6)$ & $(1,2)$ & $12$ & $(12981,10763,5551)$ & 8.852\\
8 & $(8,7)$ & $(2,2)$ & $40/3$ & $(13784,12076,8562)$ & 10.562\\
9 & $(9,8,8)$ & $(1,1)$ & $16$ & $(14593,13194,10858,1059)$ & 12.356\\
10 & $(10,9,9)$ & $(1,2)$ & $18$ & $(15303,14283,12101,6908)$ & 14.220\\
\hline
\end{tabular}
\end{center}

Here is a complete specification of the arithmetic. A path $(t_1,\ldots,t_{d-1})$
defines $s_1=S-1$ and $s_{j+1}=2s_j-t_j$. Put
$\Phi_j=\sum_{i\le j}\phi_i$ and $q_j=1-\Phi_j$.
At the terminal depth put $c=2s_d/u$ and choose nonnegative $a,b$ with
\[
 a+b+c=1,\qquad 2^{\ell_1}a+2^{\ell_2}b=1.
\]
If $\ell_1=\ell_2$, take $a=2^{-\ell_1}$ and $b=0$.
Define
\[
 A_{d,\ell}=2^\ell-\ell+(2^\ell-1)d
       -5(2^\ell-1)\Phi_d+\mathbf1_{d+\ell=6}\Phi_d,
\]
and initialize
\[
 T_d=aA_{d,\ell_1}+bA_{d,\ell_2}+cq_d\mu_u.
\]
For $j=d-1,\ldots,1$, put $p=2s_j-t_j$ and use
\[
 T_j=\frac12\left\{q_j\mu_{2t_j}+j-5\Phi_j+(2p-1)\phi_{j+1}
 +\max\left(1+j-6\Phi_j+\Phi_{j+1}+q_{j+1}\mu_{2p},T_{j+1}\right)\right\}.
\]
Finally the row output is
\[
 \mathcal T_{2S}(\boldsymbol\mu)=
 \max\left\{1+2(S-1)\phi_1+q_1\mu_{2S-2},
                         (2S-3)\phi_1+T_1\right\}.
\]
Only two nonintegral moments occur, and their meaning is fixed as
\[
 \mu_{32/3}=(2\mu_{10}+\mu_{12})/3,
 \qquad \mu_{40/3}=(\mu_{12}+2\mu_{14})/3.
\]
These are ordinary H\"older interpolations. Lemma~\ref{lem:sixth-linear-boundary} supplies the terminal estimates;
the backwards steps follow from the fundamental differencing inequality.

Substitution of the table's vector $\boldsymbol b$ gives:
\begin{center}
\begin{tabular}{@{}r r r r@{}}
\hline
$r$ & $b_r$ & $\mathcal T_r(\boldsymbol b)$ & $b_r-\mathcal T_r(\boldsymbol b)$\\
\hline
6 &3.0915&3.091234850&0.000265150\\
8 &4.334&4.333865305&0.000134695\\
10&5.726&5.725616976&0.000383024\\
12&7.233&7.232682700&0.000317300\\
14&8.852&8.851804270&0.000195730\\
16&10.562&10.561620060&0.000379940\\
18&12.356&12.355567340&0.000432660\\
20&14.220&14.219701320&0.000298680\\
\hline
\end{tabular}
\end{center}
The displayed outputs are exact terminating decimals. All parameter orderings
hold. The largest auxiliary side expression
$\Phi_I+6(\phi_{I-1}+\phi_I)$ is $1.99991<2$.

Every branch is monotone, and the coefficient sum in its variable
moment exponents is at most $23/25$. Lemma~\ref{lem:supersolution}
therefore shows that every entry is permissible. In particular,
\[
 \int|g_6|^{20}\ll X^{14.22+\epsilon}.
\]

Using the permissible excess $\Delta_{20}=0.22$, the twenty-four-variable
criterion in Section~\ref{sec:product-endgame} has
\[
 D=4-18\Delta_{20}=1/25,\qquad D/42=1/1050.
\]
Four ordinary roots and twenty product roots therefore give the required
broad-minor power saving. The primitive local bound $h_6\le12$ is more than
sufficient at twenty-four roots, completing the sixth-power assertion of
Theorem~\ref{thm:main-transfer}.

\section{Seventh powers}
\label{sec:seventh-powers}
The A and B seed processes are combined with the cubic terminal at
depth four. The auxiliary estimates are instances of
Lemma~\ref{lem:field-general-seed-aux}; all residual degrees needed
here are odd or belong to $\{2,4\}$.
For the cubic terminal, put $q_4=1-\Phi_4$ and use
\begin{equation}\label{eq:septic-cubic-terminal}
 C_4=5-6\Phi_4+
 \max\left\{-\frac14+q_4b_{s-4},\quad
 q_4\left(-\frac74+\frac34b\left(\frac{4(s-4)}3\right)\right)\right\}.
\end{equation}
Section~\ref{sec:cubic-terminal} proves this estimate when $7q_4>2$.
The remaining steps are \eqref{eq:seed-backwards} and
\eqref{eq:seed-initial-output}.

\subsection{The moment certificate}


In the table, $m_i=10^5\phi_i$.  Dashes mean that the corresponding parameter
does not occur.  Put $b_0=0$, $b_1=1$, $b_2=2$, and take $b_s$ from the table.
The last column is an integer lower bound for
$10^8\{b_s-T_s(\boldsymbol b)\}$, obtained by the exact recurrence
above.  It is strictly positive in every row.

\begin{center}
\small
\begin{tabular}{r c r r r r r r}
\hline
$s$ & Process & $b_s$ & $m_1$ & $m_2$ & $m_3$ & $m_4$ & Slack bound\\
\hline
3 & $A_{1}^{(1)}$ & 3.06493 & 3196 & -- & -- & -- & 52177 \\
4 & $A_{1}^{(2)}$ & 4.26513 & 6818 & -- & -- & -- & 8692 \\
5 & $B_{2,5}^{(2)}$ & 5.59014 & 8699 & 3541 & -- & -- & 10185 \\
6 & $A_{2}^{(2)}$ & 7.01542 & 9641 & 6937 & -- & -- & 12539 \\
7 & $B_{3,6}^{(2)}$ & 8.54214 & 10564 & 8803 & 4059 & -- & 14896 \\
8 & $A_{3}^{(2)}$ & 10.15371 & 11203 & 9889 & 6902 & -- & 12594 \\
9 & $B_{4,8}^{(2)}$ & 11.84806 & 11874 & 10797 & 8946 & 4151 & 16152 \\
10 & $A_{4}^{(2)}$ & 13.60673 & 12330 & 11453 & 9898 & 6609 & 13579 \\
11 & $A_{4}^{(2)}$ & 15.42553 & 12805 & 12028 & 10870 & 8586 & 14177 \\
12 & $A_{4}^{(2)}$ & 17.29695 & 13252 & 12659 & 10684 & 10188 & 17123 \\
13 & $M_4$ & 19.19477 & 13392 & 12953 & 12116 & 10714 & 14754 \\
14 & $M_4$ & 21.11765 & 13559 & 13070 & 12326 & 10846 & 15886 \\
\hline
\end{tabular}
\end{center}

Every A and B row satisfies the chronological divisor inequalities
of Lemma~\ref{lem:field-general-seed-aux}; these follow directly from
the displayed integers $m_i$.  Every $M_4$ row satisfies $7q_4>2$, the hypothesis of the cubic
dissection.  All required interpolation arguments lie in $[0,14]$.


The coefficient sum in every branch is at most $97/100$.
Lemma~\ref{lem:supersolution} proves
\begin{equation}\label{eq:g7-smooth28}
 \int_{\mathbb T_K}|g_7|^{28}\ll X^{21.11765+\epsilon}
              \ll X^{21+59/500+\epsilon}.
\end{equation}

\begin{corollary}\label{thm:seventh-transfer}
For every number field $K$, $G_{\mathbb K}(7)\le31$.
\end{corollary}
\begin{proof}
Use the convenient excess $\Delta=59/500$ in \eqref{eq:g7-smooth28}.
The ordinary critical moment has order $56$, and
$6/56+25/28=1$. Thus six ordinary and twenty-five product roots
have minor-arc saving
\[
 \frac6{56}-\frac{25}{28}\Delta=\frac1{560}>0.
\]
The bootstrap to order $32$ has margin
$4-24\Delta=146/125>0$.
The annular ordinary order is $192/7>14$, and $h_7\le14$.
Theorem~\ref{thm:general-transfer} now proves the assertion.
\end{proof}

\section{Eighth powers}
\label{sec:eighth-powers}
\subsection{The degree-eight and degree-sixteen seeds}
\label{sec:field-seed-moments}
The exact seed data in Appendix~\ref{app:eighth-field} consist of
eighteen rows at degree eight and twenty-eight rows at degree sixteen.
They use the A and B transformations
\eqref{eq:seed-A-terminal}--\eqref{eq:seed-initial-output}.
Each row satisfies the auxiliary hypotheses of
Lemma~\ref{lem:field-general-seed-aux}. In particular, the rows
requiring a higher even residual curve count have $\Phi_j\le1/3$;
the chronological alternative checks all prefixes.

The least strict supersolution margins are
\[
 \frac{1164945381}{800000000000000},\qquad
 \frac{21545278901}{80000000000000000},
\]
and the coefficient sums are at most
$976509627/10^9$ and $9962585803/10^{10}$, respectively.
Lemma~\ref{lem:supersolution} proves the permissibility of both
curves from the elementary second and fourth moments and trivial
higher estimates. Zero extraction exponents have nominal threshold
one; the fixed smoothness powers are absorbed in $X^\epsilon$.

\subsection{The eighth-power quadratic policies}
Apply Section~\ref{sec:quadratic-terminals} with $k=8$.
For the first policy take $u=8$, $v=20$, $a=1/5$, $c=4/5$,
and $\gamma=3/5$. Its chronological divisor inequalities
supply the required auxiliary second moment.
For the other policies take $a=0$, $c=\gamma=1$, and $v=2u$.

The fixed policies have the following terminal data:
\[
\begin{array}{c|c|c|c|c|c}
 2s&j&(r_1,\ldots,r_j)&v&h_*&2^{-(k-j-2)}\\\hline
28&6&(13,12,11,10,9,8)&20&1.69428296&1\\
30&5&(14,13,12,11,11)&22&1.83213176&1/2\\
32&5&(15,14,13,12,11)&22&1.78352592&1/2\\
34&4&(16,15,14,14)&28&2.39658688&1/4
\end{array}
\]
Thus the quadratic second resource is available in every policy.
Only quadratic orders two and four are used. Coefficient orders are even
and at most $60$; the degree-sixteen curve through half-order $30$ supplies
every requested moment. The sole cube resource occurs at depth six and
coefficient order $60$. All other higher coefficient resources are refined
ones. The exact verifier checks the two quadratic fourth resources,
the full basic minimum \eqref{eq:quadratic-terminal-basic}, and every factor
load and H\"older mass.

\subsection{The moment and the global bound}
\label{sec:eighth-global}
\begin{theorem}[A smooth eighth-power moment]
\label{thm:eighth-moment}
For every fixed number field $K$ and every $\epsilon>0$, there is
$\eta>0$ such that the degree-eight ideal-smooth sum satisfies
\[
 \int_{\mathbb T_K}|g_8(\alpha;X,R)|^{34}\,d\alpha
 \ll_{K,\epsilon}X^{26+\Delta_{34}+\epsilon}
 \qquad(1\le R\le X^\eta),
 \qquad
 \Delta_{34}=\frac{5857852793}{50000000000}.
\]
The same assertion holds for the fixed boxes and bounded nonnegative
weights used in the representation argument.
\end{theorem}

\begin{proof}
The exact program in Appendix~\ref{app:eighth-field} contains all eighteen
degree-eight seed rows, all twenty-eight degree-sixteen seed rows, and
the four terminal policies, including every height interval and resource
weight. Its first part checks every seed supersolution inequality and
all side conditions used in the seed proof. The contraction argument in
Section~\ref{sec:field-seed-moments} starts these curves from the elementary
second and fourth moments and the trivial higher bounds.

The second part makes twelve simultaneous rounds at half-orders
$14,15,16,17$. Each round reads the exponents available at its beginning.
Before and after a round, H\"older interpolation over every pair of
available integer half-orders may improve an entry. This is a valid
operation on permissible field exponents. Each policy checks all
H\"older masses and factor loads, every coefficient-moment index, each
height interval endpoint, and the backwards counted recurrence.
Section~\ref{sec:eighth-quadratic-terminal} proves that these are the
exponents of the actual field estimates. Convexity of the shell exponent
justifies checking only the two endpoints of each of the eight assigned
height intervals.

After each policy application the exponent is rounded \emph{up} to a
multiple of $10^{-12}$, and the better of that value and the previously
available interpolated value is retained. All calculations use rational
arithmetic. The final outputs are:
\[
\begin{array}{c|c|c}
 \text{moment order}&\text{full permissible exponent}&\text{excess over }2s-8\\\hline
 28&20.360308744711&0.360308744711\\
 30&22.258467366407&0.258467366407\\
 32&24.178873819111&0.178873819111\\
 34&26.117157055860&0.117157055860
\end{array}
\]
The last entry is the exponent in the theorem. Only finitely many
applications are needed for any prescribed final $\epsilon$; choose
the smoothness parameter after those applications and after the fixed
rational parameters. This absorbs the smoothness and divisor losses
without consuming the strict final margin.
\end{proof}


\begin{corollary}\label{thm:eighth-transfer}
For every number field $K$, $G_{\mathbb K}(8)\le38$.
\end{corollary}
\begin{proof}
Apply Corollary~\ref{cor:even-transfer} with $(w,S,u)=(72,38,34)$.
Its exact margin is
\[
 D=4-34\Delta_{34}=\frac{416502519}{25000000000}>0.
\]
The local bound $h_8\le35$ supplies the required primitive
representations in thirty-eight variables.
\end{proof}

\section{Ninth powers}
\label{sec:ninth-field-seeds}

\subsection{The degree-nine base and refined curve}

The degree-nine seed curve and the degree-eighteen coefficient
curve are obtained from the common A and B transformations.
The exact parameters are recorded in Appendix~\ref{app:ninth-field}.

\begin{proposition}\label{prop:ninth-field-seed-curve}
Put $b_0=0,b_1=1,b_2=2$.  Every entry $\lambda_s$ in the last column
of the following table is a permissible exponent:
\[
 \int_{\mathbb T_K}|g_9(\alpha;X,R)|^{2s}\,d\alpha
       \ll_{K,s,\epsilon}X^{\lambda_s+\epsilon}
       \qquad(1\le R\le X^\eta)
\]
for a sufficiently small $\eta>0$ chosen after $\epsilon$.
The base exponents $b_s$ in the second column are also permissible.
\end{proposition}

\begin{center}
\begin{tabular}{@{}r r c r@{}}
\hline
$s$&Base $b_s$&Process&Refined $\lambda_s$\\
\hline
3&3.035816&$A_{1}^{(1)}$&3.035807234000\\
4&4.182300&$A_{1}^{(2)}$&4.182294748390\\
5&5.419716&$A_{2}^{(1)}$&5.419711470690\\
6&6.738319&$B_{2,6}^{(2)}$&6.738313329014\\
7&8.135645&$A_{2}^{(2)}$&8.135638882429\\
8&9.603938&$B_{3,7}^{(2)}$&9.603930891210\\
9&11.142513&$A_{3}^{(2)}$&11.142506096925\\
10&12.746354&$A_{4}^{(2)}$&12.746346059899\\
11&14.410598&$A_{4}^{(2)}$&14.410586823381\\
12&16.129237&$A_{5}^{(2)}$&16.129216870240\\
13&17.896005&$A_{5}^{(2)}$&17.895964268487\\
14&19.705700&$B_{6,11}^{(2)}$&19.705619653366\\
15&21.551232&$A_{6}^{(2)}$&21.551085549629\\
16&23.428999&$A_{6}^{(2)}$&23.428755886157\\
17&25.333544&$A_{6}^{(2)}$&25.333173492456\\
18&27.259337&$A_{6}^{(2)}$&27.258814572150\\
19&29.202038&$A_{6}^{(2)}$&29.201360351041\\
20&31.157968&$A_{6}^{(2)}$&31.157131085916\\
21&33.124215&$A_{6}^{(2)}$&33.123229933630\\
22&35.098519&$A_{6}^{(2)}$&35.097377231005\\
\hline
\end{tabular}
\end{center}
Here $A_j^{(\ell)}$ denotes the A terminal at depth $j$ using auxiliary
order $2^\ell$, and $B_{j,t}^{(\ell)}$ denotes the B terminal with smooth
half-order $t$.  All extraction parameters are exact rationals, recorded
in the program of Appendix~\ref{app:ninth-field}.  The printed decimals are exact.

\begin{proof}
Use Lemma~\ref{lem:field-general-seed-aux} with $k=9$.
Every row has $1\le\ell\le2$ and $j+\ell\le8$.
The sparse alternative is $j\le\lfloor(9-j-\ell+1)/2\rfloor+1$;
otherwise the chronological inequalities are $\Phi_i+7\phi_i\le1$.
Every higher even residual degree satisfies $\Phi_j\le1/3$.

Apply the A/B terminal and backwards recurrence from
Section~\ref{sec:efficient-differencing}, now with $k=9$.
The exact checks use all twenty rows, the interpolation ranges, the
extraction inequalities $M_i\le Q_{i-1}$, and every actual appended
residual degree.  Their minimum strict supersolution slack is
\[
 \delta=\min_{3\le s\le22}\{b_s-T_s(b)\}
       =\frac{43740673471}{40000000000000000}>0.
\]
The chronological inequalities have minimum margin
$1/5000000000>0$. Choose the smoothness parameter after this margin and
all the finite applications below.

The coefficient sums are at most $982096383/10^9<1$.
Lemma~\ref{lem:supersolution} proves the base curve.

Finally, apply the same fixed twenty rows for twenty-four further
simultaneous rounds.  Every row reads only the exponents available at
the start of its round.  After a row is evaluated, round its output
\emph{up} to a multiple of $10^{-12}$ and retain the better of this value
and the previous exponent.  This gives precisely the final column.
The program replays all twenty-four rounds with rational
arithmetic and checks every entry against its exact printed value.
Each operation is an application of an established field estimate,
H\"older interpolation, or monotonicity of permissible exponents.
Only finitely many operations occur, so one sufficiently small
smoothness parameter works for the whole proof.
\end{proof}

\subsection{The degree-eighteen coefficient seed curve}
\label{sec:degree18-field-seeds}

The coefficient sums in the ninth-power terminal have degree eighteen.
We supply their moments through full order $62$ over the field itself.
As before, $X=P^{[K:\mathbb Q]}$ is the root-volume parameter, and
$\lambda_s$ denotes the full exponent at moment order $2s$.
Set $\lambda_0=0$, $\lambda_1=1$, $\lambda_2=2$, and take the following
exact decimal values for the remaining entries:
\[
\begin{array}{c|r@{\qquad}c|r}
s&\lambda_s&s&\lambda_s\\\hline
3&3.005107&18&22.723143\\
4&4.048469&19&24.306116\\
5&5.146988&20&25.918677\\
6&6.289328&21&27.559735\\
7&7.467683&22&29.228132\\
8&8.681624&23&30.922636\\
9&9.933259&24&32.641958\\
10&11.218619&25&34.384850\\
11&12.538919&26&36.150022\\
12&13.893656&27&37.936763\\
13&15.282351&28&39.742902\\
14&16.705014&29&41.567190\\
15&18.160923&30&43.408399\\
16&19.649745&31&45.265286\\
17&21.170711&&
\end{array}
\]

\begin{proposition}\label{prop:degree18-field-seeds}
For every number field $K$, each entry in the table is permissible for
the degree-eighteen ideal-smooth moment. In particular, for every
$\epsilon>0$ and sufficiently small $\eta>0$,
\[
 \int_{\mathbb T_K}|g_{18}(\alpha;X,R)|^{62}\,d\alpha
       \ll_{K,\epsilon}X^{45.265286+\epsilon}
       \qquad(1\le R\le X^\eta).
\]
The same convention applies to every order in the table and to bounded
nonnegative weights and fixed root boxes.
\end{proposition}

\begin{proof}
Use Lemma~\ref{lem:field-general-seed-aux} with $k=18$.
Each row has $1\le\ell\le3$ and $j+\ell\le17$.
Its sparse alternative is $j\le\lfloor(18-j-\ell+1)/2\rfloor+1$;
otherwise the chronological inequalities are $\Phi_i+16\phi_i\le1$.
The higher even residual degrees satisfy $\Phi_j\le1/3$.

The exact certificate in Appendix~\ref{app:ninth-field} gives one A or B row for each
$s=3,\ldots,31$. It checks all these hypotheses for every appended
residual degree, all extraction inequalities $M_i\le Q_{i-1}$, and
every interpolation argument. Among rows requiring the curve condition,
the largest coefficient-product exponent is
\[
 \Phi_j=\frac{3267856441}{10000000000}<\frac13.
\]
Every chronological condition is strict; its smallest positive margin is
$1/5000000000$. Such fixed positive margins permit the smoothness
parameter to be chosen afterwards. Zero extraction exponents occur at
half-orders $22,23,25,26,27,28$ and have the meaning explained in
Section~\ref{sec:efficient-differencing}: the nominal extraction threshold is one,
while fixed powers of the smoothness parameter are absorbed in
$X^\epsilon$.

Let $T_s$ be the exact A/B recurrence defined in
Section~\ref{sec:efficient-differencing}, with the parameters in this
certificate. All twenty-nine rows satisfy $T_s(\lambda)<\lambda_s$.
The least exact supersolution slack is
\[
 \min_{3\le s\le31}(\lambda_s-T_s(\lambda))
    =\frac{72464782233}{160000000000000000}>0.
\]
The maximum coefficient sum in the input moment exponents is at most
\[
                    \frac{9974554911}{10000000000}<1.
\]
Lemma~\ref{lem:supersolution} now proves every entry.
\end{proof}

All coefficient orders requested by the ninth-power policies are
even and at most $62$. The table therefore supplies them by direct use
or by H\"older interpolation between available orders.


\subsection{The ninth-power quadratic policies}
Apply Section~\ref{sec:quadratic-terminals} with $k=9$.
Each policy uses $a=0$, $c=\gamma=1$, and $v=2r_j$.
Only the quadratic second resource and basic or refined coefficient
resources are required.

The policies have the following exact terminal data:
\[
\begin{array}{c|c|c|c|c|c}
2s&j&(r_1,\ldots,r_j)&v&h_*&2^{j-7}\\\hline
34&6&(16,15,14,13,12,11)&22&2.00000221&1/2\\
36&6&(17,16,15,14,13,12)&24&1.88634850&1/2\\
38&6&(18,17,16,15,14,13)&26&1.84146280&1/2\\
40&5&(19,18,17,16,15)&30&2.44508497&1/4
\end{array}
\]
Every policy has $H_*>X$ and $X<Q_0^{9/2}$ strictly.  Every refined
preliminary size cap is positive.  The coefficient orders are even and
at most $62$, so the degree-eighteen seed curve through half-order $31$
suffices.  Each policy has four assigned normalized height intervals.
All parameters are specified as exact rationals in Appendix~\ref{app:ninth-field}.

\subsection{The moment and the global bound}
\label{sec:ninth-global}
\begin{theorem}[A smooth ninth-power moment]
\label{thm:ninth-moment}
For every fixed number field $K$ and every $\epsilon>0$, there is
$\eta>0$ such that
\[
 \begin{gathered}
 \int_{\mathbb T_K}|g_9(\alpha;X,R)|^{40}\,d\alpha
       \ll_{K,\epsilon}X^{31+\Delta_{40}+\epsilon}
       \qquad(1\le R\le X^\eta),\\
 \Delta_{40}=\frac{8516302869}{62500000000}
            =0.136260845904.
 \end{gathered}
\]
The assertion includes the fixed boxes and bounded nonnegative weights
used in the representation argument.
\end{theorem}

\begin{proof}
Appendix~\ref{app:ninth-field} contains all twenty degree-nine base
rows, all twenty-nine degree-eighteen rows, the twenty-four further
degree-nine seed refinement rounds, and all four fixed terminal policies.
The seed propositions prove the input moments over $K$. Their
strict supersolutions follow from Lemma~\ref{lem:supersolution}.

Apply the four policies at full moment orders $34,36,38,40$ for seven
simultaneous rounds. Each round reads only its initial curve. H\"older
interpolation over all available integer half-orders is applied at the
start of a round, and the better of an old exponent and a new exponent
is retained. Every new exponent is rounded \emph{up} to a multiple of
$10^{-12}$. The full basic coefficient estimate is used at both depths;
the equality of its minimum with the selected $\phi_i$ is checked exactly.

For every policy the program verifies the resource orders, the coefficient
sampling size caps, the common H\"older masses and factor loads, all
height interval endpoints, and the backwards counted recurrence.
Section~\ref{sec:ninth-terminal} proves these exponents for the field
sums. Convexity on each of the four assigned height intervals makes
the two endpoint checks sufficient on the entire interval.
The resulting exact upper values are
\[
\begin{array}{c|c|c}
 \text{moment order}&\text{full permissible exponent}&\text{excess over }2s-9\\\hline
 34&25.329443645750&0.329443645750\\
 36&27.249606352140&0.249606352140\\
 38&29.186123445082&0.186123445082\\
 40&31.136260845904&0.136260845904
\end{array}
\]
The last entry proves the theorem. All displayed decimals are exact
rational numbers; none is an approximation used to decide an inequality.
The finite iteration allows the smoothness parameter and the individual
error exponents to be chosen after all numerical parameters.
\end{proof}


\begin{corollary}\label{thm:ninth-transfer}
For every number field $K$, $G_{\mathbb K}(9)\le46$.
\end{corollary}
\begin{proof}
Apply Corollary~\ref{cor:even-transfer} with $(w,S,u)=(90,46,40)$.
Its exact margin is
\[
 D=6-44\Delta_{40}=\frac{70668441}{15625000000}>0.
\]
Since $h_9\le30$, the primitive local hypothesis holds in
forty-six variables. This proves the assertion.
\end{proof}

\section{Tenth powers}
\label{sec:tenth-powers}
The cubic terminal suffices to attain the bound $55$ for tenth powers.
The same seed curve supplies the lower moment used in the quintic argument.

\begin{proposition}\label{prop:tenth-moments}
For every fixed number field $K$, the smooth tenth-power sum satisfies
\[
 \int_{\mathbb T_K}|g(\alpha)|^{44}\,d\alpha
 \ll X^{34+48479/250000+\epsilon},\qquad
 \int_{\mathbb T_K}|g(\alpha)|^{46}\,d\alpha
 \ll X^{36+152339/1000000+\epsilon}.
\]
The smoothness parameter may be chosen sufficiently small for these
estimates to hold simultaneously.
\end{proposition}
\begin{proof}
Write $b_s$ for the exponent of the $2s$th smooth moment, with
$b_0=0$, $b_1=1$ and $b_2=2$. Table~\ref{tab:tenth-seeds} gives a
rational simultaneous supersolution for the $A$, $B$ and cubic terminal
processes. The notation $A_j^\ell$, $B_j^\ell(t)$ specifies the
 efficient-differencing depth, terminal auxiliary moment and, in the
second case, the smooth half-order. The symbol $M_7$ denotes the cubic
terminal at depth seven.

The terminal formula for this last process is
\begin{equation}\label{eq:tenth-cubic-terminal}
 C_7=\max\left\{
 a_0-\frac14+q_7b_{s-7},\quad
 a_0+q_7\left(-\frac52+\frac34b\left(\frac{4(s-7)}3\right)\right)
 \right\},\qquad a_0=8-9\Phi_7.
\end{equation}
Here $b(t)$ is linear interpolation between consecutive entries.
The first term follows from the cubic minor-arc saving and the second
from the fourth moment on its auxiliary major arcs. Apply the common
backward recurrence to this terminal.

All scale parameters are rational with denominator dividing $10^{10}$.
The exact certificate checks $0<\phi_i<1/10$, all chronological divisor
and small-product conditions required by the auxiliary moments, and
$10q_7>2$ for each cubic terminal. Its smallest supersolution margin is
\[
 \frac{65679537093}{40000000000000000}>0,
\]
and the common contraction factor is at most
$9859429631/10^{10}<1$. The simultaneous supersolution lemma therefore
proves all twenty-one rows. The entries at $s=22,23$ give the proposition.
\end{proof}

\begin{corollary}\label{prop:tenth44}
The exponent $34.227$ is permissible for the smooth tenth-power moment
of order $44$ over every fixed number field.
\end{corollary}
\begin{proof}
The sharper exponent above is $34.193916<34.227$.
\end{proof}

\begin{table}[htbp]
\centering\small
\begin{tabular}{r r l@{\qquad}r r l}
\hline
$s$ & $b_s$ & Process & $s$ & $b_s$ & Process\\
\hline
3 & 3.028131 & $A_{1}^{1}$ & 14 & 19.142712 & $B_{6}^{2}(11)$\\
4 & 4.156832 & $A_{1}^{2}$ & 15 & 20.930391 & $A_{6}^{2}$\\
5 & 5.365973 & $A_{2}^{1}$ & 16 & 22.754051 & $B_{6}^{3}(13)$\\
6 & 6.646618 & $B_{2}^{2}(6)$ & 17 & 24.609057 & $A_{7}^{2}$\\
7 & 7.996737 & $A_{2}^{2}$ & 18 & 26.489639 & $M_7$\\
8 & 9.414374 & $B_{3}^{2}(7)$ & 19 & 28.391489 & $M_7$\\
9 & 10.894592 & $A_{3}^{2}$ & 20 & 30.311466 & $M_7$\\
10 & 12.437588 & $A_{4}^{2}$ & 21 & 32.246495 & $M_7$\\
11 & 14.037216 & $A_{4}^{2}$ & 22 & 34.193916 & $M_7$\\
12 & 15.691422 & $A_{5}^{2}$ & 23 & 36.152339 & $M_7$\\
13 & 17.394386 & $A_{5}^{2}$ &  & & \\
\hline
\end{tabular}
\caption{A simultaneous tenth-power supersolution. All displayed
 decimals are exact rational numbers.}\label{tab:tenth-seeds}
\end{table}

\begin{proposition}\label{prop:tenth-55-analytic}
The analytic part of the number-field Waring problem for tenth powers
holds with $55$ variables.
\end{proposition}
\begin{proof}
Set $\Delta_{46}=152339/1000000$. The ordinary minor-arc moment has
order $110$ and exponent $99+\epsilon$. Apply the even-moment bootstrap
with $u=46$, $S=56$ and $w=110$. Its strict margin is
\[
 (56-46)-(110-56)\Delta_{46}=\frac{886847}{500000}>0.
\]
Thus the smooth moment of order $56$ has exponent $46$ without an
$X^\epsilon$ loss.

Use fifteen ordinary roots and forty product roots. On the ordinary
minor arcs, H\"older's inequality has weights
\[
 \frac3{22},\qquad\frac{46}{55},\qquad\frac3{110}
\]
for moments $110$, $46$ and $56$, respectively. These weights sum to one;
the ordinary load is $15$ and the product-root load is $40$.
The minor-arc contribution is consequently
\[
 \ll X^{45-\delta+\epsilon},\qquad
 \delta=\frac3{22}-\frac{46}{55}\Delta_{46}
       =\frac{246203}{27500000}>0.
\]
On the major-arc annulus, use the product moment of order $56$ and the
ordinary moment of order
\[
 p=\frac{15}{1-40/56}=\frac{105}{2}>20.
\]
The annular pruning estimate supplies a power saving. The product-root
major-arc calculation gives a positive main term under the primitive
local hypotheses, proving the assertion.
\end{proof}

The local bound for tenth powers is smaller than $55$. Combining it with
Proposition~\ref{prop:tenth-55-analytic} gives $G_{\mathbb K}(10)\le55$.

\appendix
\section{Exact parameter data}\label{app:certificates}
The following programs accompany the supplement. They perform their
inequalities over the rational numbers. The local computations enumerate
all the indicated residue classes. Run \href{https://github.com/morley988/warings-problem/blob/3752a566d4ae778c3f2effaeb7caec4fc8899514/supplements/number-field/small-powers/verify.py}{\texttt{python3 verify.py}} from
the directory containing the source of this supplement.

\subsection{Fourth powers}\label{app:quartic-exact-certificate}
\href{https://github.com/morley988/warings-problem/blob/3752a566d4ae778c3f2effaeb7caec4fc8899514/supplements/number-field/small-powers/certificates/quartic-exact-certificate.py}{\texttt{certificates/quartic-exact-certificate.py}}.

\subsection{Fifth powers}\label{app:quintic}
\href{https://github.com/morley988/warings-problem/blob/3752a566d4ae778c3f2effaeb7caec4fc8899514/supplements/number-field/small-powers/certificates/quintic.py}{\texttt{certificates/quintic.py}}.

\subsection{Sixth powers}\label{app:sixth}
\href{https://github.com/morley988/warings-problem/blob/3752a566d4ae778c3f2effaeb7caec4fc8899514/supplements/number-field/small-powers/certificates/sixth.py}{\texttt{certificates/sixth.py}}.

\subsection{Seventh powers}\label{app:septic}
\href{https://github.com/morley988/warings-problem/blob/3752a566d4ae778c3f2effaeb7caec4fc8899514/supplements/number-field/small-powers/certificates/septic.py}{\texttt{certificates/septic.py}}.

\subsection{Eighth powers}\label{app:eighth-field}
\href{https://github.com/morley988/warings-problem/blob/3752a566d4ae778c3f2effaeb7caec4fc8899514/supplements/number-field/small-powers/certificates/eighth-field.py}{\texttt{certificates/eighth-field.py}}.

\subsection{Ninth powers}\label{app:ninth-field}
\href{https://github.com/morley988/warings-problem/blob/3752a566d4ae778c3f2effaeb7caec4fc8899514/supplements/number-field/small-powers/certificates/ninth-field.py}{\texttt{certificates/ninth-field.py}}.

\subsection{Tenth powers}\label{app:tenth-field}
\href{https://github.com/morley988/warings-problem/blob/3752a566d4ae778c3f2effaeb7caec4fc8899514/supplements/number-field/small-powers/certificates/tenth-field.py}{\texttt{certificates/tenth-field.py}}.

\subsection{Eighth carry relations}\label{app:eighth-ghost}
\href{https://github.com/morley988/warings-problem/blob/3752a566d4ae778c3f2effaeb7caec4fc8899514/supplements/number-field/small-powers/certificates/eighth-ghost.py}{\texttt{certificates/eighth-ghost.py}}.

\subsection{Eighth powers in ramification degree five}\label{app:eighth-tame}
\href{https://github.com/morley988/warings-problem/blob/3752a566d4ae778c3f2effaeb7caec4fc8899514/supplements/number-field/small-powers/certificates/eighth-tame.py}{\texttt{certificates/eighth-tame.py}}.

\section*{Acknowledgement}
Some thirty years ago, as a doctoral student at the University of Michigan,
I had the good fortune to learn about Waring's problem from Robert Vaughan
and my advisor, Trevor Wooley. I remain grateful to them both for their
teaching and guidance, and for the clarity and beauty of their writing
on the subject. In revising this supplement, I hope to bring its arguments
into clearer order and to give readers something of the pleasure I found
in their papers.

\section*{Use of generative AI}
The author used OpenAI's ChatGPT in developing the arguments, checking
parameter choices and preparing the manuscript.
\begin{thebibliography}{99}
\bibitem{BCK2025}G. Binyamini, R. Cluckers and F. Kato,
\emph{Sharp bounds for the number of rational points on algebraic curves and
dimension growth, over all global fields}, Proc. London Math. Soc. (3)
\textbf{130} (2025), e70016. Corollary 12 and the height definition preceding (28).

\bibitem{BW2023}J. Br\"udern and T. D. Wooley, \emph{On Waring's problem for larger powers},
J. reine angew. Math. \textbf{805} (2023), 115--142. Theorem 4.2.

\bibitem{Cluckers2011}R. Cluckers, \emph{Analytic van der Corput lemma for $p$-adic and
$\mathbb F_q((t))$ oscillatory integrals, singular Fourier transforms,
and restriction theorems}, Exposition. Math. \textbf{29} (2011), 371--386.
Proposition 3.3.

\bibitem{Hormander2003}L. H\"ormander, \emph{The Analysis of Linear Partial Differential
Operators I: Distribution Theory and Fourier Analysis}, 2nd ed.,
Classics in Mathematics, Springer, Berlin, 2003.

\bibitem{Joly1970}J.-R. Joly, \emph{Sommes de puissances $d$-i\`emes dans un anneau commutatif},
Acta Arith. \textbf{17} (1970), 37--114. Proposition 7.16 and
Theorem 7.34, pp. 95 and 105--106.

\bibitem{Ramanujam1963}C. P. Ramanujam, \emph{Sums of $m$-th powers in $p$-adic rings},
Mathematika \textbf{10} (1963), 137--146.

\bibitem{VW1995}R. C. Vaughan and T. D. Wooley, \emph{Further improvements in Waring's problem},
Acta Math. \textbf{174} (1995), 147--240.

\bibitem{VW2000}R. C. Vaughan and T. D. Wooley,
\emph{Further improvements in Waring's problem, IV: Higher powers},
Acta Arith. \textbf{94} (2000), 203--285.

\bibitem{Weil1948}A. Weil, \emph{On some exponential sums},
Proc. Natl. Acad. Sci. USA \textbf{34} (1948), 204--207.

\bibitem{Wooley1992}T. D. Wooley, \emph{Large improvements in Waring's problem},
Ann. of Math. (2) \textbf{135} (1992), 131--164.

\bibitem{Wooley2016}T. D. Wooley, \emph{On Waring's problem for intermediate powers},
Acta Arith. \textbf{176} (2016), 241--247.

\bibitem{Wooley2019}T. D. Wooley, \emph{Nested efficient congruencing and relatives of
Vinogradov's mean value theorem}, Proc. London Math. Soc. (3)
\textbf{118} (2019), 942--1016.

\bibitem{WooleyThesis}T. D. Wooley, \emph{On simultaneous additive equations and Waring's problem},
Ph.D. thesis, University of London, 1990, Chapter 6, especially pp. 230--243.

\bibitem{hazewinkel2009}M. Hazewinkel, \emph{Witt vectors. Part~1}, in M. Hazewinkel (ed.),
\emph{Handbook of Algebra}, vol.~6, Elsevier, 2009, pp.~319--472.
Sections~5.2--5.22.
\end{thebibliography}
\end{document}
